% ===========================================================================
%  Chapter 3 -- Methodology
%  Disease-Adaptive Heterogeneous Graph Learning with Anatomically Aligned
%  Multi-Sequence MRI Representations
%
%  PROTOCOL-GRADE VERSION
%  This chapter is written prospectively because the full PhD experiment has not
%  yet been completed. After implementation, convert prospective statements to
%  past tense and replace all [TO CONFIRM] / [TO RECORD] fields with the executed
%  protocol and observed cohort counts. Do not silently fill them from intention.
%
%  Compiles standalone:   pdflatex chapter3 ; biber chapter3 ; pdflatex chapter3
%
%  To fold into a larger thesis instead, delete everything from \documentclass
%  down to \begin{document}, and the \printbibliography/\end{document} at the
%  foot, then % ===========================================================================
%  Chapter 3 -- Methodology
%  Disease-Adaptive Heterogeneous Graph Learning with Anatomically Aligned
%  Multi-Sequence MRI Representations
%
%  PROTOCOL-GRADE VERSION
%  This chapter is written prospectively because the full PhD experiment has not
%  yet been completed. After implementation, convert prospective statements to
%  past tense and replace all [TO CONFIRM] / [TO RECORD] fields with the executed
%  protocol and observed cohort counts. Do not silently fill them from intention.
%
%  Compiles standalone:   pdflatex chapter3 ; biber chapter3 ; pdflatex chapter3
%
%  To fold into a larger thesis instead, delete everything from \documentclass
%  down to \begin{document}, and the \printbibliography/\end{document} at the
%  foot, then % ===========================================================================
%  Chapter 3 -- Methodology
%  Disease-Adaptive Heterogeneous Graph Learning with Anatomically Aligned
%  Multi-Sequence MRI Representations
%
%  PROTOCOL-GRADE VERSION
%  This chapter is written prospectively because the full PhD experiment has not
%  yet been completed. After implementation, convert prospective statements to
%  past tense and replace all [TO CONFIRM] / [TO RECORD] fields with the executed
%  protocol and observed cohort counts. Do not silently fill them from intention.
%
%  Compiles standalone:   pdflatex chapter3 ; biber chapter3 ; pdflatex chapter3
%
%  To fold into a larger thesis instead, delete everything from \documentclass
%  down to \begin{document}, and the \printbibliography/\end{document} at the
%  foot, then % ===========================================================================
%  Chapter 3 -- Methodology
%  Disease-Adaptive Heterogeneous Graph Learning with Anatomically Aligned
%  Multi-Sequence MRI Representations
%
%  PROTOCOL-GRADE VERSION
%  This chapter is written prospectively because the full PhD experiment has not
%  yet been completed. After implementation, convert prospective statements to
%  past tense and replace all [TO CONFIRM] / [TO RECORD] fields with the executed
%  protocol and observed cohort counts. Do not silently fill them from intention.
%
%  Compiles standalone:   pdflatex chapter3 ; biber chapter3 ; pdflatex chapter3
%
%  To fold into a larger thesis instead, delete everything from \documentclass
%  down to \begin{document}, and the \printbibliography/\end{document} at the
%  foot, then \input{chapter3} from the main thesis.
% ===========================================================================

\documentclass[12pt,a4paper]{report}

\usepackage{lmodern}
\usepackage[T1]{fontenc}
\usepackage[utf8]{inputenc}
\usepackage[english]{babel}
\usepackage{microtype}
\usepackage[margin=2.5cm]{geometry}
\usepackage{setspace}
\onehalfspacing

\usepackage{amsmath,amssymb}
\usepackage{booktabs}
\usepackage{longtable}
\usepackage{array}
\usepackage{graphicx}
\usepackage[hidelinks]{hyperref}

\usepackage[
  backend=biber,
  style=numeric-comp,
  sorting=none,
  maxbibnames=6,
  minbibnames=3,
  giveninits=true,
  doi=true,
  url=false,
  isbn=false
]{biblatex}

% Same bibliography used by Chapter 2.
\addbibresource{../lumbar_spine_mri_ai_literature_inventory.bib}

\AtEveryBibitem{%
  \clearfield{note}%
  \clearfield{comment}%
  \clearfield{annotation}%
  \clearfield{abstract}%
  \clearfield{file}%
}

\newcommand{\toconfirm}[1]{\textbf{[TO CONFIRM: #1]}}
\newcommand{\torecord}[1]{\textbf{[TO RECORD AFTER IMPLEMENTATION: #1]}}

\begin{document}

\chapter{Methodology}
\label{ch:methodology}

% ===========================================================================
\section{Introduction and Methodological Position}
\label{sec:method-intro}

Chapter~\ref{ch:litreview} established that automated lumbar MRI grading has
moved beyond the stage at which novelty can be claimed from replacing one image
backbone with another. Anatomical localisation, multi-sequence regions of
interest, inter-level context and ordinal grading have all been demonstrated in
recent work. The methodological question in this thesis is therefore not whether
a sufficiently large network can classify lumbar MRI. It is whether the
\emph{structure of the prediction problem} can be represented more faithfully:
which disease is being graded, at which level and side, which MRI sequence
contains the relevant evidence, how neighbouring targets are related, and how
the resulting model behaves when moved to a hospital whose scanners, protocols,
population and reporting conventions were not represented during development.

The study is designed as a staged, controlled experimental programme around
three methodological contributions and one translational validation programme.
The first contribution is an anatomically aligned cross-sequence
self-supervised representation in which positive relationships are defined by
DICOM patient-space correspondence rather than by generic image similarity.
The second is a disease-conditioned sequence router that learns how much each
available MRI sequence contributes to each anatomical target and is trained
explicitly to remain operational when one or more sequences are absent. The
third is a heterogeneous disease--anatomy graph whose nodes correspond to
level--condition--laterality targets and whose typed edges encode anatomical
adjacency, within-level relationships and bilateral symmetry. The translational
programme then separates development from deployment: the model is trained on a
public multinational benchmark, frozen, evaluated zero-shot on an independent
Rizgary Teaching Hospital cohort, and only afterwards adapted with controlled
amounts of local labelled data.

Several components are deliberately treated as \emph{enabling technology}
rather than as doctoral contributions. Anatomical localisation, segmentation,
2.5D region construction, conventional convolutional or transformer encoders,
class-imbalance handling, ordinal losses and calibration are all necessary to
make the comparison scientifically fair, but none is advertised as new by
itself. This distinction follows directly from the novelty boundary established
in Chapter~\ref{ch:litreview}: the thesis tests a new formulation of the
structured task rather than presenting a collection of familiar modules as if
their combination alone constituted originality.

The chapter is written as a protocol-grade methodology. Wherever an
implementation choice has not yet been fixed by evidence, it is specified as a
controlled comparison rather than invented retrospectively. Wherever an
institutional fact remains unresolved---for example the final number of
Rizgary studies surviving data reconciliation or the formal ethics reference---
it is marked explicitly. This is intentional. A doctoral methodology should
describe what is known, what is pre-specified and what will be decided by a
transparent rule; it should not convert planning assumptions into apparently
observed facts.

\subsection{Research Questions}
\label{sec:method-rqs}

The five research questions derived in Chapter~\ref{ch:litreview} are
operationalised as follows.

\begin{description}
  \item[RQ1 --- Relational disease modelling.] Does representing lumbar disease
  targets as a typed heterogeneous graph improve grading over independent target
  heads, an ordered five-level transformer and a homogeneous graph, and are any
  gains concentrated where single-target evidence is weakest?

  \item[RQ2 --- Anatomically aligned self-supervision.] Does cross-sequence
  pretraining based on DICOM-defined anatomical correspondence improve grading,
  label efficiency and transfer compared with ImageNet initialisation, generic
  medical pretraining and ordinary augmentation-based contrastive learning?

  \item[RQ3 --- Disease-conditioned modality routing.] Can one model learn
  target- and patient-dependent sequence weights while remaining robust to
  arbitrary missing-sequence configurations, without requiring a separately
  trained network for every combination of sagittal T1, sagittal T2/STIR and
  axial T2?

  \item[RQ4 --- Cross-institutional transfer.] What is the zero-shot degradation
  when a benchmark-trained model is applied to a previously unseen Middle
  Eastern clinical cohort, and how does performance recover as the amount of
  local annotation available for adaptation increases?

  \item[RQ5 --- Clinically asymmetric error and uncertainty.] Do ordinal and
  cost-sensitive objectives reduce clinically consequential distant errors,
  particularly Severe$\rightarrow$Normal/Mild, and can calibrated uncertainty
  support selective prediction without overstating clinical safety?
\end{description}

The order is not arbitrary. RQ2 and RQ3 determine how evidence is represented;
RQ1 determines how target predictions communicate; RQ5 determines how the
ordered output and its uncertainty are represented; and RQ4 asks whether those
choices survive the transition from a benchmark to real local clinical data.

\subsection{Design Principles}
\label{sec:method-principles}

Six principles govern all experiments.

\textbf{First, anatomy precedes grading.} Every classifier and ablation receives
the same anatomically localised input so that a comparison of representations is
not confounded by one model being given a better crop than another. This follows
the consistent evidence that localisation is a precondition of reliable
per-level grading \cite{lu2018deepspine,pan2021automatically,batra2025mscan}.

\textbf{Second, DICOM geometry is preserved.} The primary scientific
representation is reconstructed from DICOM patient-space coordinates. PNG
conversion may be used for visual inspection but is not the basis of
cross-sequence alignment, because it discards the spatial metadata required to
determine which axial slices correspond to a sagittal disc level.

\textbf{Third, patient identity defines independence.} No images, slices,
levels, sides or sequences from one patient may cross between training,
validation and test partitions. The unit of statistical independence for
uncertainty intervals and resampling is the patient, not the image or the disc
level.

\textbf{Fourth, the external cohort is not a tuning set.} The primary model is
developed without Rizgary images. Zero-shot performance is measured before any
local tuning. Few-shot adaptation then uses a separately defined local
adaptation pool while evaluation remains on a fixed local test set.

\textbf{Fifth, every claimed contribution requires an ablation.} A complete
model score does not establish that anatomical self-supervision, routing or
graph topology contributed anything. Each component is added and removed under
a controlled experiment, and a negative result is reported as such.

\textbf{Sixth, clinical interpretation is bounded by the reference standard.}
The model is trained to reproduce radiological grades whose inter-reader
reliability is imperfect. Performance is therefore interpreted against human
agreement and not as evidence that the network has discovered a more objective
biological truth.

% ===========================================================================
\section{Overall Research Design}
\label{sec:method-design}

\subsection{Study Type}
\label{sec:method-study-type}

The project is a retrospective, multi-dataset medical-imaging methods study with
three linked stages: public benchmark development, controlled methodological
ablation, and independent cross-institutional transfer. It does not prospectively
alter patient care and does not use model output to make treatment decisions.
The Rizgary component is an external retrospective evaluation of imaging and
radiology reports already acquired in routine clinical care.

The core development dataset is the RSNA Lumbar Degenerative Imaging Spine
Classification benchmark (LumbarDISC), which contains multi-sequence MRI and
expert per-level severity labels \cite{rsna2024rsna,richards2026the}. SPIDER is
used as an anatomical localisation/segmentation resource where its licence and
task definition permit \cite{vandergraaf2024lumbar}. The local Rizgary cohort is
reserved for external transfer and domain-adaptation analysis rather than being
mixed into the benchmark training data.

This separation is central to the inference the study is intended to support. A
random split of one pooled dataset can estimate internal generalisation under
that dataset's acquisition distribution. It cannot answer whether a model
trained on international benchmark data transports to a different hospital.
The latter requires the hospital to remain unseen during development.

\subsection{Unit of Analysis}
\label{sec:method-unit}

The primary unit of prediction on LumbarDISC is a
\emph{level--condition--laterality target}. At five disc levels
(L1--L2, L2--L3, L3--L4, L4--L5 and L5--S1), the benchmark defines five targets:
central canal stenosis, left and right neural foraminal narrowing, and left and
right subarticular stenosis. Full benchmark scope therefore contains
$5\times5=25$ targets per examination. Each target receives the ordered label
Normal/Mild, Moderate or Severe \cite{richards2026the}.

The unit of \emph{statistical independence}, however, remains the patient.
Twenty-five outputs from the same examination are correlated measurements, not
twenty-five independent observations. This distinction affects data splitting,
confidence intervals, significance tests and error analysis throughout the
chapter.

For the local Rizgary cohort, the target space is narrower. The reports contain
disc morphology and central-canal descriptions but do not reproduce the entire
RSNA schema reliably. The primary cross-institutional comparison is therefore
restricted to the five central-canal targets, L1--L2 through L5--S1, for which a
defensible overlap can be created. Bulge, protrusion and extrusion are analysed
as a separate local morphology task rather than being relabelled as if they were
RSNA stenosis grades. Foraminal and subarticular transfer are excluded unless
the local reference standard is expanded by new expert grading.

\subsection{Experimental Staging}
\label{sec:method-staging}

The methodology is intentionally cumulative:

\begin{enumerate}
  \item establish reproducible DICOM reconstruction, patient-level splits,
  localisation and an independent ROI baseline;
  \item establish anatomically correct cross-sequence correspondence;
  \item test disease-conditioned sequence routing and missing-modality training;
  \item test anatomically aligned cross-sequence self-supervision;
  \item test homogeneous and heterogeneous graph reasoning against independent
        target prediction;
  \item test ordinal/cost-sensitive output heads and probability calibration;
  \item freeze the selected public-data model and perform zero-shot external
        evaluation at Rizgary;
  \item adapt with prespecified numbers of local cases and estimate the
        annotation--performance recovery curve.
\end{enumerate}

The staging protects the thesis from a common failure mode of complex medical AI
systems: if the final model underperforms, a monolithic experiment cannot reveal
whether the failure arose from localisation, fusion, representation, graph
reasoning, loss design or domain shift. Here each stage creates a measurable
intermediate object and a direct baseline.

\subsection{Pre-Specified Hypotheses}
\label{sec:method-hypotheses}

The principal hypotheses are directional but not expressed as arbitrary target
accuracies.

\begin{description}
  \item[H1] Anatomical cross-sequence self-supervision will improve
  label-efficient grading and cross-institutional robustness relative to generic
  pretraining, particularly when only a fraction of benchmark labels is used.

  \item[H2] Disease-conditioned routing with modality dropout will degrade less
  under missing-sequence experiments than fixed concatenation or ordinary
  cross-attention trained only on complete studies.

  \item[H3] A typed heterogeneous graph will outperform independent target heads
  and a homogeneous graph if the anatomical edge types carry useful relational
  information; performance under a shuffled-edge control should deteriorate if
  the topology itself matters.

  \item[H4] Zero-shot performance on Rizgary will be lower than internal
  benchmark performance because of scanner, protocol, population and reference-
  standard shift; the magnitude and pattern of that drop are treated as findings
  rather than failures of the project.

  \item[H5] Local adaptation will improve performance as annotation increases,
  but no predetermined percentage of ``recovery'' is required. The shape of the
  recovery curve is the result.

  \item[H6] Cost-sensitive/ordinal objectives may reduce distant
  Severe$\rightarrow$Normal/Mild errors even if they do not improve aggregate
  macro-F1; plain cross-entropy remains the mandatory baseline because previous
  lumbar work has shown that ordinal objectives do not automatically help
  \cite{niemeyer2021a}.
\end{description}

% ===========================================================================
\section{Data Sources}
\label{sec:method-data}

\subsection{Public Development Cohort: LumbarDISC}
\label{sec:method-lumbardisc}

The primary development data are drawn from the publicly released RSNA 2024
Lumbar Spine Degenerative Classification dataset. The published LumbarDISC
resource comprises 2,697 patients and 8,593 MRI series collected across eight
institutions in six countries \cite{richards2026the}. The locally held labelled
competition training subset contains approximately 1,975 studies and is used for
core development. The distinction between the published cohort size and the
available labelled subset is retained explicitly throughout the thesis rather
than treating the two numbers as interchangeable.

Eligible benchmark examinations contain sagittal T2-like imaging, sagittal T1
and axial T2, with expert severity grading at the five lumbar levels
\cite{rsna2024rsna,richards2026the}. Coordinate/localiser annotations are also
available for relevant disease locations. These properties make the dataset
appropriate for all three methodological contributions: it is multi-sequence,
level-resolved, multi-target, geographically heterogeneous and sufficiently
large to separate representation learning from local external validation.

The benchmark serves three roles:

\begin{enumerate}
  \item supervised development of per-target grading;
  \item construction of anatomically corresponding cross-sequence training pairs;
  \item internal/site-aware robustness analysis before the model is exposed to
        Rizgary.
\end{enumerate}

All dataset versions, downloaded files, checksums where feasible and the final
patient split lists will be archived in the reproducibility record. The exact
number of studies remaining after local quality-control exclusions will be
\torecord{final LumbarDISC development, validation and held-out counts}.

\subsection{Public Anatomical Resource: SPIDER}
\label{sec:method-spider}

SPIDER provides manually delineated lumbar anatomy on sagittal MRI from multiple
hospitals, including vertebrae, intervertebral discs and spinal canal
\cite{vandergraaf2024lumbar}. It is used only for anatomical localisation,
segmentation benchmarking or pretraining where permitted by its licence. It is
not used to manufacture disease labels that are absent from the dataset.

The practical purpose is to avoid making localisation the doctoral novelty.
Rather than spend the main experimental budget designing another segmentation
network, the thesis will use or reproduce a strong established anatomical model
and quantify its localisation error. If a SPIDER-pretrained model is used to
initialise the localiser, the external disease-grading experiments remain
separate because the disease targets originate in LumbarDISC.

\subsection{Independent Local Cohort: Rizgary Teaching Hospital}
\label{sec:method-rizgary}

The local cohort consists of lumbar MRI examinations and corresponding English
radiology reports supplied by Rizgary Teaching Hospital. Current project
inventory identifies 299 narrative reports and 294 candidate cases with matching
multi-sequence imaging, while a wider raw DICOM audit contains additional
studies. These counts are not treated as final until a one-to-one cohort
reconciliation is complete.

The raw imaging inspected during protocol development confirms the sequence
families required by the thesis---sagittal T1, sagittal T2 and axial T2---and
preserves standard DICOM spatial metadata. A sample examination was acquired on
a Siemens Avanto 1.5-T system; scanner/vendor distribution across the full local
cohort will be reported after complete metadata extraction rather than inferred
from one case.

The local cohort is valuable precisely because it is not curated to reproduce
the public benchmark. Its scanner, acquisition protocol, prevalence, reporting
language and label completeness differ from LumbarDISC. Those differences are
the domain shift the transfer experiment is intended to measure, not nuisances
to be eliminated by silently harmonising labels or discarding difficult cases.

\subsection{Local Radiology Reports and Historical Spreadsheet}
\label{sec:method-local-reports}

The narrative radiology reports are the primary textual source for the local
reference standard. The existing spreadsheet is treated as a secondary
historical transcription aid rather than as unquestioned ground truth. It
contains patient-level variables and comma-separated level references, but it
does not constitute a complete level-resolved RSNA-style matrix and has known
transcription inconsistencies. The methodology therefore reconstructs the local
reference matrix from source reports and uses the spreadsheet only for
cross-checking.

The canonical local matrix will preserve, at minimum:

\begin{itemize}
  \item pseudonymous case identifier;
  \item lumbar level;
  \item central-canal stenosis severity where explicitly stated;
  \item disc bulge, protrusion and extrusion as separate morphology fields;
  \item foraminal findings and side only where explicitly stated;
  \item nerve-root pressure/impingement where reported;
  \item an explicit missing/not-reported state distinct from a negative finding.
\end{itemize}

This last distinction is essential. Absence of a phrase in a narrative report is
not automatically equivalent to ``normal''. A field may be unreported because
the radiologist considered it clinically irrelevant, because reporting practice
varied, or because the report was not structured around that compartment.
``Not reported'' is therefore encoded separately and excluded from target-
specific external evaluation unless a clinician has explicitly adjudicated the
absence as normal.

\subsection{Cohort Reconciliation}
\label{sec:method-reconciliation}

Before any local model evaluation, a case-flow table will reconcile the raw
archive, reports and analysable imaging. For each candidate case the audit will
record:

\begin{itemize}
  \item whether a report exists;
  \item whether sagittal T1 exists and is usable;
  \item whether sagittal T2/STIR exists and is usable;
  \item whether axial T2 exists and is usable;
  \item whether DICOM geometry is internally consistent;
  \item whether a central-canal reference label can be derived;
  \item whether laterality-dependent labels are available;
  \item whether the case is duplicated or represents repeat imaging;
  \item the reason for any exclusion.
\end{itemize}

The final thesis will report the flow from raw studies to the fixed external
test cohort, rather than presenting the current 299/294/raw-study counts without
explanation. Repeat examinations from one patient, if identified, remain in the
same partition to preserve patient independence.

% ===========================================================================
\section{Ethics, Governance and De-Identification}
\label{sec:method-ethics}

The Rizgary component uses retrospective clinical data. No MRI is acquired for
the purposes of this study, no patient undergoes an additional intervention, and
model output is not returned to the clinical record. Nevertheless, the data are
sensitive medical information and are subject to institutional governance.

The final thesis will state the formal approval mechanism and reference:
\toconfirm{hospital research committee/IRB approval number, date and responsible
institution}. It will also state the legal/ethical basis for retrospective use:
\toconfirm{consent waiver, broad consent or other applicable basis}. These items
must be completed before the local study is described as ethically approved.

The raw local DICOM archive contains direct identifiers and therefore cannot be
given to students or uploaded to external compute in its original state.
De-identification is performed on a controlled copy before analysis. At minimum,
the process removes direct patient identifiers, replaces case identifiers with
study pseudonyms, strips or audits private tags, and handles dates according to
the approved policy while preserving only the temporal information required for
scientific analysis. The de-identification procedure must also check for
burned-in identifiers where applicable. The re-identification key, if one must
exist for hospital governance, is stored separately from the research
repository and is not available to the model-training pipeline.

External/cloud processing of local images is prohibited unless explicitly
allowed by hospital policy:
\toconfirm{whether external cloud/GPU processing is permitted and under what
conditions}. Public benchmark data are handled according to their published
licence and competition terms.

Only aggregate results are reported. Illustrative local images may be included
in the thesis only after de-identification review and only if institutional
permission permits publication. The local cohort is not assumed to be publicly
shareable; any future data release would require separate approval and a more
stringent disclosure review.


% ===========================================================================
\section{Reference Standard and Target Encoding}
\label{sec:method-reference}

\subsection{Benchmark Target Representation}
\label{sec:method-benchmark-targets}

For LumbarDISC, the target set is represented explicitly rather than collapsed
into a generic three-class image label. Let the five lumbar levels be
\[
\mathcal{L}=\{L1\!-\!L2,L2\!-\!L3,L3\!-\!L4,L4\!-\!L5,L5\!-\!S1\},
\]
and let the five condition/laterality targets be
\[
\mathcal{C}=\{
\mathrm{Canal},
\mathrm{LF},
\mathrm{RF},
\mathrm{LSA},
\mathrm{RSA}
\},
\]
where LF/RF denote left/right neural foraminal narrowing and LSA/RSA
left/right subarticular stenosis. One examination therefore has the target set
\[
\mathcal{T}=\mathcal{L}\times\mathcal{C}, \qquad |\mathcal{T}|=25.
\]

For target $t\in\mathcal{T}$ the observed grade is
\[
y_t\in\{0,1,2\},
\]
corresponding to Normal/Mild, Moderate and Severe. Encoding the targets in this
form is not merely convenient notation. The pair $(l,c)$ is retained throughout
data loading, feature extraction, graph construction and evaluation, which
prevents a severity label from losing the anatomical meaning that generated it.

Where the benchmark provides an absent or unusable annotation, a target mask
$m_t\in\{0,1\}$ is carried alongside the grade. Masked targets do not contribute
to the supervised loss. Missing labels are never imputed from neighbouring
levels for the main experiments, because doing so would introduce exactly the
relational assumption the thesis is intended to test.

\subsection{Local Reference Matrix}
\label{sec:method-local-reference}

The Rizgary reference standard is constructed at the same anatomical level but
not forced into the same disease schema. A canonical row represents one
case--level pair. The minimum schema is:

\begin{center}
\begin{tabular}{p{2.4cm}p{2.0cm}p{8.1cm}}
\toprule
\textbf{Field} & \textbf{Type} & \textbf{Interpretation}\\
\midrule
Case ID & categorical & pseudonymous patient/study identifier\\
Level & categorical & L1--L2 through L5--S1\\
Canal stenosis & ordinal/missing & explicit report-derived severity, if stated\\
Bulge & binary/missing & circumferential/generalised bulge stated at level\\
Protrusion & binary/missing & protrusion stated at level\\
Extrusion & binary/missing & extrusion stated at level\\
Foraminal finding & categorical/missing & narrowing/pressure effect if stated\\
Laterality & left/right/bilateral/NR & side only where stated\\
Nerve-root effect & categorical/missing & pressure, displacement or compression\\
Other morphology & text/coded & migration, dehydration, height loss and related findings\\
\bottomrule
\end{tabular}
\end{center}

The local matrix keeps \emph{not reported} (NR) separate from \emph{explicitly
absent}. This rule is especially important for foraminal and subarticular
findings, because local reporting completeness differs from the structured RSNA
annotation scheme. A target is eligible for external evaluation only when the
reference standard provides a defensible class label.

\subsection{Report Extraction and Human Verification}
\label{sec:method-report-verification}

Automated report extraction, if available from the parallel clinical-NLP work,
may accelerate construction of the matrix but does not define the gold
standard. The final labels used for Selar's external evaluation are manually
verified against the original report text.

A two-stage verification procedure is preferred:

\begin{enumerate}
  \item \textbf{Structured extraction audit.} Every non-normal/non-empty finding
  is traced back to the exact report phrase, level and laterality. Multi-level
  expressions such as ``L2--3 and L4--5, L5--S1 shows circumferential disc
  bulge'' are expanded into separate level records without changing their
  meaning.

  \item \textbf{Clinical reliability subset.} A pre-specified stratified subset
  of local studies is independently re-read by two qualified readers using an
  explicit grading sheet. Disagreements are adjudicated by a more senior reader
  where available. Inter-reader agreement is reported with weighted
  $\kappa$ for ordinal grades and ordinary agreement/appropriate $\kappa$ for
  binary morphology. The size of the re-read subset will be fixed before model
  results are examined, based on reader availability and the precision required
  for the agreement estimate.
\end{enumerate}

If a second-reader study cannot be completed, the report-derived matrix remains
usable as a real-world reference standard, but the absence of independent
regrading is stated explicitly as a limitation and no claim is made that local
labels reproduce the RSNA grading protocol exactly.

\subsection{Schema Alignment Between LumbarDISC and Rizgary}
\label{sec:method-schema-alignment}

Cross-institutional evaluation requires equality of \emph{task}, not merely
similar words. The mapping is therefore conservative.

\begin{center}
\begin{tabular}{p{4.1cm}p{4.0cm}p{5.0cm}}
\toprule
\textbf{LumbarDISC target} & \textbf{Rizgary evidence} & \textbf{Transfer status}\\
\midrule
Central canal stenosis & level-resolved mild/moderate/severe statements where present
& Primary external target\\
Left/right foraminal narrowing & pressure/narrowing often incompletely lateralised
& Exclude unless laterality and grade are adequate or freshly re-read\\
Left/right subarticular stenosis & not routinely represented in current reports
& Exclude unless new expert annotation is created\\
Disc bulge/protrusion/extrusion & directly described locally but not an RSNA stenosis target
& Separate local morphology task\\
\bottomrule
\end{tabular}
\end{center}

No local disc-herniation morphology is converted into a stenosis grade by
heuristic rule. Likewise, ``pressure on exit foramina'' is not automatically
treated as a left/right foraminal severity label unless the report or expert
review supplies the information required by that target. This prevents
apparent external validation from being created by label re-interpretation.

% ===========================================================================
\section{Data Partitioning and Leakage Control}
\label{sec:method-splits}

\subsection{Patient-Level Independence}
\label{sec:method-patient-split}

All splits are made at patient/study level before slice extraction or data
augmentation. For a patient $p$, every sagittal series, every axial series,
every disc level, every side and every augmented crop belongs to one and only
one partition. A programmatic assertion is run before training to confirm that
the intersection of patient identifiers across partitions is empty.

The split record is version-controlled as a list of pseudonymous IDs. Data
loaders consume those fixed lists rather than performing a new random split each
time the training script is executed.

\subsection{Public Development, Validation and Held-Out Evaluation}
\label{sec:method-public-split}

Where the available public dataset exposes institution identifiers, a
site-aware evaluation is preferred because it tests a more meaningful shift
than a random split. The exact public partitioning scheme will be chosen
according to the metadata available in the held dataset and will be frozen
before model comparison. At minimum there will be:

\begin{itemize}
  \item a development/training partition;
  \item a validation partition for hyperparameter choice, early stopping and
        calibration;
  \item a held-out public evaluation partition not used for model selection.
\end{itemize}

If site identifiers permit leave-one-site-out or institution-held-out analysis,
that analysis is reported in addition to the ordinary benchmark split. It is not
a replacement for Rizgary, because the local hospital introduces a different
reporting/reference-standard domain as well as a different imaging domain.

\subsection{Fixed Local Test Set and Adaptation Pool}
\label{sec:method-local-split}

The local transfer study is protected against adaptation leakage by separating
Rizgary into two parts \emph{before} the public model is evaluated:

\begin{enumerate}
  \item a \textbf{fixed local test set}, which is never used for adaptation,
        calibration or threshold selection; and
  \item a \textbf{local adaptation pool}, from which the $N=10,25,50,100$
        few-shot subsets are sampled.
\end{enumerate}

The exact number allocated to each part will be determined after the audited
class distribution is known. The split will be stratified by central-canal
severity and lumbar level where feasible without producing clinically
implausible or extremely small strata. The fixed local test set is used for the
zero-shot baseline and for every subsequent adapted model, making the recovery
curves directly comparable.

For each adaptation size $N$, several independent stratified subsets are sampled
from the adaptation pool. Results are averaged across these draws with variance
reported, because one set of ten cases can be unrepresentative by chance. The
test set itself never changes across adaptation sizes.

\subsection{Prevention of Information Leakage During Self-Supervision}
\label{sec:method-ssl-leakage}

Self-supervised learning does not remove the need for split discipline. A model
pretrained on the held-out test patients, even without labels, has already seen
their anatomy and acquisition characteristics. Therefore anatomical
self-supervision uses only the development partition during internal
experiments. The held-out public test set and fixed Rizgary test set remain
unseen in both supervised and self-supervised training.

The same restriction applies to intensity normalisation parameters and learned
harmonisation models: any transformation whose parameters are estimated from
data is fitted on the training/development partition and then applied unchanged
to validation/test data.

% ===========================================================================
\section{DICOM Reconstruction and Quality Control}
\label{sec:method-dicom}

\subsection{Series Identification}
\label{sec:method-series-identification}

Each study is first organised by DICOM \texttt{SeriesInstanceUID}. Candidate
series are classified into sagittal T1, sagittal T2/STIR and axial T2 using a
combination of DICOM metadata and image-orientation checks. Series descriptions
are informative but not trusted alone because naming conventions differ across
institutions. The classifier therefore considers fields such as
\texttt{SeriesDescription}, \texttt{ProtocolName}, sequence parameters where
available, and the orientation derived from \texttt{ImageOrientationPatient}.

Ambiguous or duplicated candidate series are flagged for manual review rather
than selected by filename order. The selected series identifiers are stored in
a manifest so that every experiment uses the same source series.

\subsection{Patient-Space Geometry}
\label{sec:method-geometry}

For each DICOM slice, the image plane is reconstructed from
\texttt{ImagePositionPatient}, \texttt{ImageOrientationPatient} and
\texttt{PixelSpacing}. Let $\mathbf{s}$ be the patient-space location of the
upper-left pixel, $\mathbf{r}$ and $\mathbf{c}$ the row and column direction
cosines, and $\Delta_r,\Delta_c$ the corresponding pixel spacings. The physical
coordinate of pixel $(i,j)$ is

\begin{equation}
\mathbf{x}(i,j)=
\mathbf{s}
+i\,\Delta_r\,\mathbf{r}
+j\,\Delta_c\,\mathbf{c}.
\label{eq:dicom-patient-coordinate}
\end{equation}

The slice normal is

\begin{equation}
\mathbf{n}=\mathbf{r}\times\mathbf{c}.
\label{eq:slice-normal}
\end{equation}

Slices are ordered by projection onto $\mathbf{n}$ rather than by filename.
Spacing between consecutive slice origins is inspected for discontinuities.
This procedure preserves patient coordinates across sagittal and axial
acquisitions, allowing a level localised in one plane to be projected into the
other.

For a sagittal target point $\mathbf{q}$ and candidate axial slice $k$ with
origin $\mathbf{s}_k$ and normal $\mathbf{n}_k$, the perpendicular distance to
that slice plane is

\begin{equation}
d_k =
\left|
\mathbf{n}_k^\top(\mathbf{q}-\mathbf{s}_k)
\right|.
\label{eq:plane-distance}
\end{equation}

The nearest anatomically corresponding axial slices are the slices with the
smallest $d_k$, subject to coverage and orientation checks. The method is an
implementation of physical correspondence, not a learned registration shortcut,
and follows the geometric principle used by prior lumbar pipelines
\cite{pan2021automatically,batra2025mscan}.

\subsection{Orientation Standardisation}
\label{sec:method-orientation}

Volumes are reoriented to a common anatomical convention for tensor processing
while retaining the transform to original patient coordinates. Left/right must
never be inferred after arbitrary array transposition; all laterality-dependent
targets are checked against the DICOM orientation transform. This is particularly
important for foraminal and subarticular labels, where a left/right inversion
would create a scientifically silent but catastrophic error.

\subsection{Voxel Spacing and Resampling}
\label{sec:method-resampling}

The acquisition spacing varies by institution, scanner and plane. Inputs are
therefore resampled only after anatomical coordinates have been reconstructed.
The target spacing is chosen from the training-set distribution so that it
preserves structures relevant to grading without creating unnecessary memory
cost. Sagittal and axial streams may use different in-plane resolutions because
their native geometry and diagnostic roles differ.

Interpolation uses a continuous method (e.g.\ linear or higher-order interpolation)
for MR intensities and nearest-neighbour interpolation for discrete masks or
labels. Original spacing is retained in metadata for any measurement reported
in millimetres. The exact target spacing will be \torecord{final sagittal and
axial resampling spacings selected from the development cohort}.

\subsection{Intensity Processing}
\label{sec:method-intensity}

MRI intensities are not quantitatively standardised across scanners. The
preprocessing therefore removes gross scale differences without imposing a
global histogram learned from test data. A default pipeline is:

\begin{enumerate}
  \item suppress non-finite values and obvious padding/background values;
  \item define a foreground mask from the body/spine region;
  \item clip extreme foreground intensities using training-derived robust
        percentiles;
  \item perform per-volume or per-ROI robust z-score normalisation;
  \item optionally compare with bias-field correction or histogram
        standardisation as a robustness ablation.
\end{enumerate}

Any learned harmonisation is fitted only on the development data. The external
Rizgary zero-shot result is reported both with the baseline preprocessing and,
if a separate harmonisation experiment is conducted, with that method clearly
labelled as an intervention rather than folded into the zero-shot definition.

\subsection{Quality-Control Flags}
\label{sec:method-qc}

Automated QC records the following per series:

\begin{itemize}
  \item slice count and ordering consistency;
  \item missing or duplicated instance positions;
  \item orientation consistency within series;
  \item in-plane spacing and slice spacing;
  \item field of view and lumbar coverage;
  \item intensity dynamic range;
  \item severe motion or truncation where detectable;
  \item presence of all sequences required by the particular experiment.
\end{itemize}

Cases failing geometry checks are not silently forced into the pipeline. They
are reviewed, excluded with a recorded reason, or assigned to a robustness
analysis if the failure itself represents a realistic missing-data condition.
A final inclusion/exclusion table will report every such decision.

% ===========================================================================
\section{Anatomical Localisation and ROI Construction}
\label{sec:method-localisation}

\subsection{Role of Localisation}
\label{sec:method-localisation-role}

Localisation exists to make the subsequent comparison fair. The thesis does not
claim that identifying L1--L2 through L5--S1 is itself novel; Chapter~2 showed
that this problem is mature and that strong systems already attain highly
reliable level assignment \cite{pan2021automatically}. The requirement here is
that the localiser be accurate enough that representation and graph experiments
are not dominated by wrong crops.

An established model is therefore preferred. Candidate implementations include
a SPIDER-pretrained nnU-Net-style segmentation network
\cite{isensee2021nnunet,vandergraaf2024lumbar}, a heatmap-regression detector,
or another state-of-the-art anatomical parser available when implementation
begins. The choice is made on validation localisation error and robustness, not
on whether it is architecturally fashionable.

\subsection{Outputs}
\label{sec:method-localisation-output}

The localiser produces, for each study:

\begin{itemize}
  \item centres of L1--L2 through L5--S1 in patient coordinates;
  \item vertebral/disc/canal masks or landmarks where the selected model provides
        them;
  \item a confidence or quality flag for each detected level;
  \item the transform between voxel coordinates and DICOM patient space.
\end{itemize}

A failed or low-confidence localisation is retained as a distinct failure mode.
It is not corrected manually for the primary automated pipeline without being
reported, because doing so would inflate downstream grading performance by
removing a real deployment error.

\subsection{Localisation Evaluation}
\label{sec:method-localisation-eval}

Where point annotations are available, centre localisation error is measured in
millimetres:

\begin{equation}
e_l = \|\hat{\mathbf{q}}_l-\mathbf{q}_l\|_2.
\label{eq:localisation-error}
\end{equation}

Mean/median error and percentile distributions are reported by level.
Percentage-of-correct-keypoint (PCK) or an equivalent tolerance-based measure is
also reported at pre-specified physical thresholds. Where full masks are
available, Dice and surface-distance metrics are reported. These results are
separated from disease grading, consistent with the literature showing that
detection/localisation and fine-grained grading generalise differently.

\subsection{Disease-Specific 2.5D Regions of Interest}
\label{sec:method-roi}

One rectangular crop is not assumed to be equally appropriate for every target.
The crop definition is conditioned on anatomical compartment.

For central canal stenosis, the model receives sagittal T2/STIR context around
the disc level together with corresponding axial T2 slices. For neural
foraminal narrowing, the sagittal T1 stream receives greater parasagittal
coverage and the ROI is side-aware, with sagittal T2 retained as complementary
context. For subarticular stenosis, axial T2 supplies the dominant local
evidence with sagittal context. These choices follow the clinical sequence--
pathology correspondence established in Chapter~2 rather than being learned
from arbitrary crop templates.

Each 2.5D stream contains a small ordered stack around the anatomically closest
slice,

\[
\mathcal{I}_l =
\{I_{k-r},\ldots,I_k,\ldots,I_{k+r}\},
\]

where $k$ is the level-centred slice and $r$ is a small radius determined on the
development set. A five-slice stack ($r=2$) is the initial reference
configuration because it captures through-plane context at modest memory cost;
the final radius is reported after sensitivity testing rather than treated as
a novel contribution.

Crops are defined in physical dimensions where possible, then resampled to the
network input resolution. This avoids a fixed 100-pixel crop representing
different anatomical widths on scanners with different pixel spacing.

\subsection{ROI Quality Control}
\label{sec:method-roi-qc}

A validation subset is visually inspected using overlays of the detected level,
crop boundary and corresponding axial slices. The inspection records:

\begin{itemize}
  \item correct lumbar level;
  \item inclusion of the relevant canal/foramen/recess;
  \item correct left/right orientation;
  \item adequate axial coverage;
  \item crop truncation;
  \item correspondence failure caused by unusual axial angulation.
\end{itemize}

The purpose is not to manually curate the full dataset but to detect systematic
geometry errors before they propagate into model training.

% ===========================================================================
\section{Baseline Representation and Sequence Encoders}
\label{sec:method-baseline}

\subsection{Independent ROI Classifier}
\label{sec:method-e0}

The mandatory baseline (E0) is an independent ROI classifier. Each target is
graded from its anatomically localised input without communication with other
targets. This represents the prevailing formulation criticised in
Chapter~\ref{ch:litreview} and provides the reference against which relational
reasoning is measured.

The baseline uses a well-characterised image encoder such as ResNet
\cite{he2016deep}, EfficientNet \cite{tan2019efficientnet}, Swin
\cite{liu2021swin}, or another contemporary backbone selected before final
experiments. The purpose is not to prove that one family is universally
superior. A primary backbone is selected for the ablation ladder so that
representation, routing and graph changes are measured under one stable
feature extractor. A second backbone may be used in a model-agnosticity check
for the strongest contribution.

\subsection{Sequence-Specific Feature Extraction}
\label{sec:method-sequence-encoders}

Each available MRI sequence is processed by a sequence-specific encoder
$E_m(\cdot)$,

\begin{equation}
\mathbf{f}_{p,l,m}=E_m(\mathcal{I}_{p,l,m}),
\label{eq:sequence-encoder}
\end{equation}

for patient $p$, level $l$ and modality/sequence $m$. The encoders may share
architecture while keeping separate normalisation statistics and weights, or
may partially share weights if an explicit ablation shows that doing so is
beneficial. Separate encoders are the default because T1, sagittal T2/STIR and
axial T2 have different contrast and orientation statistics.

To keep comparisons fair, the dimensionality of the output representation is
held constant across fusion methods. Fixed concatenation, simple averaging and
ordinary cross-attention all consume the same per-sequence features.

\subsection{Controlled Backbone Selection}
\label{sec:method-backbone-control}

Backbone choice is selected on the validation partition before the central
novelty experiments and then held fixed. The candidate comparison will use
identical crops, augmentations, optimiser budget and patient split. If one
backbone is selected because it gives the most stable validation performance
rather than the absolute highest point estimate, that decision is recorded.

This design prevents the thesis from becoming a repeated architecture search.
The main questions concern anatomical self-supervision, adaptive modality
allocation and target-level relational reasoning. A backbone is a controlled
feature extractor in those experiments.


% ===========================================================================
\section{Core Contribution I: Anatomically Aligned Cross-Sequence Self-Supervision}
\label{sec:method-ssl}

\subsection{Rationale}
\label{sec:method-ssl-rationale}

Conventional self-supervised vision learning defines related views by applying
augmentations to one image. That is useful but does not exploit a distinctive
property of multi-sequence MRI: sagittal T1, sagittal T2/STIR and axial T2 can
look very different while representing the same underlying patient and spinal
level. DICOM geometry supplies the correspondence directly.

The proposed objective therefore treats anatomical identity, rather than visual
similarity, as the primary source of positive pairs. For patient $p$, level
$l$, and two sequences $m_a$ and $m_b$, embeddings
$\mathbf{z}_{p,l,m_a}$ and $\mathbf{z}_{p,l,m_b}$ are encouraged to encode a
shared anatomical representation even though acquisition plane and contrast may
differ.

This idea is evaluated as a representation-learning hypothesis, not assumed to
be beneficial. The literature already shows that generic and longitudinal
self-supervision can reduce label dependence
\cite{jamaludin2017selfsupervised,chen2020simple}; the experiment here asks
whether explicitly anatomical cross-sequence correspondence is a stronger
signal for lumbar grading.

\subsection{Projection Head}
\label{sec:method-ssl-projection}

Each sequence feature is mapped through a projection head $P_m$:

\begin{equation}
\mathbf{z}_{p,l,m}=
\frac{
P_m(\mathbf{f}_{p,l,m})
}{
\|P_m(\mathbf{f}_{p,l,m})\|_2
}.
\label{eq:ssl-projection}
\end{equation}

The projection space is used for the self-supervised objective; the encoder
features before projection are retained for downstream grading. This prevents
the representation used for contrastive geometry from constraining the final
classifier unnecessarily.

\subsection{Positive-Pair Construction}
\label{sec:method-positive-pairs}

The strongest positive relation is:

\[
(p,l,m_a)\sim(p,l,m_b),
\qquad m_a\neq m_b,
\]

meaning the same patient and same anatomical level in two different sequences.
Correspondence is established through the DICOM reconstruction and level
localiser, not through nearest image appearance.

An optional softer relation may be tested for adjacent levels from the same
patient, because neighbouring levels share biomechanical context. It is not
included in the primary objective by default. Treating L3--L4 and L4--L5 as
equivalent positives would risk erasing genuine level-specific pathology, so
any adjacent-level term is isolated as a separate ablation.

Different patients at the same level are not defined as positives in the main
objective, because common anatomical level does not imply common disease state.
They may be used in a secondary supervised-contrastive variant once labels are
available, but this is conceptually distinct from the proposed anatomical
self-supervision.

\subsection{Contrastive Objective}
\label{sec:method-ssl-objective}

For anchor embedding $\mathbf{z}_i$ and its cross-sequence positive
$\mathbf{z}_{i^+}$, the reference InfoNCE-style loss is

\begin{equation}
\mathcal{L}_{\mathrm{ACSSL}}(i)
=
-\log
\frac{
\exp(\mathrm{sim}(\mathbf{z}_i,\mathbf{z}_{i^+})/\tau)
}{
\sum_{j\in\mathcal{B}\setminus\{i\}}
\exp(\mathrm{sim}(\mathbf{z}_i,\mathbf{z}_j)/\tau)
},
\label{eq:acssl}
\end{equation}

where $\mathrm{sim}$ is cosine similarity, $\tau$ is a temperature parameter
and $\mathcal{B}$ is the contrastive batch. When more than one anatomically
valid positive exists, the numerator is extended to the set of positives rather
than choosing one arbitrarily.

False negatives are a concern because two different patients may have very
similar normal anatomy. The implementation will therefore compare ordinary
InfoNCE with a variant that does not force all non-positive samples to be hard
negatives, such as a multi-positive or non-contrastive consistency objective.
The novelty claim is tied to the \emph{anatomical pairing}, not to a particular
brand of contrastive loss.

\subsection{Label-Efficiency Experiment}
\label{sec:method-label-efficiency}

The primary test of RQ2 uses fixed fractions of the labelled development set:
10\%, 25\%, 50\% and 100\%. For each fraction:

\begin{enumerate}
  \item the encoder is initialised from the candidate pretraining strategy;
  \item the same supervised grading head is fitted;
  \item patient splits and target distribution are held fixed;
  \item several random labelled subsets are evaluated where computationally
        feasible;
  \item macro-F1, weighted $\kappa$, severe-class recall and calibration are
        compared.
\end{enumerate}

The competing initialisations are:

\begin{itemize}
  \item training from random initialisation;
  \item ImageNet initialisation;
  \item generic medical/volumetric pretraining where a suitable public model is
        available;
  \item ordinary augmentation-based self-supervision;
  \item the proposed anatomically aligned cross-sequence self-supervision.
\end{itemize}

A successful result is not defined as beating every baseline at every label
fraction. The useful scientific question is where anatomical pairing helps,
where it does not, and whether any benefit is largest under label scarcity or
domain shift.

\subsection{Transfer Test of the Representation}
\label{sec:method-ssl-transfer}

The public-data model used for zero-shot Rizgary evaluation is trained once with
and once without anatomical self-supervision while the downstream architecture
is held constant. If the anatomically pretrained model loses less performance
externally despite comparable internal performance, that would support the
claim that the representation is less tied to dataset-specific appearance. If
both models degrade equally, the self-supervised contribution is limited to
label efficiency rather than domain robustness.

% ===========================================================================
\section{Core Contribution II: Disease-Conditioned Adaptive Sequence Routing}
\label{sec:method-routing}

\subsection{Problem Formulation}
\label{sec:method-routing-problem}

A conventional multi-sequence model applies the same fusion rule to every
target. That assumption is clinically implausible because the sequences are not
equally informative for each anatomical compartment. Sagittal T1 is particularly
useful for foraminal fat obliteration, sagittal T2/STIR for disc/canal context,
and axial T2 for canal and lateral recess morphology. The router turns that
known heterogeneity into a learnable, testable mechanism.

For patient $p$, target $t=(l,c)$ and available modalities
$\mathcal{M}_p$, let $\mathbf{f}_{p,l,m}$ be the feature from modality $m$.
The fused target feature is

\begin{equation}
\mathbf{u}_{p,t}
=
\sum_{m\in\mathcal{M}_p}
g_{p,t,m}\,
\mathbf{f}_{p,l,m},
\label{eq:routed-feature}
\end{equation}

where the non-negative routing weights satisfy

\begin{equation}
\sum_{m\in\mathcal{M}_p} g_{p,t,m}=1.
\label{eq:routing-simplex}
\end{equation}

The gate is conditioned on target identity and patient features:

\begin{equation}
g_{p,t,m}
=
\mathrm{softmax}_{m\in\mathcal{M}_p}
\left(
G[
\mathbf{f}_{p,l,m},
\mathbf{e}_{c},
\mathbf{e}_{l},
\mathbf{a}_{p}
]
\right),
\label{eq:routing-gate}
\end{equation}

where $\mathbf{e}_{c}$ and $\mathbf{e}_{l}$ are learnable condition and level
embeddings and $\mathbf{a}_{p}$ denotes optional acquisition/context metadata
that is available without leaking the target label. The initial implementation
uses a lightweight MLP or target-conditioned attention gate. A more complex
mixture-of-experts router is considered only if the simpler mechanism is
insufficient.

\subsection{Explicit Availability Mask}
\label{sec:method-routing-mask}

Missing sequences are represented by an availability mask
$a_{p,m}\in\{0,1\}$. Logits for unavailable modalities are masked before the
softmax:

\begin{equation}
g_{p,t,m}=0 \quad \mathrm{if}\quad a_{p,m}=0.
\label{eq:routing-missing-mask}
\end{equation}

The remaining weights are renormalised over the modalities that actually
exist. A missing sequence is therefore not represented by an all-zero image
whose meaning the network must infer. Availability is explicit in the model.

\subsection{Modality Dropout}
\label{sec:method-modality-dropout}

During training, complete studies are stochastically converted into incomplete
ones by removing one or more sequence streams. The dropout distribution includes:

\begin{itemize}
  \item all three sequences;
  \item sagittal T1 + sagittal T2;
  \item sagittal T2 + axial T2;
  \item sagittal T1 + axial T2 where clinically meaningful;
  \item each single-sequence configuration as an edge case.
\end{itemize}

The sampling probabilities are tuned so that the model sees enough complete
studies to learn complementary fusion while also receiving substantial exposure
to incomplete inputs. The exact probabilities are \torecord{final modality-
dropout schedule selected on validation data}.

No external test case is artificially repaired by copying another sequence into
a missing slot. The model's purpose is to remain functional under genuine
absence, not to conceal it.

\subsection{Routing Consistency}
\label{sec:method-routing-consistency}

An optional consistency term encourages the prediction from an incomplete
modality set to remain close to the complete-study prediction when the removed
modality is non-essential:

\begin{equation}
\mathcal{L}_{\mathrm{miss}}
=
D\!\left(
p_{\theta}(y\mid x_{\mathrm{full}}),
p_{\theta}(y\mid x_{\mathrm{subset}})
\right),
\label{eq:missing-consistency}
\end{equation}

where $D$ is a divergence such as KL divergence or Jensen--Shannon divergence.
This term is not applied blindly: for targets known to depend strongly on a
specific sequence, forcing invariance to its removal could suppress legitimate
diagnostic information. The term is therefore ablated by target and compared
with modality-dropout training alone.

\subsection{Routing Baselines}
\label{sec:method-routing-baselines}

RQ3 is evaluated against four progressively stronger controls:

\begin{enumerate}
  \item fixed feature concatenation;
  \item simple equal-weight averaging;
  \item ordinary cross-attention over available sequence features;
  \item disease-conditioned routing without modality dropout;
  \item disease-conditioned routing with modality dropout.
\end{enumerate}

For each method, encoder capacity, ROI inputs and training budget are matched as
closely as possible. Missing-sequence robustness is evaluated on the same
patients under the same artificial absence pattern.

\subsection{Interpreting Routing Weights}
\label{sec:method-routing-interpretation}

The gate weights are model allocations, not causal estimates of sequence
importance. A high axial T2 weight does not prove that axial T2 caused a correct
decision. Clinical plausibility is assessed in two complementary ways.

First, aggregate routing weights are summarised by target. The expected pattern
is that foraminal targets allocate more weight to sagittal T1 and subarticular/
canal targets more to axial T2, but this expectation is not built into the gate
as a hard prior.

Second, controlled input ablation tests whether removing the sequence with high
routing weight causes a greater performance loss than removing a low-weight
sequence. Agreement between learned allocation and intervention strengthens the
interpretation; disagreement is reported rather than rationalised post hoc.

% ===========================================================================
\section{Core Contribution III: Heterogeneous Disease--Anatomy Graph}
\label{sec:method-graph}

\subsection{Graph Construction}
\label{sec:method-graph-construction}

The proposed graph represents \emph{prediction targets}, not pixels and not a
3D disc surface. At full LumbarDISC scope, the node set is

\begin{equation}
V=\{v_{l,c}:l\in\mathcal{L},\,c\in\mathcal{C}\},
\qquad |V|=25.
\label{eq:graph-nodes}
\end{equation}

Each node begins with the routed visual representation for that target together
with embeddings that identify level and condition:

\begin{equation}
\mathbf{h}^{(0)}_{l,c}
=
[
\mathbf{u}_{l,c}
\Vert
\mathbf{e}_{l}
\Vert
\mathbf{e}_{c}
].
\label{eq:node-initial}
\end{equation}

The graph topology is defined anatomically before training. It is not learned
from correlations in the test data.

\subsection{Typed Edge Families}
\label{sec:method-edge-types}

Three primary edge types are used.

\textbf{Adjacent-level edges} connect the same target type across neighbouring
lumbar levels, for example

\[
v_{L3-L4,\mathrm{Canal}}
\leftrightarrow
v_{L4-L5,\mathrm{Canal}}.
\]

These encode the ordered chain of lumbar motion segments.

\textbf{Same-level cross-condition edges} connect anatomical compartments that
coexist at one level, for example central canal to left/right foraminal and
subarticular targets. The edge expresses that the targets share local
morphology, not that one grade deterministically implies another.

\textbf{Bilateral edges} connect left and right counterparts of the same
condition, for example

\[
v_{L4-L5,\mathrm{LF}}
\leftrightarrow
v_{L4-L5,\mathrm{RF}}.
\]

This represents bilateral anatomical symmetry without forcing equal grades.

Additional cross-type edges are added only if clinically justified in the
pre-specified topology. A denser graph is not automatically a better graph; the
study tests whether explicit relation types add value over homogeneous message
passing.

\subsection{Relation-Aware Message Passing}
\label{sec:method-message-passing}

For layer $q$, node $i$ receives messages from neighbours $j\in\mathcal{N}(i)$.
A relation-aware attention formulation is

\begin{align}
\mathbf{q}_i &= W_Q\mathbf{h}^{(q)}_i,\\
\mathbf{k}_{j,r} &= W^r_K\mathbf{h}^{(q)}_j,\\
\mathbf{v}_{j,r} &= W^r_V\mathbf{h}^{(q)}_j,
\end{align}

with attention

\begin{equation}
\alpha_{ijr}
=
\mathrm{softmax}_{j,r}
\left(
\frac{
\mathbf{q}_i^\top\mathbf{k}_{j,r}
}{
\sqrt{d}
}
+b_r
\right),
\label{eq:relational-attention}
\end{equation}

and update

\begin{equation}
\tilde{\mathbf{h}}^{(q+1)}_i
=
\sum_{(j,r)\in\mathcal{N}(i)}
\alpha_{ijr}\mathbf{v}_{j,r}.
\label{eq:graph-message}
\end{equation}

Residual connection, normalisation and feed-forward transformation produce the
final layer output. Relation-specific transformations $W^r_K,W^r_V$ allow an
adjacent-level neighbour to influence a node differently from its contralateral
counterpart.

The exact graph transformer implementation may be replaced by a computationally
simpler relation-aware GAT if validation shows no benefit from the more complex
version. The contribution is the explicit typed target graph and its controlled
evaluation, not a requirement to use the largest graph architecture.

\subsection{Graph Baselines}
\label{sec:method-graph-baselines}

The graph contribution is tested against:

\begin{enumerate}
  \item \textbf{Independent heads:} no message passing between targets.
  \item \textbf{Ordered level transformer:} features communicate across the five
        lumbar levels but do not form 25 typed condition/laterality nodes.
  \item \textbf{Homogeneous graph:} the same 25 nodes and adjacency pattern but
        one undifferentiated edge type.
  \item \textbf{Heterogeneous graph:} typed adjacent-level, cross-condition and
        bilateral relations.
  \item \textbf{Random/shuffled-edge control:} graph capacity retained while
        anatomical edges are permuted.
\end{enumerate}

The shuffled-edge control is essential. If a graph with arbitrary topology
performs as well as the anatomical graph, the study has shown that extra
message-passing capacity helps, not that the encoded anatomy matters.

\subsection{Edge-Family Ablation}
\label{sec:method-edge-ablation}

The full graph is further decomposed:

\begin{itemize}
  \item remove adjacent-level edges;
  \item remove same-level cross-condition edges;
  \item remove bilateral edges;
  \item retain only one family at a time.
\end{itemize}

Results are reported by target family. It is plausible, for example, that
bilateral edges help lateral-compartment grading but add little to central-canal
classification. A single average graph improvement would conceal this structure.

\subsection{Graph Consistency and Uncertainty}
\label{sec:method-graph-uncertainty}

Graph information is not used to hard-code clinically deterministic rules.
Severe stenosis at L4--L5 does not require moderate stenosis at L3--L4, and the
model is allowed to preserve isolated pathology. Relational message passing is
therefore soft.

For uncertainty analysis, node disagreement can be measured between an
independent-head prediction and the graph-refined prediction, or between
neighbour-conditioned and local evidence. This disagreement is explored as an
auxiliary signal for selective prediction, but it is not equated with calibrated
clinical uncertainty unless validated against error.

% ===========================================================================
\section{Supporting Method: Ordinal and Cost-Sensitive Grading}
\label{sec:method-ordinal}

\subsection{Categorical Baseline}
\label{sec:method-ce}

Standard three-class cross-entropy is retained as the mandatory reference:

\begin{equation}
\mathcal{L}_{\mathrm{CE}}
=
-\sum_{t}m_t
\sum_{k=0}^{2}
\mathbb{I}(y_t=k)\log p_{t,k}.
\label{eq:cross-entropy}
\end{equation}

Class weighting or balanced sampling is applied in separate controlled variants.
This baseline is necessary because previous lumbar work found that explicitly
ordinal objectives did not consistently outperform cross-entropy
\cite{niemeyer2021a}.

\subsection{Ordinal Threshold Formulation}
\label{sec:method-ordinal-threshold}

An ordinal head represents a three-grade target as two ordered threshold
questions:

\[
P(y>0\mid\mathbf{h}), \qquad
P(y>1\mid\mathbf{h}).
\]

A CORAL/CORN-style implementation may be used provided the ordering constraint
is respected. The exact ordinal formulation is selected before the final
ablation and compared directly with categorical cross-entropy under identical
features.

Performance is evaluated not only by exact class agreement but also by error
distance:

\begin{equation}
d_t=|\hat{y}_t-y_t|.
\label{eq:grade-distance}
\end{equation}

The proportions with $d_t=0$, $d_t=1$ and $d_t=2$ are reported separately.

\subsection{Clinically Asymmetric Cost}
\label{sec:method-cost}

A cost matrix $C$ allows distant under-grading to receive a larger training
penalty:

\begin{equation}
C=
\begin{bmatrix}
0 & c_{01} & c_{02}\\
c_{10} & 0 & c_{12}\\
c_{20} & c_{21} & 0
\end{bmatrix},
\label{eq:cost-matrix}
\end{equation}

with the clinically motivated condition

\[
c_{20}>c_{21},
\]

so that Severe$\rightarrow$Normal/Mild is more costly than
Severe$\rightarrow$Moderate. The exact numerical values are not invented by the
model developer. They are either aligned with a defensible benchmark/clinical
weighting or selected through a sensitivity analysis across a small
pre-specified family of matrices.

One implementation uses the expected cost under the predicted distribution:

\begin{equation}
\mathcal{L}_{\mathrm{cost}}
=
\sum_t m_t
\sum_{k=0}^{2}
C_{y_t,k}\,p_{t,k}.
\label{eq:expected-cost}
\end{equation}

The cost term is evaluated alongside, not instead of, standard calibration and
discrimination metrics. A model is not considered better merely because it
optimises the chosen matrix.

\subsection{Class Imbalance}
\label{sec:method-imbalance}

Because severe grades are uncommon in LumbarDISC, the study compares a small
set of established strategies rather than stacking all of them simultaneously:

\begin{itemize}
  \item unweighted cross-entropy;
  \item inverse-frequency or effective-number class weighting;
  \item weighted sampling;
  \item focal loss \cite{lin2017focal};
  \item the cost-sensitive/ordinal variant above.
\end{itemize}

MixUp/CutMix are not the default for ordinal grading because mixed labels may
have no clinical interpretation. Geometric augmentation is constrained so that
left/right labels are updated consistently; horizontal flipping is prohibited
for laterality-dependent targets unless labels and orientation are transformed
together.

% ===========================================================================
\section{Supporting Method: Calibration, Uncertainty and Selective Prediction}
\label{sec:method-calibration}

\subsection{Probability Calibration}
\label{sec:method-temperature}

Raw neural-network probabilities are not treated as calibrated confidence.
Temperature scaling is fitted on the validation set:

\begin{equation}
p_k^{(T)}
=
\frac{
\exp(z_k/T)
}{
\sum_j \exp(z_j/T)
},
\label{eq:temperature}
\end{equation}

where $T>0$ is selected without access to the test labels. Calibration is then
evaluated with Expected Calibration Error, Brier score and reliability diagrams.

\subsection{Predictive Uncertainty}
\label{sec:method-uncertainty}

The primary feasible uncertainty methods are MC dropout and/or a small deep
ensemble. If $S$ stochastic predictions are available, the predictive mean is

\begin{equation}
\bar{\mathbf{p}}
=
\frac{1}{S}\sum_{s=1}^{S}\mathbf{p}^{(s)},
\label{eq:predictive-mean}
\end{equation}

and uncertainty can be summarised by predictive entropy,

\begin{equation}
H(\bar{\mathbf{p}})
=
-\sum_k \bar{p}_k\log\bar{p}_k,
\label{eq:entropy}
\end{equation}

together with between-member variance. The study reports whether uncertainty is
higher on incorrect predictions, distant grade errors, low-agreement targets
and external-domain cases.

\subsection{Selective Prediction}
\label{sec:method-selective}

A selective system abstains when uncertainty exceeds threshold $\delta$:

\[
U(x)>\delta \Rightarrow \mathrm{refer/abstain}.
\]

Rather than choosing one arbitrary threshold, the complete risk--coverage curve
is reported: as increasingly uncertain cases are withheld, what fraction of the
cohort remains covered and what error remains among the retained predictions?
The area under, or summary points from, this curve can compare methods.

The terminology ``refer'' describes model behaviour only. It does not claim a
validated clinical workflow; a prospective reader study would be required to
show that abstention improves patient care.

% ===========================================================================
\section{Composite Training Objective}
\label{sec:method-total-loss}

The complete model may use the objective

\begin{equation}
\mathcal{L}_{\mathrm{total}}
=
\mathcal{L}_{\mathrm{grade}}
+\lambda_{\mathrm{ssl}}\mathcal{L}_{\mathrm{ACSSL}}
+\lambda_{\mathrm{miss}}\mathcal{L}_{\mathrm{miss}}
+\lambda_{\mathrm{cost}}\mathcal{L}_{\mathrm{cost}}
+\lambda_{\mathrm{reg}}\mathcal{L}_{\mathrm{reg}},
\label{eq:total-loss}
\end{equation}

where $\mathcal{L}_{\mathrm{grade}}$ is categorical or ordinal supervised loss,
$\mathcal{L}_{\mathrm{ACSSL}}$ is the anatomical cross-sequence objective,
$\mathcal{L}_{\mathrm{miss}}$ is optional missing-modality consistency,
$\mathcal{L}_{\mathrm{cost}}$ is the optional clinical cost term and
$\mathcal{L}_{\mathrm{reg}}$ collects standard regularisation.

Graph message passing is normally part of the forward model rather than a
separate loss. If an explicit graph-consistency regulariser is introduced, it
receives its own coefficient and ablation.

Every non-zero $\lambda$ is treated as a hyperparameter whose contribution must
be demonstrated. A loss term is removed if it adds no reproducible value. The
final thesis reports the selected values and the search range:
\torecord{$\lambda$ values, temperature, optimiser and regularisation settings}.

% ===========================================================================
\section{Optimisation and Training Protocol}
\label{sec:method-training}

\subsection{Training Phases}
\label{sec:method-training-phases}

Training is separated into conceptually distinct phases:

\begin{enumerate}
  \item \textbf{Anatomical/localisation training or loading.}
        The localiser is trained/reproduced and then fixed for the grading
        experiments unless a joint-training ablation is explicitly conducted.

  \item \textbf{Self-supervised pretraining.}
        Sequence encoders are trained with anatomical cross-sequence pairs using
        development patients only.

  \item \textbf{Supervised grading.}
        Encoders are fine-tuned with the routing and graph modules on
        LumbarDISC labels.

  \item \textbf{Calibration.}
        Temperature/uncertainty parameters are fitted on validation data after
        model selection.

  \item \textbf{External evaluation.}
        The selected model is frozen and evaluated on the fixed Rizgary test
        set.

  \item \textbf{Few-shot adaptation.}
        New model copies are adapted using only the local adaptation pool and
        re-evaluated on the same fixed local test set.
\end{enumerate}

\subsection{Optimiser and Schedules}
\label{sec:method-optimiser}

AdamW is the default optimiser for transformer/graph components; SGD or AdamW
may be compared for CNN baselines if required. Learning rate is selected on the
validation set within a pre-specified range, with warm-up and cosine decay or a
plateau scheduler used consistently across the main ablation. Early stopping is
based on a prevalence-robust validation metric such as macro-F1 or weighted
$\kappa$, not raw accuracy.

The final report records:

\begin{itemize}
  \item optimiser and version;
  \item initial learning rate and scheduler;
  \item batch size / effective batch size;
  \item number of epochs and early-stopping patience;
  \item weight decay and dropout;
  \item input resolution and slice-stack radius;
  \item precision mode (FP32/BF16/FP16);
  \item random seeds;
  \item hardware and software versions;
  \item wall-clock and GPU time.
\end{itemize}

No final hyperparameter value is fabricated in this protocol chapter; each is
\torecord{executed training configuration and model-selection criterion}.

\subsection{Augmentation}
\label{sec:method-augmentation}

Training augmentation is designed to simulate plausible MRI variability without
altering disease anatomy. Candidate transforms include modest intensity scaling,
gamma/contrast change, bias-field simulation, noise, small translation/rotation
and limited elastic deformation. Aggressive transformations that distort the
canal or foraminal morphology beyond plausible acquisition variation are
excluded.

Laterality-aware augmentation updates all side-specific labels. If a horizontal
flip changes anatomical left/right, left and right target labels and graph node
identity are swapped in the same transformation. Augmentation is disabled for
validation and test data except in a separately labelled test-time-augmentation
experiment.

\subsection{Model Selection}
\label{sec:method-model-selection}

Only validation data are used to choose architecture depth, routing
hyperparameters, loss weights, early stopping and calibration. The held-out
public test and fixed Rizgary test sets are not inspected for iterative model
selection. If many exploratory experiments are performed, only the pre-specified
final comparison is used for confirmatory statistical claims; exploratory
results are labelled as such.


% ===========================================================================
\section{Experimental Ladder and Ablation Programme}
\label{sec:method-ablation}

The main experiment is not one comparison between ``AMOG-Net'' and a baseline.
It is a ladder in which each additional assumption is isolated.

\begin{longtable}{p{1.1cm}p{5.0cm}p{7.2cm}}
\toprule
\textbf{ID} & \textbf{Configuration} & \textbf{Question isolated}\\
\midrule
\endfirsthead
\toprule
\textbf{ID} & \textbf{Configuration} & \textbf{Question isolated}\\
\midrule
\endhead
E0 & Independent ROI classifier & How well can anatomically localised targets be graded without relational or adaptive fusion?\\
E1 & + DICOM-aligned multi-sequence ROIs & What is gained by correct cross-plane correspondence under fixed fusion?\\
E2 & + disease-conditioned sequence routing & Does target-specific routing outperform fixed fusion?\\
E3 & + modality dropout / missing-sequence training & Does explicit training for absence improve robustness?\\
E4 & + anatomical cross-sequence SSL & Does anatomy-defined pretraining improve representation and label efficiency?\\
E5 & + homogeneous target graph & Does generic message passing help?\\
E6 & + typed heterogeneous disease--anatomy graph & Do clinically defined edge types matter beyond graph capacity?\\
E7 & + ordinal/cost-sensitive/calibrated heads & Can clinically consequential errors and probability reliability be improved?\\
E8 & Selected complete system + external transfer & Do the selected components survive institutional shift and adapt efficiently?\\
\bottomrule
\end{longtable}

The ordering may be adjusted for engineering practicality, but the comparisons
themselves remain. In particular, E6 must be compared with both E5 and a
shuffled-edge control. Otherwise any improvement can be attributed to additional
parameters rather than anatomical topology.

\subsection{Contribution-Specific Ablations}
\label{sec:method-contribution-ablation}

Three ablation families correspond directly to the doctoral contributions.

\textbf{Anatomical SSL ablation} compares no SSL, generic image SSL and
DICOM-aligned cross-sequence SSL under identical downstream training. It is
repeated at reduced label fractions.

\textbf{Routing ablation} compares fixed fusion, cross-attention, target-
conditioned routing and target-conditioned routing with modality dropout.
Performance is reported under complete and missing-sequence input.

\textbf{Graph ablation} compares independent heads, an ordered five-level
transformer, homogeneous graph, heterogeneous graph and random/shuffled edges.
Edge-family removal further identifies whether adjacency, bilateral symmetry or
same-level cross-condition interaction carries the signal.

\subsection{Negative Results}
\label{sec:method-negative-results}

A component that fails to improve the pre-specified outcome is not silently
removed from the research record. The thesis reports it, including confidence
intervals and target-specific effects. This is particularly important for
ordinal objectives and graph modelling, because both are theoretically
attractive and therefore vulnerable to confirmation bias.

A component enters the final integrated system only if it adds a reproducible
benefit, improves an outcome that the thesis has defined as clinically relevant,
or yields a meaningful robustness/calibration advantage despite similar
aggregate discrimination.

% ===========================================================================
\section{External Validation at Rizgary}
\label{sec:method-external}

\subsection{Definition of Zero-Shot Transfer}
\label{sec:method-zero-shot}

Zero-shot transfer is defined strictly. The selected model, including encoder,
router, graph, classifier and calibration strategy developed on public data, is
frozen before any Rizgary test label is used. The local images are passed
through the same preprocessing and localisation pipeline, with only
non-learned DICOM reconstruction and pre-specified per-volume normalisation
allowed.

No local fine-tuning, threshold optimisation or calibration is permitted in the
zero-shot result. If a local intensity harmonisation method is fitted on Rizgary
data, that constitutes a separate domain-adaptation experiment and is labelled
accordingly.

The primary zero-shot analysis is restricted to the central-canal targets for
which the local reports provide a defensible mapping. The principal comparison
is therefore not the full 25-target benchmark score versus a five-target local
score; it is public held-out versus local external performance on the
\emph{same central-canal task}.

\subsection{Quantifying Domain Shift}
\label{sec:method-domain-shift}

Performance degradation is reported as both absolute and relative change:

\begin{align}
\Delta M &= M_{\mathrm{Rizgary}}-M_{\mathrm{Public}},\\
\Delta M_{\%} &=
100\,
\frac{M_{\mathrm{Rizgary}}-M_{\mathrm{Public}}}
{M_{\mathrm{Public}}},
\end{align}

for each primary metric $M$. The thesis does not interpret a large negative
$\Delta M$ as failed research; the size and structure of the drop are among the
main outcomes of RQ4.

The shift is decomposed by:

\begin{itemize}
  \item lumbar level;
  \item severity class;
  \item available sequence configuration;
  \item scanner/vendor/field strength where the local cohort contains enough
        variation;
  \item slice thickness/in-plane resolution;
  \item localisation success;
  \item report/reference-standard confidence.
\end{itemize}

This analysis separates representational failure from data-quality failure
where possible. For example, if grading declines but localisation remains
stable, the pattern supports the detection-versus-grading asymmetry documented
in Chapter~2.

\subsection{Local Error Review}
\label{sec:method-local-error-review}

A blinded error audit is performed after the zero-shot metrics are locked.
Incorrect local predictions are grouped into:

\begin{itemize}
  \item wrong level/localisation;
  \item adjacent-grade disagreement;
  \item distant Severe$\leftrightarrow$Normal/Mild disagreement;
  \item apparent label/report ambiguity;
  \item missing or atypical sequence;
  \item image-quality/artifact issue;
  \item plausible model error despite adequate evidence.
\end{itemize}

Where reader time permits, a subset of model/report disagreements is re-read
without exposing the model prediction. The purpose is not to alter the test
labels opportunistically but to estimate how much apparent model error may
reflect reference-standard variability.

% ===========================================================================
\section{Few-Shot Domain Adaptation}
\label{sec:method-adaptation}

\subsection{Adaptation Sizes}
\label{sec:method-adaptation-sizes}

The adaptation pool is sampled at

\[
N\in\{10,25,50,100\}
\]

patient-level labelled cases, subject to the final number of eligible local
studies. Each $N$ is sampled multiple times, stratified where feasible by
central-canal severity. The same fixed local test set is used after every
adaptation.

The experiment therefore estimates a function

\begin{equation}
M(N)=
\text{test performance after adaptation with }N\text{ local cases},
\label{eq:adaptation-curve}
\end{equation}

rather than a single ``fine-tuned'' number.

\subsection{Adaptation Strategies}
\label{sec:method-adaptation-strategies}

Four strategies are compared:

\begin{enumerate}
  \item \textbf{No adaptation.} The zero-shot baseline.
  \item \textbf{Non-learned/simple harmonisation.} Intensity preprocessing is
        adjusted using only the adaptation pool where a defensible method exists.
  \item \textbf{Parameter-efficient adaptation.} Adapter layers or low-rank
        updates are trained while most of the public-data representation remains
        fixed.
  \item \textbf{Full fine-tuning.} All or most parameters are updated, where
        computationally feasible and where $N$ is large enough to avoid
        immediate overfitting.
\end{enumerate}

The exact PEFT implementation is chosen after a small validation comparison and
is not itself claimed as novel. Its role is to ask an operational question:
how many local labels are needed when a hospital cannot afford to retrain a full
model from scratch?

\subsection{Annotation-Efficiency Summary}
\label{sec:method-annotation-efficiency}

For each method, the recovery curve reports macro-F1, weighted $\kappa$, severe
recall and calibration as a function of $N$. A useful summary is the area under
the performance-versus-log-annotation curve, but raw points with confidence
intervals remain primary because the curve may be non-linear.

No arbitrary threshold such as ``95\% recovery'' defines success. If 100 local
cases recover little performance, that is a meaningful negative finding about
transportability. If 10 cases recover most of the loss, that is evidence for
annotation-efficient adaptation.

\subsection{Avoiding Adaptation Overfit}
\label{sec:method-adaptation-overfit}

Few-shot training is particularly vulnerable to overfitting and selection bias.
The study therefore:

\begin{itemize}
  \item repeats each $N$ across several sampled adaptation sets;
  \item uses regularisation and early stopping based only on an adaptation
        validation split or internal cross-validation within the adaptation
        pool;
  \item never selects the adaptation epoch from the fixed external test;
  \item reports the variance across sampled local subsets;
  \item compares PEFT with full fine-tuning rather than assuming fewer trainable
        parameters are always superior.
\end{itemize}

% ===========================================================================
\section{Evaluation Metrics}
\label{sec:method-metrics}

No single metric is treated as sufficient because the task is imbalanced,
ordinal and multi-target.

\subsection{Macro F1 and Balanced Accuracy}
\label{sec:method-macro}

For class $k$, precision and recall are

\begin{align}
P_k &= \frac{TP_k}{TP_k+FP_k},\\
R_k &= \frac{TP_k}{TP_k+FN_k},
\end{align}

and

\begin{equation}
F1_k=
2\frac{P_kR_k}{P_k+R_k}.
\end{equation}

Macro-F1 is the unweighted class mean,

\begin{equation}
F1_{\mathrm{macro}}
=
\frac{1}{K}\sum_{k=1}^{K}F1_k,
\label{eq:macro-f1}
\end{equation}

so the severe class cannot be hidden by the Normal/Mild majority.

Balanced accuracy is

\begin{equation}
BA=
\frac{1}{K}\sum_{k=1}^{K}R_k.
\label{eq:balanced-accuracy}
\end{equation}

These are primary discrimination metrics.

\subsection{Quadratic Weighted Kappa}
\label{sec:method-kappa}

Quadratic weighted $\kappa$ is reported because the labels are ordered and the
literature supplies human agreement values in $\kappa$. For observed confusion
matrix $O$ and expected matrix $E$, with weight

\[
w_{ij}=\frac{(i-j)^2}{(K-1)^2},
\]

the statistic is

\begin{equation}
\kappa_w
=
1-
\frac{
\sum_{ij}w_{ij}O_{ij}
}{
\sum_{ij}w_{ij}E_{ij}
}.
\label{eq:weighted-kappa}
\end{equation}

Exact agreement and within-one-grade agreement are also reported so that a high
$\kappa$ is not the only view of ordinal error.

\subsection{Severe-Class Outcomes}
\label{sec:method-severe}

The clinically decisive class is reported separately:

\begin{itemize}
  \item Severe recall/sensitivity;
  \item Severe precision;
  \item Severe-versus-rest AUROC and AUPRC where appropriate;
  \item Severe$\rightarrow$Normal/Mild error rate;
  \item Severe$\rightarrow$Moderate error rate.
\end{itemize}

The separation distinguishes distant under-grading from adjacent boundary error.

\subsection{AUROC and AUPRC}
\label{sec:method-auroc}

For three-class outcomes, AUROC is calculated one-versus-rest and macro-averaged.
AUPRC is reported alongside AUROC for rare Severe classes because it is more
sensitive to prevalence and false positives. Curves are constructed at patient-
target level, while confidence intervals respect patient clustering.

\subsection{Calibration}
\label{sec:method-calibration-metrics}

Calibration evaluation includes:

\begin{itemize}
  \item Brier score;
  \item Expected Calibration Error (ECE);
  \item classwise reliability diagrams;
  \item negative log likelihood where useful;
  \item risk--coverage curves under selective prediction.
\end{itemize}

Calibration is reported both internally and externally because a model can
retain discrimination while becoming overconfident after domain shift.

\subsection{Localisation Metrics}
\label{sec:method-localisation-metrics}

Localisation is evaluated separately using:

\begin{itemize}
  \item mean, median and 95th-percentile Euclidean error in millimetres;
  \item PCK at physically meaningful tolerances;
  \item Dice coefficient for available masks;
  \item surface-distance metrics where mask quality supports them;
  \item complete localisation-failure rate per study.
\end{itemize}

\subsection{Efficiency Metrics}
\label{sec:method-efficiency}

For each model configuration the thesis reports:

\begin{itemize}
  \item trainable and total parameter count;
  \item FLOPs/MACs where they can be estimated consistently;
  \item peak VRAM;
  \item training wall time and GPU hours;
  \item inference time per study and per target;
  \item storage size of weights.
\end{itemize}

Efficiency comparisons use the same hardware and batch/inference conditions.
Timing from different machines is not compared as if it were an architectural
property.

% ===========================================================================
\section{Statistical Analysis}
\label{sec:method-statistics}

\subsection{Patient-Level Bootstrap Confidence Intervals}
\label{sec:method-bootstrap}

Confidence intervals for primary metrics are obtained by bootstrap resampling
\emph{patients}, not individual targets. On each replicate, a patient is sampled
with replacement and all of that patient's target predictions enter together.
This preserves within-study dependence.

The default is a sufficiently large number of replicates to stabilise the
2.5th and 97.5th percentiles; the final number is
\torecord{bootstrap replicate count}. Ninety-five per cent percentile or
bias-corrected intervals are reported consistently.

\subsection{Paired Model Comparisons}
\label{sec:method-paired-tests}

Ablation models are evaluated on the same patients, so comparisons are paired.
For scalar patient-level losses or correctness indicators, a paired permutation
test or Wilcoxon signed-rank test is used where distributional assumptions are
not justified. AUCs are compared with an appropriate paired procedure such as
DeLong's test when its assumptions are satisfied. Bootstrap differences are
used for metrics without a simple analytic test.

The primary quantity is the paired difference

\[
\Delta M=M_A-M_B
\]

with a 95\% confidence interval. Point estimates alone are not used to claim
one component is superior.

\subsection{Multiple Comparisons}
\label{sec:method-multiple}

The ablation programme contains many target-level analyses. A limited set of
primary comparisons is therefore declared in advance:

\begin{enumerate}
  \item anatomical SSL versus generic/self-supervised baseline;
  \item routed+dropout fusion versus ordinary cross-attention;
  \item heterogeneous graph versus independent heads and homogeneous graph;
  \item zero-shot public versus local performance;
  \item PEFT versus no adaptation across annotation sizes.
\end{enumerate}

Target-by-target and subgroup analyses are secondary/exploratory. Where many
formal hypotheses are tested, false-discovery-rate control is applied. The
thesis distinguishes confirmatory comparisons from exploratory patterns instead
of presenting every $p<0.05$ as independent evidence.

\subsection{Repeated Training Seeds}
\label{sec:method-seeds}

Deep-learning optimisation is stochastic. Final ablation comparisons are
therefore repeated across multiple fixed seeds where computationally feasible.
The thesis reports mean and standard deviation/interval across seeds in addition
to the patient-level test uncertainty. A one-off training run is insufficient
to attribute a small improvement to a method.

\subsection{Subgroup Analysis}
\label{sec:method-subgroups}

Subgroup analysis is conducted only where sample size permits and includes:

\begin{itemize}
  \item lumbar level;
  \item severity;
  \item sex/age bands where metadata and governance allow;
  \item scanner/vendor/field strength;
  \item acquisition resolution or slice thickness;
  \item presence/absence of individual sequences.
\end{itemize}

Sparse cells are reported descriptively rather than subjected to unstable
significance testing. The subgroup analysis is intended to expose failure modes,
not to make causal claims about demographic biology.

% ===========================================================================
\section{Robustness Experiments}
\label{sec:method-robustness}

\subsection{Missing-Sequence Matrix}
\label{sec:method-missing-matrix}

The selected public-data models are evaluated under every supported sequence
configuration. For three sequences, the core matrix includes complete input,
three two-sequence combinations and three single-sequence configurations where
the required anatomical target remains observable.

Performance is reported per target family rather than only in aggregate. A
model can remain strong on central canal grading while collapsing on foraminal
narrowing when sagittal T1 is removed; an overall mean would obscure the
clinical meaning.

\subsection{Geometry Perturbation}
\label{sec:method-geometry-robust}

The reliance on DICOM coordinates creates a specific failure mode. Controlled
perturbations therefore test sensitivity to small localisation/geometry errors:
axial slice selection shifted by one or two slices, modest ROI-centre offsets,
and simulated spacing variation. The purpose is not to train on corrupted
geometry by default, but to measure how brittle the system is to the imperfect
localisation expected in practice.

\subsection{Image-Quality and Acquisition Perturbation}
\label{sec:method-acquisition-robust}

Where feasible, synthetic intensity perturbations approximating scanner shift
are applied to the public held-out set: contrast scaling, noise, bias field and
resolution degradation. These do not replace real external validation but help
identify which changes reproduce the pattern observed at Rizgary.

\subsection{Graph Perturbation}
\label{sec:method-graph-robust}

The graph is challenged by:

\begin{itemize}
  \item random edge deletion;
  \item shuffled edge types;
  \item random rewiring with matched edge count;
  \item removal of one relation family.
\end{itemize}

A robust anatomical graph should outperform these controls without depending on
one fragile edge. If random topology performs equivalently, the anatomical
interpretation is rejected.

% ===========================================================================
\section{Interpretability and Failure Analysis}
\label{sec:method-interpretability}

\subsection{Local Visual Evidence}
\label{sec:method-visual-evidence}

Grad-CAM or equivalent saliency may be used as a gross sanity check that the
encoder responds within the relevant ROI, but saliency is not presented as a
proof of reasoning. The stronger explanations in this thesis arise from
structural quantities that are part of the model itself:

\begin{itemize}
  \item which sequence received routing weight for a specific target;
  \item which graph relation contributed to target refinement;
  \item how uncertainty changed before and after relational reasoning;
  \item how prediction changed under controlled removal of a sequence or edge.
\end{itemize}

These are still model diagnostics, not causal explanations, but they are more
directly linked to the proposed mechanisms than a heatmap alone.

\subsection{Case-Level Audit Template}
\label{sec:method-case-audit}

A fixed template is used for qualitative audit of representative correct and
incorrect cases:

\begin{itemize}
  \item ground-truth target and grade;
  \item predicted probability distribution and calibrated grade;
  \item localisation overlay;
  \item available MRI sequences;
  \item routing weights;
  \item dominant incoming graph relations/attention where extractable;
  \item uncertainty;
  \item nearest error category;
  \item whether the discrepancy is within one grade.
\end{itemize}

Cases are sampled from predefined categories rather than cherry-picked from the
best-looking outputs.

% ===========================================================================
\section{Reproducibility and Software Engineering}
\label{sec:method-reproducibility}

\subsection{Environment}
\label{sec:method-environment}

The implementation will use a reproducible Python/PyTorch medical-imaging
environment with DICOM tooling such as pydicom, SimpleITK and/or MONAI.
The final thesis records exact versions of Python, PyTorch, CUDA, cuDNN, GPU
hardware and all key libraries. Environment specification is stored in a
lockfile or container recipe.

\subsection{Experiment Configuration}
\label{sec:method-config}

Each run is configured from a machine-readable file containing:

\begin{itemize}
  \item dataset manifest and split version;
  \item selected sequence IDs;
  \item preprocessing and resampling;
  \item ROI dimensions;
  \item encoder architecture;
  \item SSL/routing/graph settings;
  \item loss coefficients;
  \item optimiser and scheduler;
  \item seed;
  \item checkpoint and early-stopping rule.
\end{itemize}

The configuration is copied into the run directory with metrics and model
weights so that a table in the thesis can be traced to an executable experiment.

\subsection{Data Provenance}
\label{sec:method-provenance}

Raw data are never edited in place. A manifest records the path/identifier and
processing status of every study. Derived crops can be regenerated from the
manifest and DICOM source, and each derived sample retains its pseudonymous
patient, level, condition, side, sequence and source-series identifier.

The local reference matrix is versioned independently of model output. If a
label is corrected after clinical review, the change is documented and the
affected experiments are rerun rather than silently updating test truth.

\subsection{Code and Model Release}
\label{sec:method-release}

Code for the public-data methodology will be released where institutional and
dataset licences permit. Pretrained weights are released if the training
dataset terms allow redistribution. Local Rizgary images and report-derived
individual-level data are not included unless a separate hospital approval
authorises public release.

The repository will include:

\begin{itemize}
  \item preprocessing and geometry reconstruction;
  \item split-generation/verification code;
  \item localisation/ROI pipeline;
  \item SSL, routing and graph modules;
  \item training and evaluation scripts;
  \item configuration files for primary experiments;
  \item instructions for reproducing public benchmark results.
\end{itemize}

% ===========================================================================
\section{Methodological Safeguards and Threats to Validity}
\label{sec:method-threats}

\subsection{Reference-Standard Shift}
\label{sec:method-threat-label}

LumbarDISC uses structured expert grades; Rizgary uses routine narrative
reports. A performance drop can therefore reflect both imaging-domain shift and
reference-standard shift. The thesis mitigates this by restricting transfer to
defensible overlap, auditing report labels and, where possible, re-reading a
stratified subset. It does not claim that the two reference standards are
identical.

\subsection{Single-Centre Local External Validation}
\label{sec:method-threat-single-centre}

Rizgary is a genuinely independent institution but still only one local centre.
External performance therefore demonstrates transport to \emph{that} setting,
not universal Middle Eastern generalisation. The thesis states this boundary
explicitly. A future multi-hospital regional study is required before making a
population-wide deployment claim.

\subsection{Model Complexity}
\label{sec:method-threat-complexity}

The full architecture contains several interacting components. Complexity can
produce apparent gains merely through parameter count and tuning effort. The
ablation ladder, matched backbone, homogeneous graph and shuffled-edge controls
are specifically designed to separate mechanism from capacity.

\subsection{Localisation Error}
\label{sec:method-threat-localisation}

The graph and routing modules cannot recover evidence absent from a bad crop.
Localisation is therefore measured and failures are reported separately.
Downstream results may additionally be stratified by localisation error to
estimate how much grading depends on the enabling stage.

\subsection{Class Imbalance}
\label{sec:method-threat-imbalance}

Severe grades are uncommon, so accuracy can be misleading and few-shot local
subsets may contain few Severe examples. The primary metrics are macro-F1,
balanced accuracy, weighted $\kappa$ and per-class recall, with patient-level
confidence intervals. Adaptation subsets are stratified where feasible and
their severity counts are always reported.

\subsection{Hyperparameter Multiplicity}
\label{sec:method-threat-hyperparameters}

A sufficiently large search can overfit the validation set even without touching
test data. The study limits the number of tunable components in each stage,
uses a fixed primary backbone for the core ablation, records the search space
and freezes choices before external testing.

\subsection{Clinical Utility}
\label{sec:method-threat-clinical}

This methodology evaluates technical grading, calibration and portability. It
does not demonstrate that radiologists become faster, safer or more accurate
when using the model. Such a claim requires a prospective reader-assistance or
impact study with independent clinical approval. The thesis therefore uses the
language of ``decision-support potential'' rather than clinical benefit.

% ===========================================================================
\section{Mapping Research Questions to Evidence}
\label{sec:method-rq-mapping}

\begin{longtable}{p{1.0cm}p{4.8cm}p{4.4cm}p{3.8cm}}
\toprule
\textbf{RQ} & \textbf{Primary comparison} & \textbf{Primary outcomes} & \textbf{Failure interpretation}\\
\midrule
\endfirsthead
\toprule
\textbf{RQ} & \textbf{Primary comparison} & \textbf{Primary outcomes} & \textbf{Failure interpretation}\\
\midrule
\endhead
RQ1 & independent/level transformer/homogeneous graph vs typed heterogeneous graph & macro-F1, $\kappa$, severe recall, edge ablations & graph topology carries little useful signal or local evidence dominates\\
RQ2 & generic pretraining vs anatomical cross-sequence SSL at 10/25/50/100\% labels & macro-F1, $\kappa$, label efficiency, external drop & anatomy-defined positives do not improve transferable representation\\
RQ3 & fixed fusion/cross-attention vs routing $\pm$ modality dropout & missing-sequence degradation, target-specific performance, calibration & fixed fusion is sufficient or learned routing is unstable\\
RQ4 & public held-out vs fixed Rizgary test; no adaptation vs PEFT/full tuning & zero-shot $\Delta$, recovery curve, calibration shift & domain shift is larger or less annotation-efficient than expected\\
RQ5 & CE vs ordinal/cost-aware objectives; calibrated vs uncalibrated prediction & Severe$\rightarrow$Normal rate, grade distance, Brier/ECE, risk--coverage & ordinal/cost formulation offers no advantage over CE\\
\bottomrule
\end{longtable}

This table defines what evidence can answer each question. It also makes the
thesis robust to partial negative results: RQ1 can be answered if the graph does
not help, and RQ2 can be answered if generic pretraining is equal or superior.
The contribution is the controlled experiment and the resulting knowledge, not
a requirement that every proposed mechanism win.

% ===========================================================================
\section{Planned Reporting Standard}
\label{sec:method-reporting}

The final manuscript/thesis reporting will follow the current applicable
medical-AI and diagnostic/prediction reporting checklists at the time of
submission. At minimum the report will include:

\begin{itemize}
  \item complete cohort flow diagrams and inclusion/exclusion reasons;
  \item reference-standard source, annotator qualifications and agreement where
        available;
  \item exact patient-level split logic and IDs reproducible within the approved
        environment;
  \item all primary metrics with confidence intervals;
  \item per-class Severe performance and ordinal error distances;
  \item calibration and external-domain performance;
  \item ablations for every claimed novel component;
  \item repeated-seed variability;
  \item compute/hardware cost;
  \item explicit limitations and negative results;
  \item a statement of what was not tested clinically.
\end{itemize}

The reporting checklist itself will be versioned with the manuscript because
CLAIM/STARD/TRIPOD+AI guidance may be updated during the candidature.

% ===========================================================================
\section{Chapter Summary}
\label{sec:method-summary}

This chapter translates the research gap of Chapter~\ref{ch:litreview} into a
controlled experimental design. The study begins from a structured definition
of lumbar disease: twenty-five level--condition--laterality targets on the
public benchmark rather than a generic image-level class. DICOM patient-space
geometry is preserved so that sagittal and axial evidence can be aligned
physically; localisation and disease-specific 2.5D ROIs are treated as shared
enabling technology; and all data partitions are defined at patient level.

Three methodological contributions are then tested independently. Anatomically
aligned cross-sequence self-supervision asks whether the same spinal level
viewed under different MRI sequences provides a stronger representation-learning
signal than generic augmentation. Disease-conditioned routing asks whether one
model can allocate different sequence evidence to different targets while
remaining functional when modalities are absent. The heterogeneous
disease--anatomy graph asks whether explicit typed relations among adjacent
levels, co-occurring compartments and bilateral targets improve grading beyond
independent heads, an inter-level transformer or a homogeneous graph. Ordinal,
cost-sensitive and uncertainty-aware outputs remain supporting methods and are
retained only if their contribution is measurable.

The translational design is deliberately separated from model development.
Rizgary is not mixed into public training. A fixed local test set establishes
zero-shot domain shift, while a separate adaptation pool supplies controlled
few-shot subsets. The primary local comparison is restricted to central-canal
stenosis, where the benchmark and narrative local reports can support a
defensible common target; local herniation morphology remains a separate task
rather than being forced into the RSNA schema.

Finally, the methodology defines success as an answer to the research
questions, not as attainment of a predetermined accuracy. A graph that adds
nothing, an ordinal loss that fails to beat cross-entropy, or a large
cross-institutional performance drop are all informative if demonstrated under
the same controlled protocol. This principle is what turns the architecture
from an engineering assembly into a doctoral experiment: every claimed
mechanism is falsifiable, every comparison is tied to a research question, and
the external cohort is used to test portability rather than to decorate an
internal result.

\printbibliography

\end{document}
 from the main thesis.
% ===========================================================================

\documentclass[12pt,a4paper]{report}

\usepackage{lmodern}
\usepackage[T1]{fontenc}
\usepackage[utf8]{inputenc}
\usepackage[english]{babel}
\usepackage{microtype}
\usepackage[margin=2.5cm]{geometry}
\usepackage{setspace}
\onehalfspacing

\usepackage{amsmath,amssymb}
\usepackage{booktabs}
\usepackage{longtable}
\usepackage{array}
\usepackage{graphicx}
\usepackage[hidelinks]{hyperref}

\usepackage[
  backend=biber,
  style=numeric-comp,
  sorting=none,
  maxbibnames=6,
  minbibnames=3,
  giveninits=true,
  doi=true,
  url=false,
  isbn=false
]{biblatex}

% Same bibliography used by Chapter 2.
\addbibresource{../lumbar_spine_mri_ai_literature_inventory.bib}

\AtEveryBibitem{%
  \clearfield{note}%
  \clearfield{comment}%
  \clearfield{annotation}%
  \clearfield{abstract}%
  \clearfield{file}%
}

\newcommand{\toconfirm}[1]{\textbf{[TO CONFIRM: #1]}}
\newcommand{\torecord}[1]{\textbf{[TO RECORD AFTER IMPLEMENTATION: #1]}}

\begin{document}

\chapter{Methodology}
\label{ch:methodology}

% ===========================================================================
\section{Introduction and Methodological Position}
\label{sec:method-intro}

Chapter~\ref{ch:litreview} established that automated lumbar MRI grading has
moved beyond the stage at which novelty can be claimed from replacing one image
backbone with another. Anatomical localisation, multi-sequence regions of
interest, inter-level context and ordinal grading have all been demonstrated in
recent work. The methodological question in this thesis is therefore not whether
a sufficiently large network can classify lumbar MRI. It is whether the
\emph{structure of the prediction problem} can be represented more faithfully:
which disease is being graded, at which level and side, which MRI sequence
contains the relevant evidence, how neighbouring targets are related, and how
the resulting model behaves when moved to a hospital whose scanners, protocols,
population and reporting conventions were not represented during development.

The study is designed as a staged, controlled experimental programme around
three methodological contributions and one translational validation programme.
The first contribution is an anatomically aligned cross-sequence
self-supervised representation in which positive relationships are defined by
DICOM patient-space correspondence rather than by generic image similarity.
The second is a disease-conditioned sequence router that learns how much each
available MRI sequence contributes to each anatomical target and is trained
explicitly to remain operational when one or more sequences are absent. The
third is a heterogeneous disease--anatomy graph whose nodes correspond to
level--condition--laterality targets and whose typed edges encode anatomical
adjacency, within-level relationships and bilateral symmetry. The translational
programme then separates development from deployment: the model is trained on a
public multinational benchmark, frozen, evaluated zero-shot on an independent
Rizgary Teaching Hospital cohort, and only afterwards adapted with controlled
amounts of local labelled data.

Several components are deliberately treated as \emph{enabling technology}
rather than as doctoral contributions. Anatomical localisation, segmentation,
2.5D region construction, conventional convolutional or transformer encoders,
class-imbalance handling, ordinal losses and calibration are all necessary to
make the comparison scientifically fair, but none is advertised as new by
itself. This distinction follows directly from the novelty boundary established
in Chapter~\ref{ch:litreview}: the thesis tests a new formulation of the
structured task rather than presenting a collection of familiar modules as if
their combination alone constituted originality.

The chapter is written as a protocol-grade methodology. Wherever an
implementation choice has not yet been fixed by evidence, it is specified as a
controlled comparison rather than invented retrospectively. Wherever an
institutional fact remains unresolved---for example the final number of
Rizgary studies surviving data reconciliation or the formal ethics reference---
it is marked explicitly. This is intentional. A doctoral methodology should
describe what is known, what is pre-specified and what will be decided by a
transparent rule; it should not convert planning assumptions into apparently
observed facts.

\subsection{Research Questions}
\label{sec:method-rqs}

The five research questions derived in Chapter~\ref{ch:litreview} are
operationalised as follows.

\begin{description}
  \item[RQ1 --- Relational disease modelling.] Does representing lumbar disease
  targets as a typed heterogeneous graph improve grading over independent target
  heads, an ordered five-level transformer and a homogeneous graph, and are any
  gains concentrated where single-target evidence is weakest?

  \item[RQ2 --- Anatomically aligned self-supervision.] Does cross-sequence
  pretraining based on DICOM-defined anatomical correspondence improve grading,
  label efficiency and transfer compared with ImageNet initialisation, generic
  medical pretraining and ordinary augmentation-based contrastive learning?

  \item[RQ3 --- Disease-conditioned modality routing.] Can one model learn
  target- and patient-dependent sequence weights while remaining robust to
  arbitrary missing-sequence configurations, without requiring a separately
  trained network for every combination of sagittal T1, sagittal T2/STIR and
  axial T2?

  \item[RQ4 --- Cross-institutional transfer.] What is the zero-shot degradation
  when a benchmark-trained model is applied to a previously unseen Middle
  Eastern clinical cohort, and how does performance recover as the amount of
  local annotation available for adaptation increases?

  \item[RQ5 --- Clinically asymmetric error and uncertainty.] Do ordinal and
  cost-sensitive objectives reduce clinically consequential distant errors,
  particularly Severe$\rightarrow$Normal/Mild, and can calibrated uncertainty
  support selective prediction without overstating clinical safety?
\end{description}

The order is not arbitrary. RQ2 and RQ3 determine how evidence is represented;
RQ1 determines how target predictions communicate; RQ5 determines how the
ordered output and its uncertainty are represented; and RQ4 asks whether those
choices survive the transition from a benchmark to real local clinical data.

\subsection{Design Principles}
\label{sec:method-principles}

Six principles govern all experiments.

\textbf{First, anatomy precedes grading.} Every classifier and ablation receives
the same anatomically localised input so that a comparison of representations is
not confounded by one model being given a better crop than another. This follows
the consistent evidence that localisation is a precondition of reliable
per-level grading \cite{lu2018deepspine,pan2021automatically,batra2025mscan}.

\textbf{Second, DICOM geometry is preserved.} The primary scientific
representation is reconstructed from DICOM patient-space coordinates. PNG
conversion may be used for visual inspection but is not the basis of
cross-sequence alignment, because it discards the spatial metadata required to
determine which axial slices correspond to a sagittal disc level.

\textbf{Third, patient identity defines independence.} No images, slices,
levels, sides or sequences from one patient may cross between training,
validation and test partitions. The unit of statistical independence for
uncertainty intervals and resampling is the patient, not the image or the disc
level.

\textbf{Fourth, the external cohort is not a tuning set.} The primary model is
developed without Rizgary images. Zero-shot performance is measured before any
local tuning. Few-shot adaptation then uses a separately defined local
adaptation pool while evaluation remains on a fixed local test set.

\textbf{Fifth, every claimed contribution requires an ablation.} A complete
model score does not establish that anatomical self-supervision, routing or
graph topology contributed anything. Each component is added and removed under
a controlled experiment, and a negative result is reported as such.

\textbf{Sixth, clinical interpretation is bounded by the reference standard.}
The model is trained to reproduce radiological grades whose inter-reader
reliability is imperfect. Performance is therefore interpreted against human
agreement and not as evidence that the network has discovered a more objective
biological truth.

% ===========================================================================
\section{Overall Research Design}
\label{sec:method-design}

\subsection{Study Type}
\label{sec:method-study-type}

The project is a retrospective, multi-dataset medical-imaging methods study with
three linked stages: public benchmark development, controlled methodological
ablation, and independent cross-institutional transfer. It does not prospectively
alter patient care and does not use model output to make treatment decisions.
The Rizgary component is an external retrospective evaluation of imaging and
radiology reports already acquired in routine clinical care.

The core development dataset is the RSNA Lumbar Degenerative Imaging Spine
Classification benchmark (LumbarDISC), which contains multi-sequence MRI and
expert per-level severity labels \cite{rsna2024rsna,richards2026the}. SPIDER is
used as an anatomical localisation/segmentation resource where its licence and
task definition permit \cite{vandergraaf2024lumbar}. The local Rizgary cohort is
reserved for external transfer and domain-adaptation analysis rather than being
mixed into the benchmark training data.

This separation is central to the inference the study is intended to support. A
random split of one pooled dataset can estimate internal generalisation under
that dataset's acquisition distribution. It cannot answer whether a model
trained on international benchmark data transports to a different hospital.
The latter requires the hospital to remain unseen during development.

\subsection{Unit of Analysis}
\label{sec:method-unit}

The primary unit of prediction on LumbarDISC is a
\emph{level--condition--laterality target}. At five disc levels
(L1--L2, L2--L3, L3--L4, L4--L5 and L5--S1), the benchmark defines five targets:
central canal stenosis, left and right neural foraminal narrowing, and left and
right subarticular stenosis. Full benchmark scope therefore contains
$5\times5=25$ targets per examination. Each target receives the ordered label
Normal/Mild, Moderate or Severe \cite{richards2026the}.

The unit of \emph{statistical independence}, however, remains the patient.
Twenty-five outputs from the same examination are correlated measurements, not
twenty-five independent observations. This distinction affects data splitting,
confidence intervals, significance tests and error analysis throughout the
chapter.

For the local Rizgary cohort, the target space is narrower. The reports contain
disc morphology and central-canal descriptions but do not reproduce the entire
RSNA schema reliably. The primary cross-institutional comparison is therefore
restricted to the five central-canal targets, L1--L2 through L5--S1, for which a
defensible overlap can be created. Bulge, protrusion and extrusion are analysed
as a separate local morphology task rather than being relabelled as if they were
RSNA stenosis grades. Foraminal and subarticular transfer are excluded unless
the local reference standard is expanded by new expert grading.

\subsection{Experimental Staging}
\label{sec:method-staging}

The methodology is intentionally cumulative:

\begin{enumerate}
  \item establish reproducible DICOM reconstruction, patient-level splits,
  localisation and an independent ROI baseline;
  \item establish anatomically correct cross-sequence correspondence;
  \item test disease-conditioned sequence routing and missing-modality training;
  \item test anatomically aligned cross-sequence self-supervision;
  \item test homogeneous and heterogeneous graph reasoning against independent
        target prediction;
  \item test ordinal/cost-sensitive output heads and probability calibration;
  \item freeze the selected public-data model and perform zero-shot external
        evaluation at Rizgary;
  \item adapt with prespecified numbers of local cases and estimate the
        annotation--performance recovery curve.
\end{enumerate}

The staging protects the thesis from a common failure mode of complex medical AI
systems: if the final model underperforms, a monolithic experiment cannot reveal
whether the failure arose from localisation, fusion, representation, graph
reasoning, loss design or domain shift. Here each stage creates a measurable
intermediate object and a direct baseline.

\subsection{Pre-Specified Hypotheses}
\label{sec:method-hypotheses}

The principal hypotheses are directional but not expressed as arbitrary target
accuracies.

\begin{description}
  \item[H1] Anatomical cross-sequence self-supervision will improve
  label-efficient grading and cross-institutional robustness relative to generic
  pretraining, particularly when only a fraction of benchmark labels is used.

  \item[H2] Disease-conditioned routing with modality dropout will degrade less
  under missing-sequence experiments than fixed concatenation or ordinary
  cross-attention trained only on complete studies.

  \item[H3] A typed heterogeneous graph will outperform independent target heads
  and a homogeneous graph if the anatomical edge types carry useful relational
  information; performance under a shuffled-edge control should deteriorate if
  the topology itself matters.

  \item[H4] Zero-shot performance on Rizgary will be lower than internal
  benchmark performance because of scanner, protocol, population and reference-
  standard shift; the magnitude and pattern of that drop are treated as findings
  rather than failures of the project.

  \item[H5] Local adaptation will improve performance as annotation increases,
  but no predetermined percentage of ``recovery'' is required. The shape of the
  recovery curve is the result.

  \item[H6] Cost-sensitive/ordinal objectives may reduce distant
  Severe$\rightarrow$Normal/Mild errors even if they do not improve aggregate
  macro-F1; plain cross-entropy remains the mandatory baseline because previous
  lumbar work has shown that ordinal objectives do not automatically help
  \cite{niemeyer2021a}.
\end{description}

% ===========================================================================
\section{Data Sources}
\label{sec:method-data}

\subsection{Public Development Cohort: LumbarDISC}
\label{sec:method-lumbardisc}

The primary development data are drawn from the publicly released RSNA 2024
Lumbar Spine Degenerative Classification dataset. The published LumbarDISC
resource comprises 2,697 patients and 8,593 MRI series collected across eight
institutions in six countries \cite{richards2026the}. The locally held labelled
competition training subset contains approximately 1,975 studies and is used for
core development. The distinction between the published cohort size and the
available labelled subset is retained explicitly throughout the thesis rather
than treating the two numbers as interchangeable.

Eligible benchmark examinations contain sagittal T2-like imaging, sagittal T1
and axial T2, with expert severity grading at the five lumbar levels
\cite{rsna2024rsna,richards2026the}. Coordinate/localiser annotations are also
available for relevant disease locations. These properties make the dataset
appropriate for all three methodological contributions: it is multi-sequence,
level-resolved, multi-target, geographically heterogeneous and sufficiently
large to separate representation learning from local external validation.

The benchmark serves three roles:

\begin{enumerate}
  \item supervised development of per-target grading;
  \item construction of anatomically corresponding cross-sequence training pairs;
  \item internal/site-aware robustness analysis before the model is exposed to
        Rizgary.
\end{enumerate}

All dataset versions, downloaded files, checksums where feasible and the final
patient split lists will be archived in the reproducibility record. The exact
number of studies remaining after local quality-control exclusions will be
\torecord{final LumbarDISC development, validation and held-out counts}.

\subsection{Public Anatomical Resource: SPIDER}
\label{sec:method-spider}

SPIDER provides manually delineated lumbar anatomy on sagittal MRI from multiple
hospitals, including vertebrae, intervertebral discs and spinal canal
\cite{vandergraaf2024lumbar}. It is used only for anatomical localisation,
segmentation benchmarking or pretraining where permitted by its licence. It is
not used to manufacture disease labels that are absent from the dataset.

The practical purpose is to avoid making localisation the doctoral novelty.
Rather than spend the main experimental budget designing another segmentation
network, the thesis will use or reproduce a strong established anatomical model
and quantify its localisation error. If a SPIDER-pretrained model is used to
initialise the localiser, the external disease-grading experiments remain
separate because the disease targets originate in LumbarDISC.

\subsection{Independent Local Cohort: Rizgary Teaching Hospital}
\label{sec:method-rizgary}

The local cohort consists of lumbar MRI examinations and corresponding English
radiology reports supplied by Rizgary Teaching Hospital. Current project
inventory identifies 299 narrative reports and 294 candidate cases with matching
multi-sequence imaging, while a wider raw DICOM audit contains additional
studies. These counts are not treated as final until a one-to-one cohort
reconciliation is complete.

The raw imaging inspected during protocol development confirms the sequence
families required by the thesis---sagittal T1, sagittal T2 and axial T2---and
preserves standard DICOM spatial metadata. A sample examination was acquired on
a Siemens Avanto 1.5-T system; scanner/vendor distribution across the full local
cohort will be reported after complete metadata extraction rather than inferred
from one case.

The local cohort is valuable precisely because it is not curated to reproduce
the public benchmark. Its scanner, acquisition protocol, prevalence, reporting
language and label completeness differ from LumbarDISC. Those differences are
the domain shift the transfer experiment is intended to measure, not nuisances
to be eliminated by silently harmonising labels or discarding difficult cases.

\subsection{Local Radiology Reports and Historical Spreadsheet}
\label{sec:method-local-reports}

The narrative radiology reports are the primary textual source for the local
reference standard. The existing spreadsheet is treated as a secondary
historical transcription aid rather than as unquestioned ground truth. It
contains patient-level variables and comma-separated level references, but it
does not constitute a complete level-resolved RSNA-style matrix and has known
transcription inconsistencies. The methodology therefore reconstructs the local
reference matrix from source reports and uses the spreadsheet only for
cross-checking.

The canonical local matrix will preserve, at minimum:

\begin{itemize}
  \item pseudonymous case identifier;
  \item lumbar level;
  \item central-canal stenosis severity where explicitly stated;
  \item disc bulge, protrusion and extrusion as separate morphology fields;
  \item foraminal findings and side only where explicitly stated;
  \item nerve-root pressure/impingement where reported;
  \item an explicit missing/not-reported state distinct from a negative finding.
\end{itemize}

This last distinction is essential. Absence of a phrase in a narrative report is
not automatically equivalent to ``normal''. A field may be unreported because
the radiologist considered it clinically irrelevant, because reporting practice
varied, or because the report was not structured around that compartment.
``Not reported'' is therefore encoded separately and excluded from target-
specific external evaluation unless a clinician has explicitly adjudicated the
absence as normal.

\subsection{Cohort Reconciliation}
\label{sec:method-reconciliation}

Before any local model evaluation, a case-flow table will reconcile the raw
archive, reports and analysable imaging. For each candidate case the audit will
record:

\begin{itemize}
  \item whether a report exists;
  \item whether sagittal T1 exists and is usable;
  \item whether sagittal T2/STIR exists and is usable;
  \item whether axial T2 exists and is usable;
  \item whether DICOM geometry is internally consistent;
  \item whether a central-canal reference label can be derived;
  \item whether laterality-dependent labels are available;
  \item whether the case is duplicated or represents repeat imaging;
  \item the reason for any exclusion.
\end{itemize}

The final thesis will report the flow from raw studies to the fixed external
test cohort, rather than presenting the current 299/294/raw-study counts without
explanation. Repeat examinations from one patient, if identified, remain in the
same partition to preserve patient independence.

% ===========================================================================
\section{Ethics, Governance and De-Identification}
\label{sec:method-ethics}

The Rizgary component uses retrospective clinical data. No MRI is acquired for
the purposes of this study, no patient undergoes an additional intervention, and
model output is not returned to the clinical record. Nevertheless, the data are
sensitive medical information and are subject to institutional governance.

The final thesis will state the formal approval mechanism and reference:
\toconfirm{hospital research committee/IRB approval number, date and responsible
institution}. It will also state the legal/ethical basis for retrospective use:
\toconfirm{consent waiver, broad consent or other applicable basis}. These items
must be completed before the local study is described as ethically approved.

The raw local DICOM archive contains direct identifiers and therefore cannot be
given to students or uploaded to external compute in its original state.
De-identification is performed on a controlled copy before analysis. At minimum,
the process removes direct patient identifiers, replaces case identifiers with
study pseudonyms, strips or audits private tags, and handles dates according to
the approved policy while preserving only the temporal information required for
scientific analysis. The de-identification procedure must also check for
burned-in identifiers where applicable. The re-identification key, if one must
exist for hospital governance, is stored separately from the research
repository and is not available to the model-training pipeline.

External/cloud processing of local images is prohibited unless explicitly
allowed by hospital policy:
\toconfirm{whether external cloud/GPU processing is permitted and under what
conditions}. Public benchmark data are handled according to their published
licence and competition terms.

Only aggregate results are reported. Illustrative local images may be included
in the thesis only after de-identification review and only if institutional
permission permits publication. The local cohort is not assumed to be publicly
shareable; any future data release would require separate approval and a more
stringent disclosure review.


% ===========================================================================
\section{Reference Standard and Target Encoding}
\label{sec:method-reference}

\subsection{Benchmark Target Representation}
\label{sec:method-benchmark-targets}

For LumbarDISC, the target set is represented explicitly rather than collapsed
into a generic three-class image label. Let the five lumbar levels be
\[
\mathcal{L}=\{L1\!-\!L2,L2\!-\!L3,L3\!-\!L4,L4\!-\!L5,L5\!-\!S1\},
\]
and let the five condition/laterality targets be
\[
\mathcal{C}=\{
\mathrm{Canal},
\mathrm{LF},
\mathrm{RF},
\mathrm{LSA},
\mathrm{RSA}
\},
\]
where LF/RF denote left/right neural foraminal narrowing and LSA/RSA
left/right subarticular stenosis. One examination therefore has the target set
\[
\mathcal{T}=\mathcal{L}\times\mathcal{C}, \qquad |\mathcal{T}|=25.
\]

For target $t\in\mathcal{T}$ the observed grade is
\[
y_t\in\{0,1,2\},
\]
corresponding to Normal/Mild, Moderate and Severe. Encoding the targets in this
form is not merely convenient notation. The pair $(l,c)$ is retained throughout
data loading, feature extraction, graph construction and evaluation, which
prevents a severity label from losing the anatomical meaning that generated it.

Where the benchmark provides an absent or unusable annotation, a target mask
$m_t\in\{0,1\}$ is carried alongside the grade. Masked targets do not contribute
to the supervised loss. Missing labels are never imputed from neighbouring
levels for the main experiments, because doing so would introduce exactly the
relational assumption the thesis is intended to test.

\subsection{Local Reference Matrix}
\label{sec:method-local-reference}

The Rizgary reference standard is constructed at the same anatomical level but
not forced into the same disease schema. A canonical row represents one
case--level pair. The minimum schema is:

\begin{center}
\begin{tabular}{p{2.4cm}p{2.0cm}p{8.1cm}}
\toprule
\textbf{Field} & \textbf{Type} & \textbf{Interpretation}\\
\midrule
Case ID & categorical & pseudonymous patient/study identifier\\
Level & categorical & L1--L2 through L5--S1\\
Canal stenosis & ordinal/missing & explicit report-derived severity, if stated\\
Bulge & binary/missing & circumferential/generalised bulge stated at level\\
Protrusion & binary/missing & protrusion stated at level\\
Extrusion & binary/missing & extrusion stated at level\\
Foraminal finding & categorical/missing & narrowing/pressure effect if stated\\
Laterality & left/right/bilateral/NR & side only where stated\\
Nerve-root effect & categorical/missing & pressure, displacement or compression\\
Other morphology & text/coded & migration, dehydration, height loss and related findings\\
\bottomrule
\end{tabular}
\end{center}

The local matrix keeps \emph{not reported} (NR) separate from \emph{explicitly
absent}. This rule is especially important for foraminal and subarticular
findings, because local reporting completeness differs from the structured RSNA
annotation scheme. A target is eligible for external evaluation only when the
reference standard provides a defensible class label.

\subsection{Report Extraction and Human Verification}
\label{sec:method-report-verification}

Automated report extraction, if available from the parallel clinical-NLP work,
may accelerate construction of the matrix but does not define the gold
standard. The final labels used for Selar's external evaluation are manually
verified against the original report text.

A two-stage verification procedure is preferred:

\begin{enumerate}
  \item \textbf{Structured extraction audit.} Every non-normal/non-empty finding
  is traced back to the exact report phrase, level and laterality. Multi-level
  expressions such as ``L2--3 and L4--5, L5--S1 shows circumferential disc
  bulge'' are expanded into separate level records without changing their
  meaning.

  \item \textbf{Clinical reliability subset.} A pre-specified stratified subset
  of local studies is independently re-read by two qualified readers using an
  explicit grading sheet. Disagreements are adjudicated by a more senior reader
  where available. Inter-reader agreement is reported with weighted
  $\kappa$ for ordinal grades and ordinary agreement/appropriate $\kappa$ for
  binary morphology. The size of the re-read subset will be fixed before model
  results are examined, based on reader availability and the precision required
  for the agreement estimate.
\end{enumerate}

If a second-reader study cannot be completed, the report-derived matrix remains
usable as a real-world reference standard, but the absence of independent
regrading is stated explicitly as a limitation and no claim is made that local
labels reproduce the RSNA grading protocol exactly.

\subsection{Schema Alignment Between LumbarDISC and Rizgary}
\label{sec:method-schema-alignment}

Cross-institutional evaluation requires equality of \emph{task}, not merely
similar words. The mapping is therefore conservative.

\begin{center}
\begin{tabular}{p{4.1cm}p{4.0cm}p{5.0cm}}
\toprule
\textbf{LumbarDISC target} & \textbf{Rizgary evidence} & \textbf{Transfer status}\\
\midrule
Central canal stenosis & level-resolved mild/moderate/severe statements where present
& Primary external target\\
Left/right foraminal narrowing & pressure/narrowing often incompletely lateralised
& Exclude unless laterality and grade are adequate or freshly re-read\\
Left/right subarticular stenosis & not routinely represented in current reports
& Exclude unless new expert annotation is created\\
Disc bulge/protrusion/extrusion & directly described locally but not an RSNA stenosis target
& Separate local morphology task\\
\bottomrule
\end{tabular}
\end{center}

No local disc-herniation morphology is converted into a stenosis grade by
heuristic rule. Likewise, ``pressure on exit foramina'' is not automatically
treated as a left/right foraminal severity label unless the report or expert
review supplies the information required by that target. This prevents
apparent external validation from being created by label re-interpretation.

% ===========================================================================
\section{Data Partitioning and Leakage Control}
\label{sec:method-splits}

\subsection{Patient-Level Independence}
\label{sec:method-patient-split}

All splits are made at patient/study level before slice extraction or data
augmentation. For a patient $p$, every sagittal series, every axial series,
every disc level, every side and every augmented crop belongs to one and only
one partition. A programmatic assertion is run before training to confirm that
the intersection of patient identifiers across partitions is empty.

The split record is version-controlled as a list of pseudonymous IDs. Data
loaders consume those fixed lists rather than performing a new random split each
time the training script is executed.

\subsection{Public Development, Validation and Held-Out Evaluation}
\label{sec:method-public-split}

Where the available public dataset exposes institution identifiers, a
site-aware evaluation is preferred because it tests a more meaningful shift
than a random split. The exact public partitioning scheme will be chosen
according to the metadata available in the held dataset and will be frozen
before model comparison. At minimum there will be:

\begin{itemize}
  \item a development/training partition;
  \item a validation partition for hyperparameter choice, early stopping and
        calibration;
  \item a held-out public evaluation partition not used for model selection.
\end{itemize}

If site identifiers permit leave-one-site-out or institution-held-out analysis,
that analysis is reported in addition to the ordinary benchmark split. It is not
a replacement for Rizgary, because the local hospital introduces a different
reporting/reference-standard domain as well as a different imaging domain.

\subsection{Fixed Local Test Set and Adaptation Pool}
\label{sec:method-local-split}

The local transfer study is protected against adaptation leakage by separating
Rizgary into two parts \emph{before} the public model is evaluated:

\begin{enumerate}
  \item a \textbf{fixed local test set}, which is never used for adaptation,
        calibration or threshold selection; and
  \item a \textbf{local adaptation pool}, from which the $N=10,25,50,100$
        few-shot subsets are sampled.
\end{enumerate}

The exact number allocated to each part will be determined after the audited
class distribution is known. The split will be stratified by central-canal
severity and lumbar level where feasible without producing clinically
implausible or extremely small strata. The fixed local test set is used for the
zero-shot baseline and for every subsequent adapted model, making the recovery
curves directly comparable.

For each adaptation size $N$, several independent stratified subsets are sampled
from the adaptation pool. Results are averaged across these draws with variance
reported, because one set of ten cases can be unrepresentative by chance. The
test set itself never changes across adaptation sizes.

\subsection{Prevention of Information Leakage During Self-Supervision}
\label{sec:method-ssl-leakage}

Self-supervised learning does not remove the need for split discipline. A model
pretrained on the held-out test patients, even without labels, has already seen
their anatomy and acquisition characteristics. Therefore anatomical
self-supervision uses only the development partition during internal
experiments. The held-out public test set and fixed Rizgary test set remain
unseen in both supervised and self-supervised training.

The same restriction applies to intensity normalisation parameters and learned
harmonisation models: any transformation whose parameters are estimated from
data is fitted on the training/development partition and then applied unchanged
to validation/test data.

% ===========================================================================
\section{DICOM Reconstruction and Quality Control}
\label{sec:method-dicom}

\subsection{Series Identification}
\label{sec:method-series-identification}

Each study is first organised by DICOM \texttt{SeriesInstanceUID}. Candidate
series are classified into sagittal T1, sagittal T2/STIR and axial T2 using a
combination of DICOM metadata and image-orientation checks. Series descriptions
are informative but not trusted alone because naming conventions differ across
institutions. The classifier therefore considers fields such as
\texttt{SeriesDescription}, \texttt{ProtocolName}, sequence parameters where
available, and the orientation derived from \texttt{ImageOrientationPatient}.

Ambiguous or duplicated candidate series are flagged for manual review rather
than selected by filename order. The selected series identifiers are stored in
a manifest so that every experiment uses the same source series.

\subsection{Patient-Space Geometry}
\label{sec:method-geometry}

For each DICOM slice, the image plane is reconstructed from
\texttt{ImagePositionPatient}, \texttt{ImageOrientationPatient} and
\texttt{PixelSpacing}. Let $\mathbf{s}$ be the patient-space location of the
upper-left pixel, $\mathbf{r}$ and $\mathbf{c}$ the row and column direction
cosines, and $\Delta_r,\Delta_c$ the corresponding pixel spacings. The physical
coordinate of pixel $(i,j)$ is

\begin{equation}
\mathbf{x}(i,j)=
\mathbf{s}
+i\,\Delta_r\,\mathbf{r}
+j\,\Delta_c\,\mathbf{c}.
\label{eq:dicom-patient-coordinate}
\end{equation}

The slice normal is

\begin{equation}
\mathbf{n}=\mathbf{r}\times\mathbf{c}.
\label{eq:slice-normal}
\end{equation}

Slices are ordered by projection onto $\mathbf{n}$ rather than by filename.
Spacing between consecutive slice origins is inspected for discontinuities.
This procedure preserves patient coordinates across sagittal and axial
acquisitions, allowing a level localised in one plane to be projected into the
other.

For a sagittal target point $\mathbf{q}$ and candidate axial slice $k$ with
origin $\mathbf{s}_k$ and normal $\mathbf{n}_k$, the perpendicular distance to
that slice plane is

\begin{equation}
d_k =
\left|
\mathbf{n}_k^\top(\mathbf{q}-\mathbf{s}_k)
\right|.
\label{eq:plane-distance}
\end{equation}

The nearest anatomically corresponding axial slices are the slices with the
smallest $d_k$, subject to coverage and orientation checks. The method is an
implementation of physical correspondence, not a learned registration shortcut,
and follows the geometric principle used by prior lumbar pipelines
\cite{pan2021automatically,batra2025mscan}.

\subsection{Orientation Standardisation}
\label{sec:method-orientation}

Volumes are reoriented to a common anatomical convention for tensor processing
while retaining the transform to original patient coordinates. Left/right must
never be inferred after arbitrary array transposition; all laterality-dependent
targets are checked against the DICOM orientation transform. This is particularly
important for foraminal and subarticular labels, where a left/right inversion
would create a scientifically silent but catastrophic error.

\subsection{Voxel Spacing and Resampling}
\label{sec:method-resampling}

The acquisition spacing varies by institution, scanner and plane. Inputs are
therefore resampled only after anatomical coordinates have been reconstructed.
The target spacing is chosen from the training-set distribution so that it
preserves structures relevant to grading without creating unnecessary memory
cost. Sagittal and axial streams may use different in-plane resolutions because
their native geometry and diagnostic roles differ.

Interpolation uses a continuous method (e.g.\ linear or higher-order interpolation)
for MR intensities and nearest-neighbour interpolation for discrete masks or
labels. Original spacing is retained in metadata for any measurement reported
in millimetres. The exact target spacing will be \torecord{final sagittal and
axial resampling spacings selected from the development cohort}.

\subsection{Intensity Processing}
\label{sec:method-intensity}

MRI intensities are not quantitatively standardised across scanners. The
preprocessing therefore removes gross scale differences without imposing a
global histogram learned from test data. A default pipeline is:

\begin{enumerate}
  \item suppress non-finite values and obvious padding/background values;
  \item define a foreground mask from the body/spine region;
  \item clip extreme foreground intensities using training-derived robust
        percentiles;
  \item perform per-volume or per-ROI robust z-score normalisation;
  \item optionally compare with bias-field correction or histogram
        standardisation as a robustness ablation.
\end{enumerate}

Any learned harmonisation is fitted only on the development data. The external
Rizgary zero-shot result is reported both with the baseline preprocessing and,
if a separate harmonisation experiment is conducted, with that method clearly
labelled as an intervention rather than folded into the zero-shot definition.

\subsection{Quality-Control Flags}
\label{sec:method-qc}

Automated QC records the following per series:

\begin{itemize}
  \item slice count and ordering consistency;
  \item missing or duplicated instance positions;
  \item orientation consistency within series;
  \item in-plane spacing and slice spacing;
  \item field of view and lumbar coverage;
  \item intensity dynamic range;
  \item severe motion or truncation where detectable;
  \item presence of all sequences required by the particular experiment.
\end{itemize}

Cases failing geometry checks are not silently forced into the pipeline. They
are reviewed, excluded with a recorded reason, or assigned to a robustness
analysis if the failure itself represents a realistic missing-data condition.
A final inclusion/exclusion table will report every such decision.

% ===========================================================================
\section{Anatomical Localisation and ROI Construction}
\label{sec:method-localisation}

\subsection{Role of Localisation}
\label{sec:method-localisation-role}

Localisation exists to make the subsequent comparison fair. The thesis does not
claim that identifying L1--L2 through L5--S1 is itself novel; Chapter~2 showed
that this problem is mature and that strong systems already attain highly
reliable level assignment \cite{pan2021automatically}. The requirement here is
that the localiser be accurate enough that representation and graph experiments
are not dominated by wrong crops.

An established model is therefore preferred. Candidate implementations include
a SPIDER-pretrained nnU-Net-style segmentation network
\cite{isensee2021nnunet,vandergraaf2024lumbar}, a heatmap-regression detector,
or another state-of-the-art anatomical parser available when implementation
begins. The choice is made on validation localisation error and robustness, not
on whether it is architecturally fashionable.

\subsection{Outputs}
\label{sec:method-localisation-output}

The localiser produces, for each study:

\begin{itemize}
  \item centres of L1--L2 through L5--S1 in patient coordinates;
  \item vertebral/disc/canal masks or landmarks where the selected model provides
        them;
  \item a confidence or quality flag for each detected level;
  \item the transform between voxel coordinates and DICOM patient space.
\end{itemize}

A failed or low-confidence localisation is retained as a distinct failure mode.
It is not corrected manually for the primary automated pipeline without being
reported, because doing so would inflate downstream grading performance by
removing a real deployment error.

\subsection{Localisation Evaluation}
\label{sec:method-localisation-eval}

Where point annotations are available, centre localisation error is measured in
millimetres:

\begin{equation}
e_l = \|\hat{\mathbf{q}}_l-\mathbf{q}_l\|_2.
\label{eq:localisation-error}
\end{equation}

Mean/median error and percentile distributions are reported by level.
Percentage-of-correct-keypoint (PCK) or an equivalent tolerance-based measure is
also reported at pre-specified physical thresholds. Where full masks are
available, Dice and surface-distance metrics are reported. These results are
separated from disease grading, consistent with the literature showing that
detection/localisation and fine-grained grading generalise differently.

\subsection{Disease-Specific 2.5D Regions of Interest}
\label{sec:method-roi}

One rectangular crop is not assumed to be equally appropriate for every target.
The crop definition is conditioned on anatomical compartment.

For central canal stenosis, the model receives sagittal T2/STIR context around
the disc level together with corresponding axial T2 slices. For neural
foraminal narrowing, the sagittal T1 stream receives greater parasagittal
coverage and the ROI is side-aware, with sagittal T2 retained as complementary
context. For subarticular stenosis, axial T2 supplies the dominant local
evidence with sagittal context. These choices follow the clinical sequence--
pathology correspondence established in Chapter~2 rather than being learned
from arbitrary crop templates.

Each 2.5D stream contains a small ordered stack around the anatomically closest
slice,

\[
\mathcal{I}_l =
\{I_{k-r},\ldots,I_k,\ldots,I_{k+r}\},
\]

where $k$ is the level-centred slice and $r$ is a small radius determined on the
development set. A five-slice stack ($r=2$) is the initial reference
configuration because it captures through-plane context at modest memory cost;
the final radius is reported after sensitivity testing rather than treated as
a novel contribution.

Crops are defined in physical dimensions where possible, then resampled to the
network input resolution. This avoids a fixed 100-pixel crop representing
different anatomical widths on scanners with different pixel spacing.

\subsection{ROI Quality Control}
\label{sec:method-roi-qc}

A validation subset is visually inspected using overlays of the detected level,
crop boundary and corresponding axial slices. The inspection records:

\begin{itemize}
  \item correct lumbar level;
  \item inclusion of the relevant canal/foramen/recess;
  \item correct left/right orientation;
  \item adequate axial coverage;
  \item crop truncation;
  \item correspondence failure caused by unusual axial angulation.
\end{itemize}

The purpose is not to manually curate the full dataset but to detect systematic
geometry errors before they propagate into model training.

% ===========================================================================
\section{Baseline Representation and Sequence Encoders}
\label{sec:method-baseline}

\subsection{Independent ROI Classifier}
\label{sec:method-e0}

The mandatory baseline (E0) is an independent ROI classifier. Each target is
graded from its anatomically localised input without communication with other
targets. This represents the prevailing formulation criticised in
Chapter~\ref{ch:litreview} and provides the reference against which relational
reasoning is measured.

The baseline uses a well-characterised image encoder such as ResNet
\cite{he2016deep}, EfficientNet \cite{tan2019efficientnet}, Swin
\cite{liu2021swin}, or another contemporary backbone selected before final
experiments. The purpose is not to prove that one family is universally
superior. A primary backbone is selected for the ablation ladder so that
representation, routing and graph changes are measured under one stable
feature extractor. A second backbone may be used in a model-agnosticity check
for the strongest contribution.

\subsection{Sequence-Specific Feature Extraction}
\label{sec:method-sequence-encoders}

Each available MRI sequence is processed by a sequence-specific encoder
$E_m(\cdot)$,

\begin{equation}
\mathbf{f}_{p,l,m}=E_m(\mathcal{I}_{p,l,m}),
\label{eq:sequence-encoder}
\end{equation}

for patient $p$, level $l$ and modality/sequence $m$. The encoders may share
architecture while keeping separate normalisation statistics and weights, or
may partially share weights if an explicit ablation shows that doing so is
beneficial. Separate encoders are the default because T1, sagittal T2/STIR and
axial T2 have different contrast and orientation statistics.

To keep comparisons fair, the dimensionality of the output representation is
held constant across fusion methods. Fixed concatenation, simple averaging and
ordinary cross-attention all consume the same per-sequence features.

\subsection{Controlled Backbone Selection}
\label{sec:method-backbone-control}

Backbone choice is selected on the validation partition before the central
novelty experiments and then held fixed. The candidate comparison will use
identical crops, augmentations, optimiser budget and patient split. If one
backbone is selected because it gives the most stable validation performance
rather than the absolute highest point estimate, that decision is recorded.

This design prevents the thesis from becoming a repeated architecture search.
The main questions concern anatomical self-supervision, adaptive modality
allocation and target-level relational reasoning. A backbone is a controlled
feature extractor in those experiments.


% ===========================================================================
\section{Core Contribution I: Anatomically Aligned Cross-Sequence Self-Supervision}
\label{sec:method-ssl}

\subsection{Rationale}
\label{sec:method-ssl-rationale}

Conventional self-supervised vision learning defines related views by applying
augmentations to one image. That is useful but does not exploit a distinctive
property of multi-sequence MRI: sagittal T1, sagittal T2/STIR and axial T2 can
look very different while representing the same underlying patient and spinal
level. DICOM geometry supplies the correspondence directly.

The proposed objective therefore treats anatomical identity, rather than visual
similarity, as the primary source of positive pairs. For patient $p$, level
$l$, and two sequences $m_a$ and $m_b$, embeddings
$\mathbf{z}_{p,l,m_a}$ and $\mathbf{z}_{p,l,m_b}$ are encouraged to encode a
shared anatomical representation even though acquisition plane and contrast may
differ.

This idea is evaluated as a representation-learning hypothesis, not assumed to
be beneficial. The literature already shows that generic and longitudinal
self-supervision can reduce label dependence
\cite{jamaludin2017selfsupervised,chen2020simple}; the experiment here asks
whether explicitly anatomical cross-sequence correspondence is a stronger
signal for lumbar grading.

\subsection{Projection Head}
\label{sec:method-ssl-projection}

Each sequence feature is mapped through a projection head $P_m$:

\begin{equation}
\mathbf{z}_{p,l,m}=
\frac{
P_m(\mathbf{f}_{p,l,m})
}{
\|P_m(\mathbf{f}_{p,l,m})\|_2
}.
\label{eq:ssl-projection}
\end{equation}

The projection space is used for the self-supervised objective; the encoder
features before projection are retained for downstream grading. This prevents
the representation used for contrastive geometry from constraining the final
classifier unnecessarily.

\subsection{Positive-Pair Construction}
\label{sec:method-positive-pairs}

The strongest positive relation is:

\[
(p,l,m_a)\sim(p,l,m_b),
\qquad m_a\neq m_b,
\]

meaning the same patient and same anatomical level in two different sequences.
Correspondence is established through the DICOM reconstruction and level
localiser, not through nearest image appearance.

An optional softer relation may be tested for adjacent levels from the same
patient, because neighbouring levels share biomechanical context. It is not
included in the primary objective by default. Treating L3--L4 and L4--L5 as
equivalent positives would risk erasing genuine level-specific pathology, so
any adjacent-level term is isolated as a separate ablation.

Different patients at the same level are not defined as positives in the main
objective, because common anatomical level does not imply common disease state.
They may be used in a secondary supervised-contrastive variant once labels are
available, but this is conceptually distinct from the proposed anatomical
self-supervision.

\subsection{Contrastive Objective}
\label{sec:method-ssl-objective}

For anchor embedding $\mathbf{z}_i$ and its cross-sequence positive
$\mathbf{z}_{i^+}$, the reference InfoNCE-style loss is

\begin{equation}
\mathcal{L}_{\mathrm{ACSSL}}(i)
=
-\log
\frac{
\exp(\mathrm{sim}(\mathbf{z}_i,\mathbf{z}_{i^+})/\tau)
}{
\sum_{j\in\mathcal{B}\setminus\{i\}}
\exp(\mathrm{sim}(\mathbf{z}_i,\mathbf{z}_j)/\tau)
},
\label{eq:acssl}
\end{equation}

where $\mathrm{sim}$ is cosine similarity, $\tau$ is a temperature parameter
and $\mathcal{B}$ is the contrastive batch. When more than one anatomically
valid positive exists, the numerator is extended to the set of positives rather
than choosing one arbitrarily.

False negatives are a concern because two different patients may have very
similar normal anatomy. The implementation will therefore compare ordinary
InfoNCE with a variant that does not force all non-positive samples to be hard
negatives, such as a multi-positive or non-contrastive consistency objective.
The novelty claim is tied to the \emph{anatomical pairing}, not to a particular
brand of contrastive loss.

\subsection{Label-Efficiency Experiment}
\label{sec:method-label-efficiency}

The primary test of RQ2 uses fixed fractions of the labelled development set:
10\%, 25\%, 50\% and 100\%. For each fraction:

\begin{enumerate}
  \item the encoder is initialised from the candidate pretraining strategy;
  \item the same supervised grading head is fitted;
  \item patient splits and target distribution are held fixed;
  \item several random labelled subsets are evaluated where computationally
        feasible;
  \item macro-F1, weighted $\kappa$, severe-class recall and calibration are
        compared.
\end{enumerate}

The competing initialisations are:

\begin{itemize}
  \item training from random initialisation;
  \item ImageNet initialisation;
  \item generic medical/volumetric pretraining where a suitable public model is
        available;
  \item ordinary augmentation-based self-supervision;
  \item the proposed anatomically aligned cross-sequence self-supervision.
\end{itemize}

A successful result is not defined as beating every baseline at every label
fraction. The useful scientific question is where anatomical pairing helps,
where it does not, and whether any benefit is largest under label scarcity or
domain shift.

\subsection{Transfer Test of the Representation}
\label{sec:method-ssl-transfer}

The public-data model used for zero-shot Rizgary evaluation is trained once with
and once without anatomical self-supervision while the downstream architecture
is held constant. If the anatomically pretrained model loses less performance
externally despite comparable internal performance, that would support the
claim that the representation is less tied to dataset-specific appearance. If
both models degrade equally, the self-supervised contribution is limited to
label efficiency rather than domain robustness.

% ===========================================================================
\section{Core Contribution II: Disease-Conditioned Adaptive Sequence Routing}
\label{sec:method-routing}

\subsection{Problem Formulation}
\label{sec:method-routing-problem}

A conventional multi-sequence model applies the same fusion rule to every
target. That assumption is clinically implausible because the sequences are not
equally informative for each anatomical compartment. Sagittal T1 is particularly
useful for foraminal fat obliteration, sagittal T2/STIR for disc/canal context,
and axial T2 for canal and lateral recess morphology. The router turns that
known heterogeneity into a learnable, testable mechanism.

For patient $p$, target $t=(l,c)$ and available modalities
$\mathcal{M}_p$, let $\mathbf{f}_{p,l,m}$ be the feature from modality $m$.
The fused target feature is

\begin{equation}
\mathbf{u}_{p,t}
=
\sum_{m\in\mathcal{M}_p}
g_{p,t,m}\,
\mathbf{f}_{p,l,m},
\label{eq:routed-feature}
\end{equation}

where the non-negative routing weights satisfy

\begin{equation}
\sum_{m\in\mathcal{M}_p} g_{p,t,m}=1.
\label{eq:routing-simplex}
\end{equation}

The gate is conditioned on target identity and patient features:

\begin{equation}
g_{p,t,m}
=
\mathrm{softmax}_{m\in\mathcal{M}_p}
\left(
G[
\mathbf{f}_{p,l,m},
\mathbf{e}_{c},
\mathbf{e}_{l},
\mathbf{a}_{p}
]
\right),
\label{eq:routing-gate}
\end{equation}

where $\mathbf{e}_{c}$ and $\mathbf{e}_{l}$ are learnable condition and level
embeddings and $\mathbf{a}_{p}$ denotes optional acquisition/context metadata
that is available without leaking the target label. The initial implementation
uses a lightweight MLP or target-conditioned attention gate. A more complex
mixture-of-experts router is considered only if the simpler mechanism is
insufficient.

\subsection{Explicit Availability Mask}
\label{sec:method-routing-mask}

Missing sequences are represented by an availability mask
$a_{p,m}\in\{0,1\}$. Logits for unavailable modalities are masked before the
softmax:

\begin{equation}
g_{p,t,m}=0 \quad \mathrm{if}\quad a_{p,m}=0.
\label{eq:routing-missing-mask}
\end{equation}

The remaining weights are renormalised over the modalities that actually
exist. A missing sequence is therefore not represented by an all-zero image
whose meaning the network must infer. Availability is explicit in the model.

\subsection{Modality Dropout}
\label{sec:method-modality-dropout}

During training, complete studies are stochastically converted into incomplete
ones by removing one or more sequence streams. The dropout distribution includes:

\begin{itemize}
  \item all three sequences;
  \item sagittal T1 + sagittal T2;
  \item sagittal T2 + axial T2;
  \item sagittal T1 + axial T2 where clinically meaningful;
  \item each single-sequence configuration as an edge case.
\end{itemize}

The sampling probabilities are tuned so that the model sees enough complete
studies to learn complementary fusion while also receiving substantial exposure
to incomplete inputs. The exact probabilities are \torecord{final modality-
dropout schedule selected on validation data}.

No external test case is artificially repaired by copying another sequence into
a missing slot. The model's purpose is to remain functional under genuine
absence, not to conceal it.

\subsection{Routing Consistency}
\label{sec:method-routing-consistency}

An optional consistency term encourages the prediction from an incomplete
modality set to remain close to the complete-study prediction when the removed
modality is non-essential:

\begin{equation}
\mathcal{L}_{\mathrm{miss}}
=
D\!\left(
p_{\theta}(y\mid x_{\mathrm{full}}),
p_{\theta}(y\mid x_{\mathrm{subset}})
\right),
\label{eq:missing-consistency}
\end{equation}

where $D$ is a divergence such as KL divergence or Jensen--Shannon divergence.
This term is not applied blindly: for targets known to depend strongly on a
specific sequence, forcing invariance to its removal could suppress legitimate
diagnostic information. The term is therefore ablated by target and compared
with modality-dropout training alone.

\subsection{Routing Baselines}
\label{sec:method-routing-baselines}

RQ3 is evaluated against four progressively stronger controls:

\begin{enumerate}
  \item fixed feature concatenation;
  \item simple equal-weight averaging;
  \item ordinary cross-attention over available sequence features;
  \item disease-conditioned routing without modality dropout;
  \item disease-conditioned routing with modality dropout.
\end{enumerate}

For each method, encoder capacity, ROI inputs and training budget are matched as
closely as possible. Missing-sequence robustness is evaluated on the same
patients under the same artificial absence pattern.

\subsection{Interpreting Routing Weights}
\label{sec:method-routing-interpretation}

The gate weights are model allocations, not causal estimates of sequence
importance. A high axial T2 weight does not prove that axial T2 caused a correct
decision. Clinical plausibility is assessed in two complementary ways.

First, aggregate routing weights are summarised by target. The expected pattern
is that foraminal targets allocate more weight to sagittal T1 and subarticular/
canal targets more to axial T2, but this expectation is not built into the gate
as a hard prior.

Second, controlled input ablation tests whether removing the sequence with high
routing weight causes a greater performance loss than removing a low-weight
sequence. Agreement between learned allocation and intervention strengthens the
interpretation; disagreement is reported rather than rationalised post hoc.

% ===========================================================================
\section{Core Contribution III: Heterogeneous Disease--Anatomy Graph}
\label{sec:method-graph}

\subsection{Graph Construction}
\label{sec:method-graph-construction}

The proposed graph represents \emph{prediction targets}, not pixels and not a
3D disc surface. At full LumbarDISC scope, the node set is

\begin{equation}
V=\{v_{l,c}:l\in\mathcal{L},\,c\in\mathcal{C}\},
\qquad |V|=25.
\label{eq:graph-nodes}
\end{equation}

Each node begins with the routed visual representation for that target together
with embeddings that identify level and condition:

\begin{equation}
\mathbf{h}^{(0)}_{l,c}
=
[
\mathbf{u}_{l,c}
\Vert
\mathbf{e}_{l}
\Vert
\mathbf{e}_{c}
].
\label{eq:node-initial}
\end{equation}

The graph topology is defined anatomically before training. It is not learned
from correlations in the test data.

\subsection{Typed Edge Families}
\label{sec:method-edge-types}

Three primary edge types are used.

\textbf{Adjacent-level edges} connect the same target type across neighbouring
lumbar levels, for example

\[
v_{L3-L4,\mathrm{Canal}}
\leftrightarrow
v_{L4-L5,\mathrm{Canal}}.
\]

These encode the ordered chain of lumbar motion segments.

\textbf{Same-level cross-condition edges} connect anatomical compartments that
coexist at one level, for example central canal to left/right foraminal and
subarticular targets. The edge expresses that the targets share local
morphology, not that one grade deterministically implies another.

\textbf{Bilateral edges} connect left and right counterparts of the same
condition, for example

\[
v_{L4-L5,\mathrm{LF}}
\leftrightarrow
v_{L4-L5,\mathrm{RF}}.
\]

This represents bilateral anatomical symmetry without forcing equal grades.

Additional cross-type edges are added only if clinically justified in the
pre-specified topology. A denser graph is not automatically a better graph; the
study tests whether explicit relation types add value over homogeneous message
passing.

\subsection{Relation-Aware Message Passing}
\label{sec:method-message-passing}

For layer $q$, node $i$ receives messages from neighbours $j\in\mathcal{N}(i)$.
A relation-aware attention formulation is

\begin{align}
\mathbf{q}_i &= W_Q\mathbf{h}^{(q)}_i,\\
\mathbf{k}_{j,r} &= W^r_K\mathbf{h}^{(q)}_j,\\
\mathbf{v}_{j,r} &= W^r_V\mathbf{h}^{(q)}_j,
\end{align}

with attention

\begin{equation}
\alpha_{ijr}
=
\mathrm{softmax}_{j,r}
\left(
\frac{
\mathbf{q}_i^\top\mathbf{k}_{j,r}
}{
\sqrt{d}
}
+b_r
\right),
\label{eq:relational-attention}
\end{equation}

and update

\begin{equation}
\tilde{\mathbf{h}}^{(q+1)}_i
=
\sum_{(j,r)\in\mathcal{N}(i)}
\alpha_{ijr}\mathbf{v}_{j,r}.
\label{eq:graph-message}
\end{equation}

Residual connection, normalisation and feed-forward transformation produce the
final layer output. Relation-specific transformations $W^r_K,W^r_V$ allow an
adjacent-level neighbour to influence a node differently from its contralateral
counterpart.

The exact graph transformer implementation may be replaced by a computationally
simpler relation-aware GAT if validation shows no benefit from the more complex
version. The contribution is the explicit typed target graph and its controlled
evaluation, not a requirement to use the largest graph architecture.

\subsection{Graph Baselines}
\label{sec:method-graph-baselines}

The graph contribution is tested against:

\begin{enumerate}
  \item \textbf{Independent heads:} no message passing between targets.
  \item \textbf{Ordered level transformer:} features communicate across the five
        lumbar levels but do not form 25 typed condition/laterality nodes.
  \item \textbf{Homogeneous graph:} the same 25 nodes and adjacency pattern but
        one undifferentiated edge type.
  \item \textbf{Heterogeneous graph:} typed adjacent-level, cross-condition and
        bilateral relations.
  \item \textbf{Random/shuffled-edge control:} graph capacity retained while
        anatomical edges are permuted.
\end{enumerate}

The shuffled-edge control is essential. If a graph with arbitrary topology
performs as well as the anatomical graph, the study has shown that extra
message-passing capacity helps, not that the encoded anatomy matters.

\subsection{Edge-Family Ablation}
\label{sec:method-edge-ablation}

The full graph is further decomposed:

\begin{itemize}
  \item remove adjacent-level edges;
  \item remove same-level cross-condition edges;
  \item remove bilateral edges;
  \item retain only one family at a time.
\end{itemize}

Results are reported by target family. It is plausible, for example, that
bilateral edges help lateral-compartment grading but add little to central-canal
classification. A single average graph improvement would conceal this structure.

\subsection{Graph Consistency and Uncertainty}
\label{sec:method-graph-uncertainty}

Graph information is not used to hard-code clinically deterministic rules.
Severe stenosis at L4--L5 does not require moderate stenosis at L3--L4, and the
model is allowed to preserve isolated pathology. Relational message passing is
therefore soft.

For uncertainty analysis, node disagreement can be measured between an
independent-head prediction and the graph-refined prediction, or between
neighbour-conditioned and local evidence. This disagreement is explored as an
auxiliary signal for selective prediction, but it is not equated with calibrated
clinical uncertainty unless validated against error.

% ===========================================================================
\section{Supporting Method: Ordinal and Cost-Sensitive Grading}
\label{sec:method-ordinal}

\subsection{Categorical Baseline}
\label{sec:method-ce}

Standard three-class cross-entropy is retained as the mandatory reference:

\begin{equation}
\mathcal{L}_{\mathrm{CE}}
=
-\sum_{t}m_t
\sum_{k=0}^{2}
\mathbb{I}(y_t=k)\log p_{t,k}.
\label{eq:cross-entropy}
\end{equation}

Class weighting or balanced sampling is applied in separate controlled variants.
This baseline is necessary because previous lumbar work found that explicitly
ordinal objectives did not consistently outperform cross-entropy
\cite{niemeyer2021a}.

\subsection{Ordinal Threshold Formulation}
\label{sec:method-ordinal-threshold}

An ordinal head represents a three-grade target as two ordered threshold
questions:

\[
P(y>0\mid\mathbf{h}), \qquad
P(y>1\mid\mathbf{h}).
\]

A CORAL/CORN-style implementation may be used provided the ordering constraint
is respected. The exact ordinal formulation is selected before the final
ablation and compared directly with categorical cross-entropy under identical
features.

Performance is evaluated not only by exact class agreement but also by error
distance:

\begin{equation}
d_t=|\hat{y}_t-y_t|.
\label{eq:grade-distance}
\end{equation}

The proportions with $d_t=0$, $d_t=1$ and $d_t=2$ are reported separately.

\subsection{Clinically Asymmetric Cost}
\label{sec:method-cost}

A cost matrix $C$ allows distant under-grading to receive a larger training
penalty:

\begin{equation}
C=
\begin{bmatrix}
0 & c_{01} & c_{02}\\
c_{10} & 0 & c_{12}\\
c_{20} & c_{21} & 0
\end{bmatrix},
\label{eq:cost-matrix}
\end{equation}

with the clinically motivated condition

\[
c_{20}>c_{21},
\]

so that Severe$\rightarrow$Normal/Mild is more costly than
Severe$\rightarrow$Moderate. The exact numerical values are not invented by the
model developer. They are either aligned with a defensible benchmark/clinical
weighting or selected through a sensitivity analysis across a small
pre-specified family of matrices.

One implementation uses the expected cost under the predicted distribution:

\begin{equation}
\mathcal{L}_{\mathrm{cost}}
=
\sum_t m_t
\sum_{k=0}^{2}
C_{y_t,k}\,p_{t,k}.
\label{eq:expected-cost}
\end{equation}

The cost term is evaluated alongside, not instead of, standard calibration and
discrimination metrics. A model is not considered better merely because it
optimises the chosen matrix.

\subsection{Class Imbalance}
\label{sec:method-imbalance}

Because severe grades are uncommon in LumbarDISC, the study compares a small
set of established strategies rather than stacking all of them simultaneously:

\begin{itemize}
  \item unweighted cross-entropy;
  \item inverse-frequency or effective-number class weighting;
  \item weighted sampling;
  \item focal loss \cite{lin2017focal};
  \item the cost-sensitive/ordinal variant above.
\end{itemize}

MixUp/CutMix are not the default for ordinal grading because mixed labels may
have no clinical interpretation. Geometric augmentation is constrained so that
left/right labels are updated consistently; horizontal flipping is prohibited
for laterality-dependent targets unless labels and orientation are transformed
together.

% ===========================================================================
\section{Supporting Method: Calibration, Uncertainty and Selective Prediction}
\label{sec:method-calibration}

\subsection{Probability Calibration}
\label{sec:method-temperature}

Raw neural-network probabilities are not treated as calibrated confidence.
Temperature scaling is fitted on the validation set:

\begin{equation}
p_k^{(T)}
=
\frac{
\exp(z_k/T)
}{
\sum_j \exp(z_j/T)
},
\label{eq:temperature}
\end{equation}

where $T>0$ is selected without access to the test labels. Calibration is then
evaluated with Expected Calibration Error, Brier score and reliability diagrams.

\subsection{Predictive Uncertainty}
\label{sec:method-uncertainty}

The primary feasible uncertainty methods are MC dropout and/or a small deep
ensemble. If $S$ stochastic predictions are available, the predictive mean is

\begin{equation}
\bar{\mathbf{p}}
=
\frac{1}{S}\sum_{s=1}^{S}\mathbf{p}^{(s)},
\label{eq:predictive-mean}
\end{equation}

and uncertainty can be summarised by predictive entropy,

\begin{equation}
H(\bar{\mathbf{p}})
=
-\sum_k \bar{p}_k\log\bar{p}_k,
\label{eq:entropy}
\end{equation}

together with between-member variance. The study reports whether uncertainty is
higher on incorrect predictions, distant grade errors, low-agreement targets
and external-domain cases.

\subsection{Selective Prediction}
\label{sec:method-selective}

A selective system abstains when uncertainty exceeds threshold $\delta$:

\[
U(x)>\delta \Rightarrow \mathrm{refer/abstain}.
\]

Rather than choosing one arbitrary threshold, the complete risk--coverage curve
is reported: as increasingly uncertain cases are withheld, what fraction of the
cohort remains covered and what error remains among the retained predictions?
The area under, or summary points from, this curve can compare methods.

The terminology ``refer'' describes model behaviour only. It does not claim a
validated clinical workflow; a prospective reader study would be required to
show that abstention improves patient care.

% ===========================================================================
\section{Composite Training Objective}
\label{sec:method-total-loss}

The complete model may use the objective

\begin{equation}
\mathcal{L}_{\mathrm{total}}
=
\mathcal{L}_{\mathrm{grade}}
+\lambda_{\mathrm{ssl}}\mathcal{L}_{\mathrm{ACSSL}}
+\lambda_{\mathrm{miss}}\mathcal{L}_{\mathrm{miss}}
+\lambda_{\mathrm{cost}}\mathcal{L}_{\mathrm{cost}}
+\lambda_{\mathrm{reg}}\mathcal{L}_{\mathrm{reg}},
\label{eq:total-loss}
\end{equation}

where $\mathcal{L}_{\mathrm{grade}}$ is categorical or ordinal supervised loss,
$\mathcal{L}_{\mathrm{ACSSL}}$ is the anatomical cross-sequence objective,
$\mathcal{L}_{\mathrm{miss}}$ is optional missing-modality consistency,
$\mathcal{L}_{\mathrm{cost}}$ is the optional clinical cost term and
$\mathcal{L}_{\mathrm{reg}}$ collects standard regularisation.

Graph message passing is normally part of the forward model rather than a
separate loss. If an explicit graph-consistency regulariser is introduced, it
receives its own coefficient and ablation.

Every non-zero $\lambda$ is treated as a hyperparameter whose contribution must
be demonstrated. A loss term is removed if it adds no reproducible value. The
final thesis reports the selected values and the search range:
\torecord{$\lambda$ values, temperature, optimiser and regularisation settings}.

% ===========================================================================
\section{Optimisation and Training Protocol}
\label{sec:method-training}

\subsection{Training Phases}
\label{sec:method-training-phases}

Training is separated into conceptually distinct phases:

\begin{enumerate}
  \item \textbf{Anatomical/localisation training or loading.}
        The localiser is trained/reproduced and then fixed for the grading
        experiments unless a joint-training ablation is explicitly conducted.

  \item \textbf{Self-supervised pretraining.}
        Sequence encoders are trained with anatomical cross-sequence pairs using
        development patients only.

  \item \textbf{Supervised grading.}
        Encoders are fine-tuned with the routing and graph modules on
        LumbarDISC labels.

  \item \textbf{Calibration.}
        Temperature/uncertainty parameters are fitted on validation data after
        model selection.

  \item \textbf{External evaluation.}
        The selected model is frozen and evaluated on the fixed Rizgary test
        set.

  \item \textbf{Few-shot adaptation.}
        New model copies are adapted using only the local adaptation pool and
        re-evaluated on the same fixed local test set.
\end{enumerate}

\subsection{Optimiser and Schedules}
\label{sec:method-optimiser}

AdamW is the default optimiser for transformer/graph components; SGD or AdamW
may be compared for CNN baselines if required. Learning rate is selected on the
validation set within a pre-specified range, with warm-up and cosine decay or a
plateau scheduler used consistently across the main ablation. Early stopping is
based on a prevalence-robust validation metric such as macro-F1 or weighted
$\kappa$, not raw accuracy.

The final report records:

\begin{itemize}
  \item optimiser and version;
  \item initial learning rate and scheduler;
  \item batch size / effective batch size;
  \item number of epochs and early-stopping patience;
  \item weight decay and dropout;
  \item input resolution and slice-stack radius;
  \item precision mode (FP32/BF16/FP16);
  \item random seeds;
  \item hardware and software versions;
  \item wall-clock and GPU time.
\end{itemize}

No final hyperparameter value is fabricated in this protocol chapter; each is
\torecord{executed training configuration and model-selection criterion}.

\subsection{Augmentation}
\label{sec:method-augmentation}

Training augmentation is designed to simulate plausible MRI variability without
altering disease anatomy. Candidate transforms include modest intensity scaling,
gamma/contrast change, bias-field simulation, noise, small translation/rotation
and limited elastic deformation. Aggressive transformations that distort the
canal or foraminal morphology beyond plausible acquisition variation are
excluded.

Laterality-aware augmentation updates all side-specific labels. If a horizontal
flip changes anatomical left/right, left and right target labels and graph node
identity are swapped in the same transformation. Augmentation is disabled for
validation and test data except in a separately labelled test-time-augmentation
experiment.

\subsection{Model Selection}
\label{sec:method-model-selection}

Only validation data are used to choose architecture depth, routing
hyperparameters, loss weights, early stopping and calibration. The held-out
public test and fixed Rizgary test sets are not inspected for iterative model
selection. If many exploratory experiments are performed, only the pre-specified
final comparison is used for confirmatory statistical claims; exploratory
results are labelled as such.


% ===========================================================================
\section{Experimental Ladder and Ablation Programme}
\label{sec:method-ablation}

The main experiment is not one comparison between ``AMOG-Net'' and a baseline.
It is a ladder in which each additional assumption is isolated.

\begin{longtable}{p{1.1cm}p{5.0cm}p{7.2cm}}
\toprule
\textbf{ID} & \textbf{Configuration} & \textbf{Question isolated}\\
\midrule
\endfirsthead
\toprule
\textbf{ID} & \textbf{Configuration} & \textbf{Question isolated}\\
\midrule
\endhead
E0 & Independent ROI classifier & How well can anatomically localised targets be graded without relational or adaptive fusion?\\
E1 & + DICOM-aligned multi-sequence ROIs & What is gained by correct cross-plane correspondence under fixed fusion?\\
E2 & + disease-conditioned sequence routing & Does target-specific routing outperform fixed fusion?\\
E3 & + modality dropout / missing-sequence training & Does explicit training for absence improve robustness?\\
E4 & + anatomical cross-sequence SSL & Does anatomy-defined pretraining improve representation and label efficiency?\\
E5 & + homogeneous target graph & Does generic message passing help?\\
E6 & + typed heterogeneous disease--anatomy graph & Do clinically defined edge types matter beyond graph capacity?\\
E7 & + ordinal/cost-sensitive/calibrated heads & Can clinically consequential errors and probability reliability be improved?\\
E8 & Selected complete system + external transfer & Do the selected components survive institutional shift and adapt efficiently?\\
\bottomrule
\end{longtable}

The ordering may be adjusted for engineering practicality, but the comparisons
themselves remain. In particular, E6 must be compared with both E5 and a
shuffled-edge control. Otherwise any improvement can be attributed to additional
parameters rather than anatomical topology.

\subsection{Contribution-Specific Ablations}
\label{sec:method-contribution-ablation}

Three ablation families correspond directly to the doctoral contributions.

\textbf{Anatomical SSL ablation} compares no SSL, generic image SSL and
DICOM-aligned cross-sequence SSL under identical downstream training. It is
repeated at reduced label fractions.

\textbf{Routing ablation} compares fixed fusion, cross-attention, target-
conditioned routing and target-conditioned routing with modality dropout.
Performance is reported under complete and missing-sequence input.

\textbf{Graph ablation} compares independent heads, an ordered five-level
transformer, homogeneous graph, heterogeneous graph and random/shuffled edges.
Edge-family removal further identifies whether adjacency, bilateral symmetry or
same-level cross-condition interaction carries the signal.

\subsection{Negative Results}
\label{sec:method-negative-results}

A component that fails to improve the pre-specified outcome is not silently
removed from the research record. The thesis reports it, including confidence
intervals and target-specific effects. This is particularly important for
ordinal objectives and graph modelling, because both are theoretically
attractive and therefore vulnerable to confirmation bias.

A component enters the final integrated system only if it adds a reproducible
benefit, improves an outcome that the thesis has defined as clinically relevant,
or yields a meaningful robustness/calibration advantage despite similar
aggregate discrimination.

% ===========================================================================
\section{External Validation at Rizgary}
\label{sec:method-external}

\subsection{Definition of Zero-Shot Transfer}
\label{sec:method-zero-shot}

Zero-shot transfer is defined strictly. The selected model, including encoder,
router, graph, classifier and calibration strategy developed on public data, is
frozen before any Rizgary test label is used. The local images are passed
through the same preprocessing and localisation pipeline, with only
non-learned DICOM reconstruction and pre-specified per-volume normalisation
allowed.

No local fine-tuning, threshold optimisation or calibration is permitted in the
zero-shot result. If a local intensity harmonisation method is fitted on Rizgary
data, that constitutes a separate domain-adaptation experiment and is labelled
accordingly.

The primary zero-shot analysis is restricted to the central-canal targets for
which the local reports provide a defensible mapping. The principal comparison
is therefore not the full 25-target benchmark score versus a five-target local
score; it is public held-out versus local external performance on the
\emph{same central-canal task}.

\subsection{Quantifying Domain Shift}
\label{sec:method-domain-shift}

Performance degradation is reported as both absolute and relative change:

\begin{align}
\Delta M &= M_{\mathrm{Rizgary}}-M_{\mathrm{Public}},\\
\Delta M_{\%} &=
100\,
\frac{M_{\mathrm{Rizgary}}-M_{\mathrm{Public}}}
{M_{\mathrm{Public}}},
\end{align}

for each primary metric $M$. The thesis does not interpret a large negative
$\Delta M$ as failed research; the size and structure of the drop are among the
main outcomes of RQ4.

The shift is decomposed by:

\begin{itemize}
  \item lumbar level;
  \item severity class;
  \item available sequence configuration;
  \item scanner/vendor/field strength where the local cohort contains enough
        variation;
  \item slice thickness/in-plane resolution;
  \item localisation success;
  \item report/reference-standard confidence.
\end{itemize}

This analysis separates representational failure from data-quality failure
where possible. For example, if grading declines but localisation remains
stable, the pattern supports the detection-versus-grading asymmetry documented
in Chapter~2.

\subsection{Local Error Review}
\label{sec:method-local-error-review}

A blinded error audit is performed after the zero-shot metrics are locked.
Incorrect local predictions are grouped into:

\begin{itemize}
  \item wrong level/localisation;
  \item adjacent-grade disagreement;
  \item distant Severe$\leftrightarrow$Normal/Mild disagreement;
  \item apparent label/report ambiguity;
  \item missing or atypical sequence;
  \item image-quality/artifact issue;
  \item plausible model error despite adequate evidence.
\end{itemize}

Where reader time permits, a subset of model/report disagreements is re-read
without exposing the model prediction. The purpose is not to alter the test
labels opportunistically but to estimate how much apparent model error may
reflect reference-standard variability.

% ===========================================================================
\section{Few-Shot Domain Adaptation}
\label{sec:method-adaptation}

\subsection{Adaptation Sizes}
\label{sec:method-adaptation-sizes}

The adaptation pool is sampled at

\[
N\in\{10,25,50,100\}
\]

patient-level labelled cases, subject to the final number of eligible local
studies. Each $N$ is sampled multiple times, stratified where feasible by
central-canal severity. The same fixed local test set is used after every
adaptation.

The experiment therefore estimates a function

\begin{equation}
M(N)=
\text{test performance after adaptation with }N\text{ local cases},
\label{eq:adaptation-curve}
\end{equation}

rather than a single ``fine-tuned'' number.

\subsection{Adaptation Strategies}
\label{sec:method-adaptation-strategies}

Four strategies are compared:

\begin{enumerate}
  \item \textbf{No adaptation.} The zero-shot baseline.
  \item \textbf{Non-learned/simple harmonisation.} Intensity preprocessing is
        adjusted using only the adaptation pool where a defensible method exists.
  \item \textbf{Parameter-efficient adaptation.} Adapter layers or low-rank
        updates are trained while most of the public-data representation remains
        fixed.
  \item \textbf{Full fine-tuning.} All or most parameters are updated, where
        computationally feasible and where $N$ is large enough to avoid
        immediate overfitting.
\end{enumerate}

The exact PEFT implementation is chosen after a small validation comparison and
is not itself claimed as novel. Its role is to ask an operational question:
how many local labels are needed when a hospital cannot afford to retrain a full
model from scratch?

\subsection{Annotation-Efficiency Summary}
\label{sec:method-annotation-efficiency}

For each method, the recovery curve reports macro-F1, weighted $\kappa$, severe
recall and calibration as a function of $N$. A useful summary is the area under
the performance-versus-log-annotation curve, but raw points with confidence
intervals remain primary because the curve may be non-linear.

No arbitrary threshold such as ``95\% recovery'' defines success. If 100 local
cases recover little performance, that is a meaningful negative finding about
transportability. If 10 cases recover most of the loss, that is evidence for
annotation-efficient adaptation.

\subsection{Avoiding Adaptation Overfit}
\label{sec:method-adaptation-overfit}

Few-shot training is particularly vulnerable to overfitting and selection bias.
The study therefore:

\begin{itemize}
  \item repeats each $N$ across several sampled adaptation sets;
  \item uses regularisation and early stopping based only on an adaptation
        validation split or internal cross-validation within the adaptation
        pool;
  \item never selects the adaptation epoch from the fixed external test;
  \item reports the variance across sampled local subsets;
  \item compares PEFT with full fine-tuning rather than assuming fewer trainable
        parameters are always superior.
\end{itemize}

% ===========================================================================
\section{Evaluation Metrics}
\label{sec:method-metrics}

No single metric is treated as sufficient because the task is imbalanced,
ordinal and multi-target.

\subsection{Macro F1 and Balanced Accuracy}
\label{sec:method-macro}

For class $k$, precision and recall are

\begin{align}
P_k &= \frac{TP_k}{TP_k+FP_k},\\
R_k &= \frac{TP_k}{TP_k+FN_k},
\end{align}

and

\begin{equation}
F1_k=
2\frac{P_kR_k}{P_k+R_k}.
\end{equation}

Macro-F1 is the unweighted class mean,

\begin{equation}
F1_{\mathrm{macro}}
=
\frac{1}{K}\sum_{k=1}^{K}F1_k,
\label{eq:macro-f1}
\end{equation}

so the severe class cannot be hidden by the Normal/Mild majority.

Balanced accuracy is

\begin{equation}
BA=
\frac{1}{K}\sum_{k=1}^{K}R_k.
\label{eq:balanced-accuracy}
\end{equation}

These are primary discrimination metrics.

\subsection{Quadratic Weighted Kappa}
\label{sec:method-kappa}

Quadratic weighted $\kappa$ is reported because the labels are ordered and the
literature supplies human agreement values in $\kappa$. For observed confusion
matrix $O$ and expected matrix $E$, with weight

\[
w_{ij}=\frac{(i-j)^2}{(K-1)^2},
\]

the statistic is

\begin{equation}
\kappa_w
=
1-
\frac{
\sum_{ij}w_{ij}O_{ij}
}{
\sum_{ij}w_{ij}E_{ij}
}.
\label{eq:weighted-kappa}
\end{equation}

Exact agreement and within-one-grade agreement are also reported so that a high
$\kappa$ is not the only view of ordinal error.

\subsection{Severe-Class Outcomes}
\label{sec:method-severe}

The clinically decisive class is reported separately:

\begin{itemize}
  \item Severe recall/sensitivity;
  \item Severe precision;
  \item Severe-versus-rest AUROC and AUPRC where appropriate;
  \item Severe$\rightarrow$Normal/Mild error rate;
  \item Severe$\rightarrow$Moderate error rate.
\end{itemize}

The separation distinguishes distant under-grading from adjacent boundary error.

\subsection{AUROC and AUPRC}
\label{sec:method-auroc}

For three-class outcomes, AUROC is calculated one-versus-rest and macro-averaged.
AUPRC is reported alongside AUROC for rare Severe classes because it is more
sensitive to prevalence and false positives. Curves are constructed at patient-
target level, while confidence intervals respect patient clustering.

\subsection{Calibration}
\label{sec:method-calibration-metrics}

Calibration evaluation includes:

\begin{itemize}
  \item Brier score;
  \item Expected Calibration Error (ECE);
  \item classwise reliability diagrams;
  \item negative log likelihood where useful;
  \item risk--coverage curves under selective prediction.
\end{itemize}

Calibration is reported both internally and externally because a model can
retain discrimination while becoming overconfident after domain shift.

\subsection{Localisation Metrics}
\label{sec:method-localisation-metrics}

Localisation is evaluated separately using:

\begin{itemize}
  \item mean, median and 95th-percentile Euclidean error in millimetres;
  \item PCK at physically meaningful tolerances;
  \item Dice coefficient for available masks;
  \item surface-distance metrics where mask quality supports them;
  \item complete localisation-failure rate per study.
\end{itemize}

\subsection{Efficiency Metrics}
\label{sec:method-efficiency}

For each model configuration the thesis reports:

\begin{itemize}
  \item trainable and total parameter count;
  \item FLOPs/MACs where they can be estimated consistently;
  \item peak VRAM;
  \item training wall time and GPU hours;
  \item inference time per study and per target;
  \item storage size of weights.
\end{itemize}

Efficiency comparisons use the same hardware and batch/inference conditions.
Timing from different machines is not compared as if it were an architectural
property.

% ===========================================================================
\section{Statistical Analysis}
\label{sec:method-statistics}

\subsection{Patient-Level Bootstrap Confidence Intervals}
\label{sec:method-bootstrap}

Confidence intervals for primary metrics are obtained by bootstrap resampling
\emph{patients}, not individual targets. On each replicate, a patient is sampled
with replacement and all of that patient's target predictions enter together.
This preserves within-study dependence.

The default is a sufficiently large number of replicates to stabilise the
2.5th and 97.5th percentiles; the final number is
\torecord{bootstrap replicate count}. Ninety-five per cent percentile or
bias-corrected intervals are reported consistently.

\subsection{Paired Model Comparisons}
\label{sec:method-paired-tests}

Ablation models are evaluated on the same patients, so comparisons are paired.
For scalar patient-level losses or correctness indicators, a paired permutation
test or Wilcoxon signed-rank test is used where distributional assumptions are
not justified. AUCs are compared with an appropriate paired procedure such as
DeLong's test when its assumptions are satisfied. Bootstrap differences are
used for metrics without a simple analytic test.

The primary quantity is the paired difference

\[
\Delta M=M_A-M_B
\]

with a 95\% confidence interval. Point estimates alone are not used to claim
one component is superior.

\subsection{Multiple Comparisons}
\label{sec:method-multiple}

The ablation programme contains many target-level analyses. A limited set of
primary comparisons is therefore declared in advance:

\begin{enumerate}
  \item anatomical SSL versus generic/self-supervised baseline;
  \item routed+dropout fusion versus ordinary cross-attention;
  \item heterogeneous graph versus independent heads and homogeneous graph;
  \item zero-shot public versus local performance;
  \item PEFT versus no adaptation across annotation sizes.
\end{enumerate}

Target-by-target and subgroup analyses are secondary/exploratory. Where many
formal hypotheses are tested, false-discovery-rate control is applied. The
thesis distinguishes confirmatory comparisons from exploratory patterns instead
of presenting every $p<0.05$ as independent evidence.

\subsection{Repeated Training Seeds}
\label{sec:method-seeds}

Deep-learning optimisation is stochastic. Final ablation comparisons are
therefore repeated across multiple fixed seeds where computationally feasible.
The thesis reports mean and standard deviation/interval across seeds in addition
to the patient-level test uncertainty. A one-off training run is insufficient
to attribute a small improvement to a method.

\subsection{Subgroup Analysis}
\label{sec:method-subgroups}

Subgroup analysis is conducted only where sample size permits and includes:

\begin{itemize}
  \item lumbar level;
  \item severity;
  \item sex/age bands where metadata and governance allow;
  \item scanner/vendor/field strength;
  \item acquisition resolution or slice thickness;
  \item presence/absence of individual sequences.
\end{itemize}

Sparse cells are reported descriptively rather than subjected to unstable
significance testing. The subgroup analysis is intended to expose failure modes,
not to make causal claims about demographic biology.

% ===========================================================================
\section{Robustness Experiments}
\label{sec:method-robustness}

\subsection{Missing-Sequence Matrix}
\label{sec:method-missing-matrix}

The selected public-data models are evaluated under every supported sequence
configuration. For three sequences, the core matrix includes complete input,
three two-sequence combinations and three single-sequence configurations where
the required anatomical target remains observable.

Performance is reported per target family rather than only in aggregate. A
model can remain strong on central canal grading while collapsing on foraminal
narrowing when sagittal T1 is removed; an overall mean would obscure the
clinical meaning.

\subsection{Geometry Perturbation}
\label{sec:method-geometry-robust}

The reliance on DICOM coordinates creates a specific failure mode. Controlled
perturbations therefore test sensitivity to small localisation/geometry errors:
axial slice selection shifted by one or two slices, modest ROI-centre offsets,
and simulated spacing variation. The purpose is not to train on corrupted
geometry by default, but to measure how brittle the system is to the imperfect
localisation expected in practice.

\subsection{Image-Quality and Acquisition Perturbation}
\label{sec:method-acquisition-robust}

Where feasible, synthetic intensity perturbations approximating scanner shift
are applied to the public held-out set: contrast scaling, noise, bias field and
resolution degradation. These do not replace real external validation but help
identify which changes reproduce the pattern observed at Rizgary.

\subsection{Graph Perturbation}
\label{sec:method-graph-robust}

The graph is challenged by:

\begin{itemize}
  \item random edge deletion;
  \item shuffled edge types;
  \item random rewiring with matched edge count;
  \item removal of one relation family.
\end{itemize}

A robust anatomical graph should outperform these controls without depending on
one fragile edge. If random topology performs equivalently, the anatomical
interpretation is rejected.

% ===========================================================================
\section{Interpretability and Failure Analysis}
\label{sec:method-interpretability}

\subsection{Local Visual Evidence}
\label{sec:method-visual-evidence}

Grad-CAM or equivalent saliency may be used as a gross sanity check that the
encoder responds within the relevant ROI, but saliency is not presented as a
proof of reasoning. The stronger explanations in this thesis arise from
structural quantities that are part of the model itself:

\begin{itemize}
  \item which sequence received routing weight for a specific target;
  \item which graph relation contributed to target refinement;
  \item how uncertainty changed before and after relational reasoning;
  \item how prediction changed under controlled removal of a sequence or edge.
\end{itemize}

These are still model diagnostics, not causal explanations, but they are more
directly linked to the proposed mechanisms than a heatmap alone.

\subsection{Case-Level Audit Template}
\label{sec:method-case-audit}

A fixed template is used for qualitative audit of representative correct and
incorrect cases:

\begin{itemize}
  \item ground-truth target and grade;
  \item predicted probability distribution and calibrated grade;
  \item localisation overlay;
  \item available MRI sequences;
  \item routing weights;
  \item dominant incoming graph relations/attention where extractable;
  \item uncertainty;
  \item nearest error category;
  \item whether the discrepancy is within one grade.
\end{itemize}

Cases are sampled from predefined categories rather than cherry-picked from the
best-looking outputs.

% ===========================================================================
\section{Reproducibility and Software Engineering}
\label{sec:method-reproducibility}

\subsection{Environment}
\label{sec:method-environment}

The implementation will use a reproducible Python/PyTorch medical-imaging
environment with DICOM tooling such as pydicom, SimpleITK and/or MONAI.
The final thesis records exact versions of Python, PyTorch, CUDA, cuDNN, GPU
hardware and all key libraries. Environment specification is stored in a
lockfile or container recipe.

\subsection{Experiment Configuration}
\label{sec:method-config}

Each run is configured from a machine-readable file containing:

\begin{itemize}
  \item dataset manifest and split version;
  \item selected sequence IDs;
  \item preprocessing and resampling;
  \item ROI dimensions;
  \item encoder architecture;
  \item SSL/routing/graph settings;
  \item loss coefficients;
  \item optimiser and scheduler;
  \item seed;
  \item checkpoint and early-stopping rule.
\end{itemize}

The configuration is copied into the run directory with metrics and model
weights so that a table in the thesis can be traced to an executable experiment.

\subsection{Data Provenance}
\label{sec:method-provenance}

Raw data are never edited in place. A manifest records the path/identifier and
processing status of every study. Derived crops can be regenerated from the
manifest and DICOM source, and each derived sample retains its pseudonymous
patient, level, condition, side, sequence and source-series identifier.

The local reference matrix is versioned independently of model output. If a
label is corrected after clinical review, the change is documented and the
affected experiments are rerun rather than silently updating test truth.

\subsection{Code and Model Release}
\label{sec:method-release}

Code for the public-data methodology will be released where institutional and
dataset licences permit. Pretrained weights are released if the training
dataset terms allow redistribution. Local Rizgary images and report-derived
individual-level data are not included unless a separate hospital approval
authorises public release.

The repository will include:

\begin{itemize}
  \item preprocessing and geometry reconstruction;
  \item split-generation/verification code;
  \item localisation/ROI pipeline;
  \item SSL, routing and graph modules;
  \item training and evaluation scripts;
  \item configuration files for primary experiments;
  \item instructions for reproducing public benchmark results.
\end{itemize}

% ===========================================================================
\section{Methodological Safeguards and Threats to Validity}
\label{sec:method-threats}

\subsection{Reference-Standard Shift}
\label{sec:method-threat-label}

LumbarDISC uses structured expert grades; Rizgary uses routine narrative
reports. A performance drop can therefore reflect both imaging-domain shift and
reference-standard shift. The thesis mitigates this by restricting transfer to
defensible overlap, auditing report labels and, where possible, re-reading a
stratified subset. It does not claim that the two reference standards are
identical.

\subsection{Single-Centre Local External Validation}
\label{sec:method-threat-single-centre}

Rizgary is a genuinely independent institution but still only one local centre.
External performance therefore demonstrates transport to \emph{that} setting,
not universal Middle Eastern generalisation. The thesis states this boundary
explicitly. A future multi-hospital regional study is required before making a
population-wide deployment claim.

\subsection{Model Complexity}
\label{sec:method-threat-complexity}

The full architecture contains several interacting components. Complexity can
produce apparent gains merely through parameter count and tuning effort. The
ablation ladder, matched backbone, homogeneous graph and shuffled-edge controls
are specifically designed to separate mechanism from capacity.

\subsection{Localisation Error}
\label{sec:method-threat-localisation}

The graph and routing modules cannot recover evidence absent from a bad crop.
Localisation is therefore measured and failures are reported separately.
Downstream results may additionally be stratified by localisation error to
estimate how much grading depends on the enabling stage.

\subsection{Class Imbalance}
\label{sec:method-threat-imbalance}

Severe grades are uncommon, so accuracy can be misleading and few-shot local
subsets may contain few Severe examples. The primary metrics are macro-F1,
balanced accuracy, weighted $\kappa$ and per-class recall, with patient-level
confidence intervals. Adaptation subsets are stratified where feasible and
their severity counts are always reported.

\subsection{Hyperparameter Multiplicity}
\label{sec:method-threat-hyperparameters}

A sufficiently large search can overfit the validation set even without touching
test data. The study limits the number of tunable components in each stage,
uses a fixed primary backbone for the core ablation, records the search space
and freezes choices before external testing.

\subsection{Clinical Utility}
\label{sec:method-threat-clinical}

This methodology evaluates technical grading, calibration and portability. It
does not demonstrate that radiologists become faster, safer or more accurate
when using the model. Such a claim requires a prospective reader-assistance or
impact study with independent clinical approval. The thesis therefore uses the
language of ``decision-support potential'' rather than clinical benefit.

% ===========================================================================
\section{Mapping Research Questions to Evidence}
\label{sec:method-rq-mapping}

\begin{longtable}{p{1.0cm}p{4.8cm}p{4.4cm}p{3.8cm}}
\toprule
\textbf{RQ} & \textbf{Primary comparison} & \textbf{Primary outcomes} & \textbf{Failure interpretation}\\
\midrule
\endfirsthead
\toprule
\textbf{RQ} & \textbf{Primary comparison} & \textbf{Primary outcomes} & \textbf{Failure interpretation}\\
\midrule
\endhead
RQ1 & independent/level transformer/homogeneous graph vs typed heterogeneous graph & macro-F1, $\kappa$, severe recall, edge ablations & graph topology carries little useful signal or local evidence dominates\\
RQ2 & generic pretraining vs anatomical cross-sequence SSL at 10/25/50/100\% labels & macro-F1, $\kappa$, label efficiency, external drop & anatomy-defined positives do not improve transferable representation\\
RQ3 & fixed fusion/cross-attention vs routing $\pm$ modality dropout & missing-sequence degradation, target-specific performance, calibration & fixed fusion is sufficient or learned routing is unstable\\
RQ4 & public held-out vs fixed Rizgary test; no adaptation vs PEFT/full tuning & zero-shot $\Delta$, recovery curve, calibration shift & domain shift is larger or less annotation-efficient than expected\\
RQ5 & CE vs ordinal/cost-aware objectives; calibrated vs uncalibrated prediction & Severe$\rightarrow$Normal rate, grade distance, Brier/ECE, risk--coverage & ordinal/cost formulation offers no advantage over CE\\
\bottomrule
\end{longtable}

This table defines what evidence can answer each question. It also makes the
thesis robust to partial negative results: RQ1 can be answered if the graph does
not help, and RQ2 can be answered if generic pretraining is equal or superior.
The contribution is the controlled experiment and the resulting knowledge, not
a requirement that every proposed mechanism win.

% ===========================================================================
\section{Planned Reporting Standard}
\label{sec:method-reporting}

The final manuscript/thesis reporting will follow the current applicable
medical-AI and diagnostic/prediction reporting checklists at the time of
submission. At minimum the report will include:

\begin{itemize}
  \item complete cohort flow diagrams and inclusion/exclusion reasons;
  \item reference-standard source, annotator qualifications and agreement where
        available;
  \item exact patient-level split logic and IDs reproducible within the approved
        environment;
  \item all primary metrics with confidence intervals;
  \item per-class Severe performance and ordinal error distances;
  \item calibration and external-domain performance;
  \item ablations for every claimed novel component;
  \item repeated-seed variability;
  \item compute/hardware cost;
  \item explicit limitations and negative results;
  \item a statement of what was not tested clinically.
\end{itemize}

The reporting checklist itself will be versioned with the manuscript because
CLAIM/STARD/TRIPOD+AI guidance may be updated during the candidature.

% ===========================================================================
\section{Chapter Summary}
\label{sec:method-summary}

This chapter translates the research gap of Chapter~\ref{ch:litreview} into a
controlled experimental design. The study begins from a structured definition
of lumbar disease: twenty-five level--condition--laterality targets on the
public benchmark rather than a generic image-level class. DICOM patient-space
geometry is preserved so that sagittal and axial evidence can be aligned
physically; localisation and disease-specific 2.5D ROIs are treated as shared
enabling technology; and all data partitions are defined at patient level.

Three methodological contributions are then tested independently. Anatomically
aligned cross-sequence self-supervision asks whether the same spinal level
viewed under different MRI sequences provides a stronger representation-learning
signal than generic augmentation. Disease-conditioned routing asks whether one
model can allocate different sequence evidence to different targets while
remaining functional when modalities are absent. The heterogeneous
disease--anatomy graph asks whether explicit typed relations among adjacent
levels, co-occurring compartments and bilateral targets improve grading beyond
independent heads, an inter-level transformer or a homogeneous graph. Ordinal,
cost-sensitive and uncertainty-aware outputs remain supporting methods and are
retained only if their contribution is measurable.

The translational design is deliberately separated from model development.
Rizgary is not mixed into public training. A fixed local test set establishes
zero-shot domain shift, while a separate adaptation pool supplies controlled
few-shot subsets. The primary local comparison is restricted to central-canal
stenosis, where the benchmark and narrative local reports can support a
defensible common target; local herniation morphology remains a separate task
rather than being forced into the RSNA schema.

Finally, the methodology defines success as an answer to the research
questions, not as attainment of a predetermined accuracy. A graph that adds
nothing, an ordinal loss that fails to beat cross-entropy, or a large
cross-institutional performance drop are all informative if demonstrated under
the same controlled protocol. This principle is what turns the architecture
from an engineering assembly into a doctoral experiment: every claimed
mechanism is falsifiable, every comparison is tied to a research question, and
the external cohort is used to test portability rather than to decorate an
internal result.

\printbibliography

\end{document}
 from the main thesis.
% ===========================================================================

\documentclass[12pt,a4paper]{report}

\usepackage{lmodern}
\usepackage[T1]{fontenc}
\usepackage[utf8]{inputenc}
\usepackage[english]{babel}
\usepackage{microtype}
\usepackage[margin=2.5cm]{geometry}
\usepackage{setspace}
\onehalfspacing

\usepackage{amsmath,amssymb}
\usepackage{booktabs}
\usepackage{longtable}
\usepackage{array}
\usepackage{graphicx}
\usepackage[hidelinks]{hyperref}

\usepackage[
  backend=biber,
  style=numeric-comp,
  sorting=none,
  maxbibnames=6,
  minbibnames=3,
  giveninits=true,
  doi=true,
  url=false,
  isbn=false
]{biblatex}

% Same bibliography used by Chapter 2.
\addbibresource{../lumbar_spine_mri_ai_literature_inventory.bib}

\AtEveryBibitem{%
  \clearfield{note}%
  \clearfield{comment}%
  \clearfield{annotation}%
  \clearfield{abstract}%
  \clearfield{file}%
}

\newcommand{\toconfirm}[1]{\textbf{[TO CONFIRM: #1]}}
\newcommand{\torecord}[1]{\textbf{[TO RECORD AFTER IMPLEMENTATION: #1]}}

\begin{document}

\chapter{Methodology}
\label{ch:methodology}

% ===========================================================================
\section{Introduction and Methodological Position}
\label{sec:method-intro}

Chapter~\ref{ch:litreview} established that automated lumbar MRI grading has
moved beyond the stage at which novelty can be claimed from replacing one image
backbone with another. Anatomical localisation, multi-sequence regions of
interest, inter-level context and ordinal grading have all been demonstrated in
recent work. The methodological question in this thesis is therefore not whether
a sufficiently large network can classify lumbar MRI. It is whether the
\emph{structure of the prediction problem} can be represented more faithfully:
which disease is being graded, at which level and side, which MRI sequence
contains the relevant evidence, how neighbouring targets are related, and how
the resulting model behaves when moved to a hospital whose scanners, protocols,
population and reporting conventions were not represented during development.

The study is designed as a staged, controlled experimental programme around
three methodological contributions and one translational validation programme.
The first contribution is an anatomically aligned cross-sequence
self-supervised representation in which positive relationships are defined by
DICOM patient-space correspondence rather than by generic image similarity.
The second is a disease-conditioned sequence router that learns how much each
available MRI sequence contributes to each anatomical target and is trained
explicitly to remain operational when one or more sequences are absent. The
third is a heterogeneous disease--anatomy graph whose nodes correspond to
level--condition--laterality targets and whose typed edges encode anatomical
adjacency, within-level relationships and bilateral symmetry. The translational
programme then separates development from deployment: the model is trained on a
public multinational benchmark, frozen, evaluated zero-shot on an independent
Rizgary Teaching Hospital cohort, and only afterwards adapted with controlled
amounts of local labelled data.

Several components are deliberately treated as \emph{enabling technology}
rather than as doctoral contributions. Anatomical localisation, segmentation,
2.5D region construction, conventional convolutional or transformer encoders,
class-imbalance handling, ordinal losses and calibration are all necessary to
make the comparison scientifically fair, but none is advertised as new by
itself. This distinction follows directly from the novelty boundary established
in Chapter~\ref{ch:litreview}: the thesis tests a new formulation of the
structured task rather than presenting a collection of familiar modules as if
their combination alone constituted originality.

The chapter is written as a protocol-grade methodology. Wherever an
implementation choice has not yet been fixed by evidence, it is specified as a
controlled comparison rather than invented retrospectively. Wherever an
institutional fact remains unresolved---for example the final number of
Rizgary studies surviving data reconciliation or the formal ethics reference---
it is marked explicitly. This is intentional. A doctoral methodology should
describe what is known, what is pre-specified and what will be decided by a
transparent rule; it should not convert planning assumptions into apparently
observed facts.

\subsection{Research Questions}
\label{sec:method-rqs}

The five research questions derived in Chapter~\ref{ch:litreview} are
operationalised as follows.

\begin{description}
  \item[RQ1 --- Relational disease modelling.] Does representing lumbar disease
  targets as a typed heterogeneous graph improve grading over independent target
  heads, an ordered five-level transformer and a homogeneous graph, and are any
  gains concentrated where single-target evidence is weakest?

  \item[RQ2 --- Anatomically aligned self-supervision.] Does cross-sequence
  pretraining based on DICOM-defined anatomical correspondence improve grading,
  label efficiency and transfer compared with ImageNet initialisation, generic
  medical pretraining and ordinary augmentation-based contrastive learning?

  \item[RQ3 --- Disease-conditioned modality routing.] Can one model learn
  target- and patient-dependent sequence weights while remaining robust to
  arbitrary missing-sequence configurations, without requiring a separately
  trained network for every combination of sagittal T1, sagittal T2/STIR and
  axial T2?

  \item[RQ4 --- Cross-institutional transfer.] What is the zero-shot degradation
  when a benchmark-trained model is applied to a previously unseen Middle
  Eastern clinical cohort, and how does performance recover as the amount of
  local annotation available for adaptation increases?

  \item[RQ5 --- Clinically asymmetric error and uncertainty.] Do ordinal and
  cost-sensitive objectives reduce clinically consequential distant errors,
  particularly Severe$\rightarrow$Normal/Mild, and can calibrated uncertainty
  support selective prediction without overstating clinical safety?
\end{description}

The order is not arbitrary. RQ2 and RQ3 determine how evidence is represented;
RQ1 determines how target predictions communicate; RQ5 determines how the
ordered output and its uncertainty are represented; and RQ4 asks whether those
choices survive the transition from a benchmark to real local clinical data.

\subsection{Design Principles}
\label{sec:method-principles}

Six principles govern all experiments.

\textbf{First, anatomy precedes grading.} Every classifier and ablation receives
the same anatomically localised input so that a comparison of representations is
not confounded by one model being given a better crop than another. This follows
the consistent evidence that localisation is a precondition of reliable
per-level grading \cite{lu2018deepspine,pan2021automatically,batra2025mscan}.

\textbf{Second, DICOM geometry is preserved.} The primary scientific
representation is reconstructed from DICOM patient-space coordinates. PNG
conversion may be used for visual inspection but is not the basis of
cross-sequence alignment, because it discards the spatial metadata required to
determine which axial slices correspond to a sagittal disc level.

\textbf{Third, patient identity defines independence.} No images, slices,
levels, sides or sequences from one patient may cross between training,
validation and test partitions. The unit of statistical independence for
uncertainty intervals and resampling is the patient, not the image or the disc
level.

\textbf{Fourth, the external cohort is not a tuning set.} The primary model is
developed without Rizgary images. Zero-shot performance is measured before any
local tuning. Few-shot adaptation then uses a separately defined local
adaptation pool while evaluation remains on a fixed local test set.

\textbf{Fifth, every claimed contribution requires an ablation.} A complete
model score does not establish that anatomical self-supervision, routing or
graph topology contributed anything. Each component is added and removed under
a controlled experiment, and a negative result is reported as such.

\textbf{Sixth, clinical interpretation is bounded by the reference standard.}
The model is trained to reproduce radiological grades whose inter-reader
reliability is imperfect. Performance is therefore interpreted against human
agreement and not as evidence that the network has discovered a more objective
biological truth.

% ===========================================================================
\section{Overall Research Design}
\label{sec:method-design}

\subsection{Study Type}
\label{sec:method-study-type}

The project is a retrospective, multi-dataset medical-imaging methods study with
three linked stages: public benchmark development, controlled methodological
ablation, and independent cross-institutional transfer. It does not prospectively
alter patient care and does not use model output to make treatment decisions.
The Rizgary component is an external retrospective evaluation of imaging and
radiology reports already acquired in routine clinical care.

The core development dataset is the RSNA Lumbar Degenerative Imaging Spine
Classification benchmark (LumbarDISC), which contains multi-sequence MRI and
expert per-level severity labels \cite{rsna2024rsna,richards2026the}. SPIDER is
used as an anatomical localisation/segmentation resource where its licence and
task definition permit \cite{vandergraaf2024lumbar}. The local Rizgary cohort is
reserved for external transfer and domain-adaptation analysis rather than being
mixed into the benchmark training data.

This separation is central to the inference the study is intended to support. A
random split of one pooled dataset can estimate internal generalisation under
that dataset's acquisition distribution. It cannot answer whether a model
trained on international benchmark data transports to a different hospital.
The latter requires the hospital to remain unseen during development.

\subsection{Unit of Analysis}
\label{sec:method-unit}

The primary unit of prediction on LumbarDISC is a
\emph{level--condition--laterality target}. At five disc levels
(L1--L2, L2--L3, L3--L4, L4--L5 and L5--S1), the benchmark defines five targets:
central canal stenosis, left and right neural foraminal narrowing, and left and
right subarticular stenosis. Full benchmark scope therefore contains
$5\times5=25$ targets per examination. Each target receives the ordered label
Normal/Mild, Moderate or Severe \cite{richards2026the}.

The unit of \emph{statistical independence}, however, remains the patient.
Twenty-five outputs from the same examination are correlated measurements, not
twenty-five independent observations. This distinction affects data splitting,
confidence intervals, significance tests and error analysis throughout the
chapter.

For the local Rizgary cohort, the target space is narrower. The reports contain
disc morphology and central-canal descriptions but do not reproduce the entire
RSNA schema reliably. The primary cross-institutional comparison is therefore
restricted to the five central-canal targets, L1--L2 through L5--S1, for which a
defensible overlap can be created. Bulge, protrusion and extrusion are analysed
as a separate local morphology task rather than being relabelled as if they were
RSNA stenosis grades. Foraminal and subarticular transfer are excluded unless
the local reference standard is expanded by new expert grading.

\subsection{Experimental Staging}
\label{sec:method-staging}

The methodology is intentionally cumulative:

\begin{enumerate}
  \item establish reproducible DICOM reconstruction, patient-level splits,
  localisation and an independent ROI baseline;
  \item establish anatomically correct cross-sequence correspondence;
  \item test disease-conditioned sequence routing and missing-modality training;
  \item test anatomically aligned cross-sequence self-supervision;
  \item test homogeneous and heterogeneous graph reasoning against independent
        target prediction;
  \item test ordinal/cost-sensitive output heads and probability calibration;
  \item freeze the selected public-data model and perform zero-shot external
        evaluation at Rizgary;
  \item adapt with prespecified numbers of local cases and estimate the
        annotation--performance recovery curve.
\end{enumerate}

The staging protects the thesis from a common failure mode of complex medical AI
systems: if the final model underperforms, a monolithic experiment cannot reveal
whether the failure arose from localisation, fusion, representation, graph
reasoning, loss design or domain shift. Here each stage creates a measurable
intermediate object and a direct baseline.

\subsection{Pre-Specified Hypotheses}
\label{sec:method-hypotheses}

The principal hypotheses are directional but not expressed as arbitrary target
accuracies.

\begin{description}
  \item[H1] Anatomical cross-sequence self-supervision will improve
  label-efficient grading and cross-institutional robustness relative to generic
  pretraining, particularly when only a fraction of benchmark labels is used.

  \item[H2] Disease-conditioned routing with modality dropout will degrade less
  under missing-sequence experiments than fixed concatenation or ordinary
  cross-attention trained only on complete studies.

  \item[H3] A typed heterogeneous graph will outperform independent target heads
  and a homogeneous graph if the anatomical edge types carry useful relational
  information; performance under a shuffled-edge control should deteriorate if
  the topology itself matters.

  \item[H4] Zero-shot performance on Rizgary will be lower than internal
  benchmark performance because of scanner, protocol, population and reference-
  standard shift; the magnitude and pattern of that drop are treated as findings
  rather than failures of the project.

  \item[H5] Local adaptation will improve performance as annotation increases,
  but no predetermined percentage of ``recovery'' is required. The shape of the
  recovery curve is the result.

  \item[H6] Cost-sensitive/ordinal objectives may reduce distant
  Severe$\rightarrow$Normal/Mild errors even if they do not improve aggregate
  macro-F1; plain cross-entropy remains the mandatory baseline because previous
  lumbar work has shown that ordinal objectives do not automatically help
  \cite{niemeyer2021a}.
\end{description}

% ===========================================================================
\section{Data Sources}
\label{sec:method-data}

\subsection{Public Development Cohort: LumbarDISC}
\label{sec:method-lumbardisc}

The primary development data are drawn from the publicly released RSNA 2024
Lumbar Spine Degenerative Classification dataset. The published LumbarDISC
resource comprises 2,697 patients and 8,593 MRI series collected across eight
institutions in six countries \cite{richards2026the}. The locally held labelled
competition training subset contains approximately 1,975 studies and is used for
core development. The distinction between the published cohort size and the
available labelled subset is retained explicitly throughout the thesis rather
than treating the two numbers as interchangeable.

Eligible benchmark examinations contain sagittal T2-like imaging, sagittal T1
and axial T2, with expert severity grading at the five lumbar levels
\cite{rsna2024rsna,richards2026the}. Coordinate/localiser annotations are also
available for relevant disease locations. These properties make the dataset
appropriate for all three methodological contributions: it is multi-sequence,
level-resolved, multi-target, geographically heterogeneous and sufficiently
large to separate representation learning from local external validation.

The benchmark serves three roles:

\begin{enumerate}
  \item supervised development of per-target grading;
  \item construction of anatomically corresponding cross-sequence training pairs;
  \item internal/site-aware robustness analysis before the model is exposed to
        Rizgary.
\end{enumerate}

All dataset versions, downloaded files, checksums where feasible and the final
patient split lists will be archived in the reproducibility record. The exact
number of studies remaining after local quality-control exclusions will be
\torecord{final LumbarDISC development, validation and held-out counts}.

\subsection{Public Anatomical Resource: SPIDER}
\label{sec:method-spider}

SPIDER provides manually delineated lumbar anatomy on sagittal MRI from multiple
hospitals, including vertebrae, intervertebral discs and spinal canal
\cite{vandergraaf2024lumbar}. It is used only for anatomical localisation,
segmentation benchmarking or pretraining where permitted by its licence. It is
not used to manufacture disease labels that are absent from the dataset.

The practical purpose is to avoid making localisation the doctoral novelty.
Rather than spend the main experimental budget designing another segmentation
network, the thesis will use or reproduce a strong established anatomical model
and quantify its localisation error. If a SPIDER-pretrained model is used to
initialise the localiser, the external disease-grading experiments remain
separate because the disease targets originate in LumbarDISC.

\subsection{Independent Local Cohort: Rizgary Teaching Hospital}
\label{sec:method-rizgary}

The local cohort consists of lumbar MRI examinations and corresponding English
radiology reports supplied by Rizgary Teaching Hospital. Current project
inventory identifies 299 narrative reports and 294 candidate cases with matching
multi-sequence imaging, while a wider raw DICOM audit contains additional
studies. These counts are not treated as final until a one-to-one cohort
reconciliation is complete.

The raw imaging inspected during protocol development confirms the sequence
families required by the thesis---sagittal T1, sagittal T2 and axial T2---and
preserves standard DICOM spatial metadata. A sample examination was acquired on
a Siemens Avanto 1.5-T system; scanner/vendor distribution across the full local
cohort will be reported after complete metadata extraction rather than inferred
from one case.

The local cohort is valuable precisely because it is not curated to reproduce
the public benchmark. Its scanner, acquisition protocol, prevalence, reporting
language and label completeness differ from LumbarDISC. Those differences are
the domain shift the transfer experiment is intended to measure, not nuisances
to be eliminated by silently harmonising labels or discarding difficult cases.

\subsection{Local Radiology Reports and Historical Spreadsheet}
\label{sec:method-local-reports}

The narrative radiology reports are the primary textual source for the local
reference standard. The existing spreadsheet is treated as a secondary
historical transcription aid rather than as unquestioned ground truth. It
contains patient-level variables and comma-separated level references, but it
does not constitute a complete level-resolved RSNA-style matrix and has known
transcription inconsistencies. The methodology therefore reconstructs the local
reference matrix from source reports and uses the spreadsheet only for
cross-checking.

The canonical local matrix will preserve, at minimum:

\begin{itemize}
  \item pseudonymous case identifier;
  \item lumbar level;
  \item central-canal stenosis severity where explicitly stated;
  \item disc bulge, protrusion and extrusion as separate morphology fields;
  \item foraminal findings and side only where explicitly stated;
  \item nerve-root pressure/impingement where reported;
  \item an explicit missing/not-reported state distinct from a negative finding.
\end{itemize}

This last distinction is essential. Absence of a phrase in a narrative report is
not automatically equivalent to ``normal''. A field may be unreported because
the radiologist considered it clinically irrelevant, because reporting practice
varied, or because the report was not structured around that compartment.
``Not reported'' is therefore encoded separately and excluded from target-
specific external evaluation unless a clinician has explicitly adjudicated the
absence as normal.

\subsection{Cohort Reconciliation}
\label{sec:method-reconciliation}

Before any local model evaluation, a case-flow table will reconcile the raw
archive, reports and analysable imaging. For each candidate case the audit will
record:

\begin{itemize}
  \item whether a report exists;
  \item whether sagittal T1 exists and is usable;
  \item whether sagittal T2/STIR exists and is usable;
  \item whether axial T2 exists and is usable;
  \item whether DICOM geometry is internally consistent;
  \item whether a central-canal reference label can be derived;
  \item whether laterality-dependent labels are available;
  \item whether the case is duplicated or represents repeat imaging;
  \item the reason for any exclusion.
\end{itemize}

The final thesis will report the flow from raw studies to the fixed external
test cohort, rather than presenting the current 299/294/raw-study counts without
explanation. Repeat examinations from one patient, if identified, remain in the
same partition to preserve patient independence.

% ===========================================================================
\section{Ethics, Governance and De-Identification}
\label{sec:method-ethics}

The Rizgary component uses retrospective clinical data. No MRI is acquired for
the purposes of this study, no patient undergoes an additional intervention, and
model output is not returned to the clinical record. Nevertheless, the data are
sensitive medical information and are subject to institutional governance.

The final thesis will state the formal approval mechanism and reference:
\toconfirm{hospital research committee/IRB approval number, date and responsible
institution}. It will also state the legal/ethical basis for retrospective use:
\toconfirm{consent waiver, broad consent or other applicable basis}. These items
must be completed before the local study is described as ethically approved.

The raw local DICOM archive contains direct identifiers and therefore cannot be
given to students or uploaded to external compute in its original state.
De-identification is performed on a controlled copy before analysis. At minimum,
the process removes direct patient identifiers, replaces case identifiers with
study pseudonyms, strips or audits private tags, and handles dates according to
the approved policy while preserving only the temporal information required for
scientific analysis. The de-identification procedure must also check for
burned-in identifiers where applicable. The re-identification key, if one must
exist for hospital governance, is stored separately from the research
repository and is not available to the model-training pipeline.

External/cloud processing of local images is prohibited unless explicitly
allowed by hospital policy:
\toconfirm{whether external cloud/GPU processing is permitted and under what
conditions}. Public benchmark data are handled according to their published
licence and competition terms.

Only aggregate results are reported. Illustrative local images may be included
in the thesis only after de-identification review and only if institutional
permission permits publication. The local cohort is not assumed to be publicly
shareable; any future data release would require separate approval and a more
stringent disclosure review.


% ===========================================================================
\section{Reference Standard and Target Encoding}
\label{sec:method-reference}

\subsection{Benchmark Target Representation}
\label{sec:method-benchmark-targets}

For LumbarDISC, the target set is represented explicitly rather than collapsed
into a generic three-class image label. Let the five lumbar levels be
\[
\mathcal{L}=\{L1\!-\!L2,L2\!-\!L3,L3\!-\!L4,L4\!-\!L5,L5\!-\!S1\},
\]
and let the five condition/laterality targets be
\[
\mathcal{C}=\{
\mathrm{Canal},
\mathrm{LF},
\mathrm{RF},
\mathrm{LSA},
\mathrm{RSA}
\},
\]
where LF/RF denote left/right neural foraminal narrowing and LSA/RSA
left/right subarticular stenosis. One examination therefore has the target set
\[
\mathcal{T}=\mathcal{L}\times\mathcal{C}, \qquad |\mathcal{T}|=25.
\]

For target $t\in\mathcal{T}$ the observed grade is
\[
y_t\in\{0,1,2\},
\]
corresponding to Normal/Mild, Moderate and Severe. Encoding the targets in this
form is not merely convenient notation. The pair $(l,c)$ is retained throughout
data loading, feature extraction, graph construction and evaluation, which
prevents a severity label from losing the anatomical meaning that generated it.

Where the benchmark provides an absent or unusable annotation, a target mask
$m_t\in\{0,1\}$ is carried alongside the grade. Masked targets do not contribute
to the supervised loss. Missing labels are never imputed from neighbouring
levels for the main experiments, because doing so would introduce exactly the
relational assumption the thesis is intended to test.

\subsection{Local Reference Matrix}
\label{sec:method-local-reference}

The Rizgary reference standard is constructed at the same anatomical level but
not forced into the same disease schema. A canonical row represents one
case--level pair. The minimum schema is:

\begin{center}
\begin{tabular}{p{2.4cm}p{2.0cm}p{8.1cm}}
\toprule
\textbf{Field} & \textbf{Type} & \textbf{Interpretation}\\
\midrule
Case ID & categorical & pseudonymous patient/study identifier\\
Level & categorical & L1--L2 through L5--S1\\
Canal stenosis & ordinal/missing & explicit report-derived severity, if stated\\
Bulge & binary/missing & circumferential/generalised bulge stated at level\\
Protrusion & binary/missing & protrusion stated at level\\
Extrusion & binary/missing & extrusion stated at level\\
Foraminal finding & categorical/missing & narrowing/pressure effect if stated\\
Laterality & left/right/bilateral/NR & side only where stated\\
Nerve-root effect & categorical/missing & pressure, displacement or compression\\
Other morphology & text/coded & migration, dehydration, height loss and related findings\\
\bottomrule
\end{tabular}
\end{center}

The local matrix keeps \emph{not reported} (NR) separate from \emph{explicitly
absent}. This rule is especially important for foraminal and subarticular
findings, because local reporting completeness differs from the structured RSNA
annotation scheme. A target is eligible for external evaluation only when the
reference standard provides a defensible class label.

\subsection{Report Extraction and Human Verification}
\label{sec:method-report-verification}

Automated report extraction, if available from the parallel clinical-NLP work,
may accelerate construction of the matrix but does not define the gold
standard. The final labels used for Selar's external evaluation are manually
verified against the original report text.

A two-stage verification procedure is preferred:

\begin{enumerate}
  \item \textbf{Structured extraction audit.} Every non-normal/non-empty finding
  is traced back to the exact report phrase, level and laterality. Multi-level
  expressions such as ``L2--3 and L4--5, L5--S1 shows circumferential disc
  bulge'' are expanded into separate level records without changing their
  meaning.

  \item \textbf{Clinical reliability subset.} A pre-specified stratified subset
  of local studies is independently re-read by two qualified readers using an
  explicit grading sheet. Disagreements are adjudicated by a more senior reader
  where available. Inter-reader agreement is reported with weighted
  $\kappa$ for ordinal grades and ordinary agreement/appropriate $\kappa$ for
  binary morphology. The size of the re-read subset will be fixed before model
  results are examined, based on reader availability and the precision required
  for the agreement estimate.
\end{enumerate}

If a second-reader study cannot be completed, the report-derived matrix remains
usable as a real-world reference standard, but the absence of independent
regrading is stated explicitly as a limitation and no claim is made that local
labels reproduce the RSNA grading protocol exactly.

\subsection{Schema Alignment Between LumbarDISC and Rizgary}
\label{sec:method-schema-alignment}

Cross-institutional evaluation requires equality of \emph{task}, not merely
similar words. The mapping is therefore conservative.

\begin{center}
\begin{tabular}{p{4.1cm}p{4.0cm}p{5.0cm}}
\toprule
\textbf{LumbarDISC target} & \textbf{Rizgary evidence} & \textbf{Transfer status}\\
\midrule
Central canal stenosis & level-resolved mild/moderate/severe statements where present
& Primary external target\\
Left/right foraminal narrowing & pressure/narrowing often incompletely lateralised
& Exclude unless laterality and grade are adequate or freshly re-read\\
Left/right subarticular stenosis & not routinely represented in current reports
& Exclude unless new expert annotation is created\\
Disc bulge/protrusion/extrusion & directly described locally but not an RSNA stenosis target
& Separate local morphology task\\
\bottomrule
\end{tabular}
\end{center}

No local disc-herniation morphology is converted into a stenosis grade by
heuristic rule. Likewise, ``pressure on exit foramina'' is not automatically
treated as a left/right foraminal severity label unless the report or expert
review supplies the information required by that target. This prevents
apparent external validation from being created by label re-interpretation.

% ===========================================================================
\section{Data Partitioning and Leakage Control}
\label{sec:method-splits}

\subsection{Patient-Level Independence}
\label{sec:method-patient-split}

All splits are made at patient/study level before slice extraction or data
augmentation. For a patient $p$, every sagittal series, every axial series,
every disc level, every side and every augmented crop belongs to one and only
one partition. A programmatic assertion is run before training to confirm that
the intersection of patient identifiers across partitions is empty.

The split record is version-controlled as a list of pseudonymous IDs. Data
loaders consume those fixed lists rather than performing a new random split each
time the training script is executed.

\subsection{Public Development, Validation and Held-Out Evaluation}
\label{sec:method-public-split}

Where the available public dataset exposes institution identifiers, a
site-aware evaluation is preferred because it tests a more meaningful shift
than a random split. The exact public partitioning scheme will be chosen
according to the metadata available in the held dataset and will be frozen
before model comparison. At minimum there will be:

\begin{itemize}
  \item a development/training partition;
  \item a validation partition for hyperparameter choice, early stopping and
        calibration;
  \item a held-out public evaluation partition not used for model selection.
\end{itemize}

If site identifiers permit leave-one-site-out or institution-held-out analysis,
that analysis is reported in addition to the ordinary benchmark split. It is not
a replacement for Rizgary, because the local hospital introduces a different
reporting/reference-standard domain as well as a different imaging domain.

\subsection{Fixed Local Test Set and Adaptation Pool}
\label{sec:method-local-split}

The local transfer study is protected against adaptation leakage by separating
Rizgary into two parts \emph{before} the public model is evaluated:

\begin{enumerate}
  \item a \textbf{fixed local test set}, which is never used for adaptation,
        calibration or threshold selection; and
  \item a \textbf{local adaptation pool}, from which the $N=10,25,50,100$
        few-shot subsets are sampled.
\end{enumerate}

The exact number allocated to each part will be determined after the audited
class distribution is known. The split will be stratified by central-canal
severity and lumbar level where feasible without producing clinically
implausible or extremely small strata. The fixed local test set is used for the
zero-shot baseline and for every subsequent adapted model, making the recovery
curves directly comparable.

For each adaptation size $N$, several independent stratified subsets are sampled
from the adaptation pool. Results are averaged across these draws with variance
reported, because one set of ten cases can be unrepresentative by chance. The
test set itself never changes across adaptation sizes.

\subsection{Prevention of Information Leakage During Self-Supervision}
\label{sec:method-ssl-leakage}

Self-supervised learning does not remove the need for split discipline. A model
pretrained on the held-out test patients, even without labels, has already seen
their anatomy and acquisition characteristics. Therefore anatomical
self-supervision uses only the development partition during internal
experiments. The held-out public test set and fixed Rizgary test set remain
unseen in both supervised and self-supervised training.

The same restriction applies to intensity normalisation parameters and learned
harmonisation models: any transformation whose parameters are estimated from
data is fitted on the training/development partition and then applied unchanged
to validation/test data.

% ===========================================================================
\section{DICOM Reconstruction and Quality Control}
\label{sec:method-dicom}

\subsection{Series Identification}
\label{sec:method-series-identification}

Each study is first organised by DICOM \texttt{SeriesInstanceUID}. Candidate
series are classified into sagittal T1, sagittal T2/STIR and axial T2 using a
combination of DICOM metadata and image-orientation checks. Series descriptions
are informative but not trusted alone because naming conventions differ across
institutions. The classifier therefore considers fields such as
\texttt{SeriesDescription}, \texttt{ProtocolName}, sequence parameters where
available, and the orientation derived from \texttt{ImageOrientationPatient}.

Ambiguous or duplicated candidate series are flagged for manual review rather
than selected by filename order. The selected series identifiers are stored in
a manifest so that every experiment uses the same source series.

\subsection{Patient-Space Geometry}
\label{sec:method-geometry}

For each DICOM slice, the image plane is reconstructed from
\texttt{ImagePositionPatient}, \texttt{ImageOrientationPatient} and
\texttt{PixelSpacing}. Let $\mathbf{s}$ be the patient-space location of the
upper-left pixel, $\mathbf{r}$ and $\mathbf{c}$ the row and column direction
cosines, and $\Delta_r,\Delta_c$ the corresponding pixel spacings. The physical
coordinate of pixel $(i,j)$ is

\begin{equation}
\mathbf{x}(i,j)=
\mathbf{s}
+i\,\Delta_r\,\mathbf{r}
+j\,\Delta_c\,\mathbf{c}.
\label{eq:dicom-patient-coordinate}
\end{equation}

The slice normal is

\begin{equation}
\mathbf{n}=\mathbf{r}\times\mathbf{c}.
\label{eq:slice-normal}
\end{equation}

Slices are ordered by projection onto $\mathbf{n}$ rather than by filename.
Spacing between consecutive slice origins is inspected for discontinuities.
This procedure preserves patient coordinates across sagittal and axial
acquisitions, allowing a level localised in one plane to be projected into the
other.

For a sagittal target point $\mathbf{q}$ and candidate axial slice $k$ with
origin $\mathbf{s}_k$ and normal $\mathbf{n}_k$, the perpendicular distance to
that slice plane is

\begin{equation}
d_k =
\left|
\mathbf{n}_k^\top(\mathbf{q}-\mathbf{s}_k)
\right|.
\label{eq:plane-distance}
\end{equation}

The nearest anatomically corresponding axial slices are the slices with the
smallest $d_k$, subject to coverage and orientation checks. The method is an
implementation of physical correspondence, not a learned registration shortcut,
and follows the geometric principle used by prior lumbar pipelines
\cite{pan2021automatically,batra2025mscan}.

\subsection{Orientation Standardisation}
\label{sec:method-orientation}

Volumes are reoriented to a common anatomical convention for tensor processing
while retaining the transform to original patient coordinates. Left/right must
never be inferred after arbitrary array transposition; all laterality-dependent
targets are checked against the DICOM orientation transform. This is particularly
important for foraminal and subarticular labels, where a left/right inversion
would create a scientifically silent but catastrophic error.

\subsection{Voxel Spacing and Resampling}
\label{sec:method-resampling}

The acquisition spacing varies by institution, scanner and plane. Inputs are
therefore resampled only after anatomical coordinates have been reconstructed.
The target spacing is chosen from the training-set distribution so that it
preserves structures relevant to grading without creating unnecessary memory
cost. Sagittal and axial streams may use different in-plane resolutions because
their native geometry and diagnostic roles differ.

Interpolation uses a continuous method (e.g.\ linear or higher-order interpolation)
for MR intensities and nearest-neighbour interpolation for discrete masks or
labels. Original spacing is retained in metadata for any measurement reported
in millimetres. The exact target spacing will be \torecord{final sagittal and
axial resampling spacings selected from the development cohort}.

\subsection{Intensity Processing}
\label{sec:method-intensity}

MRI intensities are not quantitatively standardised across scanners. The
preprocessing therefore removes gross scale differences without imposing a
global histogram learned from test data. A default pipeline is:

\begin{enumerate}
  \item suppress non-finite values and obvious padding/background values;
  \item define a foreground mask from the body/spine region;
  \item clip extreme foreground intensities using training-derived robust
        percentiles;
  \item perform per-volume or per-ROI robust z-score normalisation;
  \item optionally compare with bias-field correction or histogram
        standardisation as a robustness ablation.
\end{enumerate}

Any learned harmonisation is fitted only on the development data. The external
Rizgary zero-shot result is reported both with the baseline preprocessing and,
if a separate harmonisation experiment is conducted, with that method clearly
labelled as an intervention rather than folded into the zero-shot definition.

\subsection{Quality-Control Flags}
\label{sec:method-qc}

Automated QC records the following per series:

\begin{itemize}
  \item slice count and ordering consistency;
  \item missing or duplicated instance positions;
  \item orientation consistency within series;
  \item in-plane spacing and slice spacing;
  \item field of view and lumbar coverage;
  \item intensity dynamic range;
  \item severe motion or truncation where detectable;
  \item presence of all sequences required by the particular experiment.
\end{itemize}

Cases failing geometry checks are not silently forced into the pipeline. They
are reviewed, excluded with a recorded reason, or assigned to a robustness
analysis if the failure itself represents a realistic missing-data condition.
A final inclusion/exclusion table will report every such decision.

% ===========================================================================
\section{Anatomical Localisation and ROI Construction}
\label{sec:method-localisation}

\subsection{Role of Localisation}
\label{sec:method-localisation-role}

Localisation exists to make the subsequent comparison fair. The thesis does not
claim that identifying L1--L2 through L5--S1 is itself novel; Chapter~2 showed
that this problem is mature and that strong systems already attain highly
reliable level assignment \cite{pan2021automatically}. The requirement here is
that the localiser be accurate enough that representation and graph experiments
are not dominated by wrong crops.

An established model is therefore preferred. Candidate implementations include
a SPIDER-pretrained nnU-Net-style segmentation network
\cite{isensee2021nnunet,vandergraaf2024lumbar}, a heatmap-regression detector,
or another state-of-the-art anatomical parser available when implementation
begins. The choice is made on validation localisation error and robustness, not
on whether it is architecturally fashionable.

\subsection{Outputs}
\label{sec:method-localisation-output}

The localiser produces, for each study:

\begin{itemize}
  \item centres of L1--L2 through L5--S1 in patient coordinates;
  \item vertebral/disc/canal masks or landmarks where the selected model provides
        them;
  \item a confidence or quality flag for each detected level;
  \item the transform between voxel coordinates and DICOM patient space.
\end{itemize}

A failed or low-confidence localisation is retained as a distinct failure mode.
It is not corrected manually for the primary automated pipeline without being
reported, because doing so would inflate downstream grading performance by
removing a real deployment error.

\subsection{Localisation Evaluation}
\label{sec:method-localisation-eval}

Where point annotations are available, centre localisation error is measured in
millimetres:

\begin{equation}
e_l = \|\hat{\mathbf{q}}_l-\mathbf{q}_l\|_2.
\label{eq:localisation-error}
\end{equation}

Mean/median error and percentile distributions are reported by level.
Percentage-of-correct-keypoint (PCK) or an equivalent tolerance-based measure is
also reported at pre-specified physical thresholds. Where full masks are
available, Dice and surface-distance metrics are reported. These results are
separated from disease grading, consistent with the literature showing that
detection/localisation and fine-grained grading generalise differently.

\subsection{Disease-Specific 2.5D Regions of Interest}
\label{sec:method-roi}

One rectangular crop is not assumed to be equally appropriate for every target.
The crop definition is conditioned on anatomical compartment.

For central canal stenosis, the model receives sagittal T2/STIR context around
the disc level together with corresponding axial T2 slices. For neural
foraminal narrowing, the sagittal T1 stream receives greater parasagittal
coverage and the ROI is side-aware, with sagittal T2 retained as complementary
context. For subarticular stenosis, axial T2 supplies the dominant local
evidence with sagittal context. These choices follow the clinical sequence--
pathology correspondence established in Chapter~2 rather than being learned
from arbitrary crop templates.

Each 2.5D stream contains a small ordered stack around the anatomically closest
slice,

\[
\mathcal{I}_l =
\{I_{k-r},\ldots,I_k,\ldots,I_{k+r}\},
\]

where $k$ is the level-centred slice and $r$ is a small radius determined on the
development set. A five-slice stack ($r=2$) is the initial reference
configuration because it captures through-plane context at modest memory cost;
the final radius is reported after sensitivity testing rather than treated as
a novel contribution.

Crops are defined in physical dimensions where possible, then resampled to the
network input resolution. This avoids a fixed 100-pixel crop representing
different anatomical widths on scanners with different pixel spacing.

\subsection{ROI Quality Control}
\label{sec:method-roi-qc}

A validation subset is visually inspected using overlays of the detected level,
crop boundary and corresponding axial slices. The inspection records:

\begin{itemize}
  \item correct lumbar level;
  \item inclusion of the relevant canal/foramen/recess;
  \item correct left/right orientation;
  \item adequate axial coverage;
  \item crop truncation;
  \item correspondence failure caused by unusual axial angulation.
\end{itemize}

The purpose is not to manually curate the full dataset but to detect systematic
geometry errors before they propagate into model training.

% ===========================================================================
\section{Baseline Representation and Sequence Encoders}
\label{sec:method-baseline}

\subsection{Independent ROI Classifier}
\label{sec:method-e0}

The mandatory baseline (E0) is an independent ROI classifier. Each target is
graded from its anatomically localised input without communication with other
targets. This represents the prevailing formulation criticised in
Chapter~\ref{ch:litreview} and provides the reference against which relational
reasoning is measured.

The baseline uses a well-characterised image encoder such as ResNet
\cite{he2016deep}, EfficientNet \cite{tan2019efficientnet}, Swin
\cite{liu2021swin}, or another contemporary backbone selected before final
experiments. The purpose is not to prove that one family is universally
superior. A primary backbone is selected for the ablation ladder so that
representation, routing and graph changes are measured under one stable
feature extractor. A second backbone may be used in a model-agnosticity check
for the strongest contribution.

\subsection{Sequence-Specific Feature Extraction}
\label{sec:method-sequence-encoders}

Each available MRI sequence is processed by a sequence-specific encoder
$E_m(\cdot)$,

\begin{equation}
\mathbf{f}_{p,l,m}=E_m(\mathcal{I}_{p,l,m}),
\label{eq:sequence-encoder}
\end{equation}

for patient $p$, level $l$ and modality/sequence $m$. The encoders may share
architecture while keeping separate normalisation statistics and weights, or
may partially share weights if an explicit ablation shows that doing so is
beneficial. Separate encoders are the default because T1, sagittal T2/STIR and
axial T2 have different contrast and orientation statistics.

To keep comparisons fair, the dimensionality of the output representation is
held constant across fusion methods. Fixed concatenation, simple averaging and
ordinary cross-attention all consume the same per-sequence features.

\subsection{Controlled Backbone Selection}
\label{sec:method-backbone-control}

Backbone choice is selected on the validation partition before the central
novelty experiments and then held fixed. The candidate comparison will use
identical crops, augmentations, optimiser budget and patient split. If one
backbone is selected because it gives the most stable validation performance
rather than the absolute highest point estimate, that decision is recorded.

This design prevents the thesis from becoming a repeated architecture search.
The main questions concern anatomical self-supervision, adaptive modality
allocation and target-level relational reasoning. A backbone is a controlled
feature extractor in those experiments.


% ===========================================================================
\section{Core Contribution I: Anatomically Aligned Cross-Sequence Self-Supervision}
\label{sec:method-ssl}

\subsection{Rationale}
\label{sec:method-ssl-rationale}

Conventional self-supervised vision learning defines related views by applying
augmentations to one image. That is useful but does not exploit a distinctive
property of multi-sequence MRI: sagittal T1, sagittal T2/STIR and axial T2 can
look very different while representing the same underlying patient and spinal
level. DICOM geometry supplies the correspondence directly.

The proposed objective therefore treats anatomical identity, rather than visual
similarity, as the primary source of positive pairs. For patient $p$, level
$l$, and two sequences $m_a$ and $m_b$, embeddings
$\mathbf{z}_{p,l,m_a}$ and $\mathbf{z}_{p,l,m_b}$ are encouraged to encode a
shared anatomical representation even though acquisition plane and contrast may
differ.

This idea is evaluated as a representation-learning hypothesis, not assumed to
be beneficial. The literature already shows that generic and longitudinal
self-supervision can reduce label dependence
\cite{jamaludin2017selfsupervised,chen2020simple}; the experiment here asks
whether explicitly anatomical cross-sequence correspondence is a stronger
signal for lumbar grading.

\subsection{Projection Head}
\label{sec:method-ssl-projection}

Each sequence feature is mapped through a projection head $P_m$:

\begin{equation}
\mathbf{z}_{p,l,m}=
\frac{
P_m(\mathbf{f}_{p,l,m})
}{
\|P_m(\mathbf{f}_{p,l,m})\|_2
}.
\label{eq:ssl-projection}
\end{equation}

The projection space is used for the self-supervised objective; the encoder
features before projection are retained for downstream grading. This prevents
the representation used for contrastive geometry from constraining the final
classifier unnecessarily.

\subsection{Positive-Pair Construction}
\label{sec:method-positive-pairs}

The strongest positive relation is:

\[
(p,l,m_a)\sim(p,l,m_b),
\qquad m_a\neq m_b,
\]

meaning the same patient and same anatomical level in two different sequences.
Correspondence is established through the DICOM reconstruction and level
localiser, not through nearest image appearance.

An optional softer relation may be tested for adjacent levels from the same
patient, because neighbouring levels share biomechanical context. It is not
included in the primary objective by default. Treating L3--L4 and L4--L5 as
equivalent positives would risk erasing genuine level-specific pathology, so
any adjacent-level term is isolated as a separate ablation.

Different patients at the same level are not defined as positives in the main
objective, because common anatomical level does not imply common disease state.
They may be used in a secondary supervised-contrastive variant once labels are
available, but this is conceptually distinct from the proposed anatomical
self-supervision.

\subsection{Contrastive Objective}
\label{sec:method-ssl-objective}

For anchor embedding $\mathbf{z}_i$ and its cross-sequence positive
$\mathbf{z}_{i^+}$, the reference InfoNCE-style loss is

\begin{equation}
\mathcal{L}_{\mathrm{ACSSL}}(i)
=
-\log
\frac{
\exp(\mathrm{sim}(\mathbf{z}_i,\mathbf{z}_{i^+})/\tau)
}{
\sum_{j\in\mathcal{B}\setminus\{i\}}
\exp(\mathrm{sim}(\mathbf{z}_i,\mathbf{z}_j)/\tau)
},
\label{eq:acssl}
\end{equation}

where $\mathrm{sim}$ is cosine similarity, $\tau$ is a temperature parameter
and $\mathcal{B}$ is the contrastive batch. When more than one anatomically
valid positive exists, the numerator is extended to the set of positives rather
than choosing one arbitrarily.

False negatives are a concern because two different patients may have very
similar normal anatomy. The implementation will therefore compare ordinary
InfoNCE with a variant that does not force all non-positive samples to be hard
negatives, such as a multi-positive or non-contrastive consistency objective.
The novelty claim is tied to the \emph{anatomical pairing}, not to a particular
brand of contrastive loss.

\subsection{Label-Efficiency Experiment}
\label{sec:method-label-efficiency}

The primary test of RQ2 uses fixed fractions of the labelled development set:
10\%, 25\%, 50\% and 100\%. For each fraction:

\begin{enumerate}
  \item the encoder is initialised from the candidate pretraining strategy;
  \item the same supervised grading head is fitted;
  \item patient splits and target distribution are held fixed;
  \item several random labelled subsets are evaluated where computationally
        feasible;
  \item macro-F1, weighted $\kappa$, severe-class recall and calibration are
        compared.
\end{enumerate}

The competing initialisations are:

\begin{itemize}
  \item training from random initialisation;
  \item ImageNet initialisation;
  \item generic medical/volumetric pretraining where a suitable public model is
        available;
  \item ordinary augmentation-based self-supervision;
  \item the proposed anatomically aligned cross-sequence self-supervision.
\end{itemize}

A successful result is not defined as beating every baseline at every label
fraction. The useful scientific question is where anatomical pairing helps,
where it does not, and whether any benefit is largest under label scarcity or
domain shift.

\subsection{Transfer Test of the Representation}
\label{sec:method-ssl-transfer}

The public-data model used for zero-shot Rizgary evaluation is trained once with
and once without anatomical self-supervision while the downstream architecture
is held constant. If the anatomically pretrained model loses less performance
externally despite comparable internal performance, that would support the
claim that the representation is less tied to dataset-specific appearance. If
both models degrade equally, the self-supervised contribution is limited to
label efficiency rather than domain robustness.

% ===========================================================================
\section{Core Contribution II: Disease-Conditioned Adaptive Sequence Routing}
\label{sec:method-routing}

\subsection{Problem Formulation}
\label{sec:method-routing-problem}

A conventional multi-sequence model applies the same fusion rule to every
target. That assumption is clinically implausible because the sequences are not
equally informative for each anatomical compartment. Sagittal T1 is particularly
useful for foraminal fat obliteration, sagittal T2/STIR for disc/canal context,
and axial T2 for canal and lateral recess morphology. The router turns that
known heterogeneity into a learnable, testable mechanism.

For patient $p$, target $t=(l,c)$ and available modalities
$\mathcal{M}_p$, let $\mathbf{f}_{p,l,m}$ be the feature from modality $m$.
The fused target feature is

\begin{equation}
\mathbf{u}_{p,t}
=
\sum_{m\in\mathcal{M}_p}
g_{p,t,m}\,
\mathbf{f}_{p,l,m},
\label{eq:routed-feature}
\end{equation}

where the non-negative routing weights satisfy

\begin{equation}
\sum_{m\in\mathcal{M}_p} g_{p,t,m}=1.
\label{eq:routing-simplex}
\end{equation}

The gate is conditioned on target identity and patient features:

\begin{equation}
g_{p,t,m}
=
\mathrm{softmax}_{m\in\mathcal{M}_p}
\left(
G[
\mathbf{f}_{p,l,m},
\mathbf{e}_{c},
\mathbf{e}_{l},
\mathbf{a}_{p}
]
\right),
\label{eq:routing-gate}
\end{equation}

where $\mathbf{e}_{c}$ and $\mathbf{e}_{l}$ are learnable condition and level
embeddings and $\mathbf{a}_{p}$ denotes optional acquisition/context metadata
that is available without leaking the target label. The initial implementation
uses a lightweight MLP or target-conditioned attention gate. A more complex
mixture-of-experts router is considered only if the simpler mechanism is
insufficient.

\subsection{Explicit Availability Mask}
\label{sec:method-routing-mask}

Missing sequences are represented by an availability mask
$a_{p,m}\in\{0,1\}$. Logits for unavailable modalities are masked before the
softmax:

\begin{equation}
g_{p,t,m}=0 \quad \mathrm{if}\quad a_{p,m}=0.
\label{eq:routing-missing-mask}
\end{equation}

The remaining weights are renormalised over the modalities that actually
exist. A missing sequence is therefore not represented by an all-zero image
whose meaning the network must infer. Availability is explicit in the model.

\subsection{Modality Dropout}
\label{sec:method-modality-dropout}

During training, complete studies are stochastically converted into incomplete
ones by removing one or more sequence streams. The dropout distribution includes:

\begin{itemize}
  \item all three sequences;
  \item sagittal T1 + sagittal T2;
  \item sagittal T2 + axial T2;
  \item sagittal T1 + axial T2 where clinically meaningful;
  \item each single-sequence configuration as an edge case.
\end{itemize}

The sampling probabilities are tuned so that the model sees enough complete
studies to learn complementary fusion while also receiving substantial exposure
to incomplete inputs. The exact probabilities are \torecord{final modality-
dropout schedule selected on validation data}.

No external test case is artificially repaired by copying another sequence into
a missing slot. The model's purpose is to remain functional under genuine
absence, not to conceal it.

\subsection{Routing Consistency}
\label{sec:method-routing-consistency}

An optional consistency term encourages the prediction from an incomplete
modality set to remain close to the complete-study prediction when the removed
modality is non-essential:

\begin{equation}
\mathcal{L}_{\mathrm{miss}}
=
D\!\left(
p_{\theta}(y\mid x_{\mathrm{full}}),
p_{\theta}(y\mid x_{\mathrm{subset}})
\right),
\label{eq:missing-consistency}
\end{equation}

where $D$ is a divergence such as KL divergence or Jensen--Shannon divergence.
This term is not applied blindly: for targets known to depend strongly on a
specific sequence, forcing invariance to its removal could suppress legitimate
diagnostic information. The term is therefore ablated by target and compared
with modality-dropout training alone.

\subsection{Routing Baselines}
\label{sec:method-routing-baselines}

RQ3 is evaluated against four progressively stronger controls:

\begin{enumerate}
  \item fixed feature concatenation;
  \item simple equal-weight averaging;
  \item ordinary cross-attention over available sequence features;
  \item disease-conditioned routing without modality dropout;
  \item disease-conditioned routing with modality dropout.
\end{enumerate}

For each method, encoder capacity, ROI inputs and training budget are matched as
closely as possible. Missing-sequence robustness is evaluated on the same
patients under the same artificial absence pattern.

\subsection{Interpreting Routing Weights}
\label{sec:method-routing-interpretation}

The gate weights are model allocations, not causal estimates of sequence
importance. A high axial T2 weight does not prove that axial T2 caused a correct
decision. Clinical plausibility is assessed in two complementary ways.

First, aggregate routing weights are summarised by target. The expected pattern
is that foraminal targets allocate more weight to sagittal T1 and subarticular/
canal targets more to axial T2, but this expectation is not built into the gate
as a hard prior.

Second, controlled input ablation tests whether removing the sequence with high
routing weight causes a greater performance loss than removing a low-weight
sequence. Agreement between learned allocation and intervention strengthens the
interpretation; disagreement is reported rather than rationalised post hoc.

% ===========================================================================
\section{Core Contribution III: Heterogeneous Disease--Anatomy Graph}
\label{sec:method-graph}

\subsection{Graph Construction}
\label{sec:method-graph-construction}

The proposed graph represents \emph{prediction targets}, not pixels and not a
3D disc surface. At full LumbarDISC scope, the node set is

\begin{equation}
V=\{v_{l,c}:l\in\mathcal{L},\,c\in\mathcal{C}\},
\qquad |V|=25.
\label{eq:graph-nodes}
\end{equation}

Each node begins with the routed visual representation for that target together
with embeddings that identify level and condition:

\begin{equation}
\mathbf{h}^{(0)}_{l,c}
=
[
\mathbf{u}_{l,c}
\Vert
\mathbf{e}_{l}
\Vert
\mathbf{e}_{c}
].
\label{eq:node-initial}
\end{equation}

The graph topology is defined anatomically before training. It is not learned
from correlations in the test data.

\subsection{Typed Edge Families}
\label{sec:method-edge-types}

Three primary edge types are used.

\textbf{Adjacent-level edges} connect the same target type across neighbouring
lumbar levels, for example

\[
v_{L3-L4,\mathrm{Canal}}
\leftrightarrow
v_{L4-L5,\mathrm{Canal}}.
\]

These encode the ordered chain of lumbar motion segments.

\textbf{Same-level cross-condition edges} connect anatomical compartments that
coexist at one level, for example central canal to left/right foraminal and
subarticular targets. The edge expresses that the targets share local
morphology, not that one grade deterministically implies another.

\textbf{Bilateral edges} connect left and right counterparts of the same
condition, for example

\[
v_{L4-L5,\mathrm{LF}}
\leftrightarrow
v_{L4-L5,\mathrm{RF}}.
\]

This represents bilateral anatomical symmetry without forcing equal grades.

Additional cross-type edges are added only if clinically justified in the
pre-specified topology. A denser graph is not automatically a better graph; the
study tests whether explicit relation types add value over homogeneous message
passing.

\subsection{Relation-Aware Message Passing}
\label{sec:method-message-passing}

For layer $q$, node $i$ receives messages from neighbours $j\in\mathcal{N}(i)$.
A relation-aware attention formulation is

\begin{align}
\mathbf{q}_i &= W_Q\mathbf{h}^{(q)}_i,\\
\mathbf{k}_{j,r} &= W^r_K\mathbf{h}^{(q)}_j,\\
\mathbf{v}_{j,r} &= W^r_V\mathbf{h}^{(q)}_j,
\end{align}

with attention

\begin{equation}
\alpha_{ijr}
=
\mathrm{softmax}_{j,r}
\left(
\frac{
\mathbf{q}_i^\top\mathbf{k}_{j,r}
}{
\sqrt{d}
}
+b_r
\right),
\label{eq:relational-attention}
\end{equation}

and update

\begin{equation}
\tilde{\mathbf{h}}^{(q+1)}_i
=
\sum_{(j,r)\in\mathcal{N}(i)}
\alpha_{ijr}\mathbf{v}_{j,r}.
\label{eq:graph-message}
\end{equation}

Residual connection, normalisation and feed-forward transformation produce the
final layer output. Relation-specific transformations $W^r_K,W^r_V$ allow an
adjacent-level neighbour to influence a node differently from its contralateral
counterpart.

The exact graph transformer implementation may be replaced by a computationally
simpler relation-aware GAT if validation shows no benefit from the more complex
version. The contribution is the explicit typed target graph and its controlled
evaluation, not a requirement to use the largest graph architecture.

\subsection{Graph Baselines}
\label{sec:method-graph-baselines}

The graph contribution is tested against:

\begin{enumerate}
  \item \textbf{Independent heads:} no message passing between targets.
  \item \textbf{Ordered level transformer:} features communicate across the five
        lumbar levels but do not form 25 typed condition/laterality nodes.
  \item \textbf{Homogeneous graph:} the same 25 nodes and adjacency pattern but
        one undifferentiated edge type.
  \item \textbf{Heterogeneous graph:} typed adjacent-level, cross-condition and
        bilateral relations.
  \item \textbf{Random/shuffled-edge control:} graph capacity retained while
        anatomical edges are permuted.
\end{enumerate}

The shuffled-edge control is essential. If a graph with arbitrary topology
performs as well as the anatomical graph, the study has shown that extra
message-passing capacity helps, not that the encoded anatomy matters.

\subsection{Edge-Family Ablation}
\label{sec:method-edge-ablation}

The full graph is further decomposed:

\begin{itemize}
  \item remove adjacent-level edges;
  \item remove same-level cross-condition edges;
  \item remove bilateral edges;
  \item retain only one family at a time.
\end{itemize}

Results are reported by target family. It is plausible, for example, that
bilateral edges help lateral-compartment grading but add little to central-canal
classification. A single average graph improvement would conceal this structure.

\subsection{Graph Consistency and Uncertainty}
\label{sec:method-graph-uncertainty}

Graph information is not used to hard-code clinically deterministic rules.
Severe stenosis at L4--L5 does not require moderate stenosis at L3--L4, and the
model is allowed to preserve isolated pathology. Relational message passing is
therefore soft.

For uncertainty analysis, node disagreement can be measured between an
independent-head prediction and the graph-refined prediction, or between
neighbour-conditioned and local evidence. This disagreement is explored as an
auxiliary signal for selective prediction, but it is not equated with calibrated
clinical uncertainty unless validated against error.

% ===========================================================================
\section{Supporting Method: Ordinal and Cost-Sensitive Grading}
\label{sec:method-ordinal}

\subsection{Categorical Baseline}
\label{sec:method-ce}

Standard three-class cross-entropy is retained as the mandatory reference:

\begin{equation}
\mathcal{L}_{\mathrm{CE}}
=
-\sum_{t}m_t
\sum_{k=0}^{2}
\mathbb{I}(y_t=k)\log p_{t,k}.
\label{eq:cross-entropy}
\end{equation}

Class weighting or balanced sampling is applied in separate controlled variants.
This baseline is necessary because previous lumbar work found that explicitly
ordinal objectives did not consistently outperform cross-entropy
\cite{niemeyer2021a}.

\subsection{Ordinal Threshold Formulation}
\label{sec:method-ordinal-threshold}

An ordinal head represents a three-grade target as two ordered threshold
questions:

\[
P(y>0\mid\mathbf{h}), \qquad
P(y>1\mid\mathbf{h}).
\]

A CORAL/CORN-style implementation may be used provided the ordering constraint
is respected. The exact ordinal formulation is selected before the final
ablation and compared directly with categorical cross-entropy under identical
features.

Performance is evaluated not only by exact class agreement but also by error
distance:

\begin{equation}
d_t=|\hat{y}_t-y_t|.
\label{eq:grade-distance}
\end{equation}

The proportions with $d_t=0$, $d_t=1$ and $d_t=2$ are reported separately.

\subsection{Clinically Asymmetric Cost}
\label{sec:method-cost}

A cost matrix $C$ allows distant under-grading to receive a larger training
penalty:

\begin{equation}
C=
\begin{bmatrix}
0 & c_{01} & c_{02}\\
c_{10} & 0 & c_{12}\\
c_{20} & c_{21} & 0
\end{bmatrix},
\label{eq:cost-matrix}
\end{equation}

with the clinically motivated condition

\[
c_{20}>c_{21},
\]

so that Severe$\rightarrow$Normal/Mild is more costly than
Severe$\rightarrow$Moderate. The exact numerical values are not invented by the
model developer. They are either aligned with a defensible benchmark/clinical
weighting or selected through a sensitivity analysis across a small
pre-specified family of matrices.

One implementation uses the expected cost under the predicted distribution:

\begin{equation}
\mathcal{L}_{\mathrm{cost}}
=
\sum_t m_t
\sum_{k=0}^{2}
C_{y_t,k}\,p_{t,k}.
\label{eq:expected-cost}
\end{equation}

The cost term is evaluated alongside, not instead of, standard calibration and
discrimination metrics. A model is not considered better merely because it
optimises the chosen matrix.

\subsection{Class Imbalance}
\label{sec:method-imbalance}

Because severe grades are uncommon in LumbarDISC, the study compares a small
set of established strategies rather than stacking all of them simultaneously:

\begin{itemize}
  \item unweighted cross-entropy;
  \item inverse-frequency or effective-number class weighting;
  \item weighted sampling;
  \item focal loss \cite{lin2017focal};
  \item the cost-sensitive/ordinal variant above.
\end{itemize}

MixUp/CutMix are not the default for ordinal grading because mixed labels may
have no clinical interpretation. Geometric augmentation is constrained so that
left/right labels are updated consistently; horizontal flipping is prohibited
for laterality-dependent targets unless labels and orientation are transformed
together.

% ===========================================================================
\section{Supporting Method: Calibration, Uncertainty and Selective Prediction}
\label{sec:method-calibration}

\subsection{Probability Calibration}
\label{sec:method-temperature}

Raw neural-network probabilities are not treated as calibrated confidence.
Temperature scaling is fitted on the validation set:

\begin{equation}
p_k^{(T)}
=
\frac{
\exp(z_k/T)
}{
\sum_j \exp(z_j/T)
},
\label{eq:temperature}
\end{equation}

where $T>0$ is selected without access to the test labels. Calibration is then
evaluated with Expected Calibration Error, Brier score and reliability diagrams.

\subsection{Predictive Uncertainty}
\label{sec:method-uncertainty}

The primary feasible uncertainty methods are MC dropout and/or a small deep
ensemble. If $S$ stochastic predictions are available, the predictive mean is

\begin{equation}
\bar{\mathbf{p}}
=
\frac{1}{S}\sum_{s=1}^{S}\mathbf{p}^{(s)},
\label{eq:predictive-mean}
\end{equation}

and uncertainty can be summarised by predictive entropy,

\begin{equation}
H(\bar{\mathbf{p}})
=
-\sum_k \bar{p}_k\log\bar{p}_k,
\label{eq:entropy}
\end{equation}

together with between-member variance. The study reports whether uncertainty is
higher on incorrect predictions, distant grade errors, low-agreement targets
and external-domain cases.

\subsection{Selective Prediction}
\label{sec:method-selective}

A selective system abstains when uncertainty exceeds threshold $\delta$:

\[
U(x)>\delta \Rightarrow \mathrm{refer/abstain}.
\]

Rather than choosing one arbitrary threshold, the complete risk--coverage curve
is reported: as increasingly uncertain cases are withheld, what fraction of the
cohort remains covered and what error remains among the retained predictions?
The area under, or summary points from, this curve can compare methods.

The terminology ``refer'' describes model behaviour only. It does not claim a
validated clinical workflow; a prospective reader study would be required to
show that abstention improves patient care.

% ===========================================================================
\section{Composite Training Objective}
\label{sec:method-total-loss}

The complete model may use the objective

\begin{equation}
\mathcal{L}_{\mathrm{total}}
=
\mathcal{L}_{\mathrm{grade}}
+\lambda_{\mathrm{ssl}}\mathcal{L}_{\mathrm{ACSSL}}
+\lambda_{\mathrm{miss}}\mathcal{L}_{\mathrm{miss}}
+\lambda_{\mathrm{cost}}\mathcal{L}_{\mathrm{cost}}
+\lambda_{\mathrm{reg}}\mathcal{L}_{\mathrm{reg}},
\label{eq:total-loss}
\end{equation}

where $\mathcal{L}_{\mathrm{grade}}$ is categorical or ordinal supervised loss,
$\mathcal{L}_{\mathrm{ACSSL}}$ is the anatomical cross-sequence objective,
$\mathcal{L}_{\mathrm{miss}}$ is optional missing-modality consistency,
$\mathcal{L}_{\mathrm{cost}}$ is the optional clinical cost term and
$\mathcal{L}_{\mathrm{reg}}$ collects standard regularisation.

Graph message passing is normally part of the forward model rather than a
separate loss. If an explicit graph-consistency regulariser is introduced, it
receives its own coefficient and ablation.

Every non-zero $\lambda$ is treated as a hyperparameter whose contribution must
be demonstrated. A loss term is removed if it adds no reproducible value. The
final thesis reports the selected values and the search range:
\torecord{$\lambda$ values, temperature, optimiser and regularisation settings}.

% ===========================================================================
\section{Optimisation and Training Protocol}
\label{sec:method-training}

\subsection{Training Phases}
\label{sec:method-training-phases}

Training is separated into conceptually distinct phases:

\begin{enumerate}
  \item \textbf{Anatomical/localisation training or loading.}
        The localiser is trained/reproduced and then fixed for the grading
        experiments unless a joint-training ablation is explicitly conducted.

  \item \textbf{Self-supervised pretraining.}
        Sequence encoders are trained with anatomical cross-sequence pairs using
        development patients only.

  \item \textbf{Supervised grading.}
        Encoders are fine-tuned with the routing and graph modules on
        LumbarDISC labels.

  \item \textbf{Calibration.}
        Temperature/uncertainty parameters are fitted on validation data after
        model selection.

  \item \textbf{External evaluation.}
        The selected model is frozen and evaluated on the fixed Rizgary test
        set.

  \item \textbf{Few-shot adaptation.}
        New model copies are adapted using only the local adaptation pool and
        re-evaluated on the same fixed local test set.
\end{enumerate}

\subsection{Optimiser and Schedules}
\label{sec:method-optimiser}

AdamW is the default optimiser for transformer/graph components; SGD or AdamW
may be compared for CNN baselines if required. Learning rate is selected on the
validation set within a pre-specified range, with warm-up and cosine decay or a
plateau scheduler used consistently across the main ablation. Early stopping is
based on a prevalence-robust validation metric such as macro-F1 or weighted
$\kappa$, not raw accuracy.

The final report records:

\begin{itemize}
  \item optimiser and version;
  \item initial learning rate and scheduler;
  \item batch size / effective batch size;
  \item number of epochs and early-stopping patience;
  \item weight decay and dropout;
  \item input resolution and slice-stack radius;
  \item precision mode (FP32/BF16/FP16);
  \item random seeds;
  \item hardware and software versions;
  \item wall-clock and GPU time.
\end{itemize}

No final hyperparameter value is fabricated in this protocol chapter; each is
\torecord{executed training configuration and model-selection criterion}.

\subsection{Augmentation}
\label{sec:method-augmentation}

Training augmentation is designed to simulate plausible MRI variability without
altering disease anatomy. Candidate transforms include modest intensity scaling,
gamma/contrast change, bias-field simulation, noise, small translation/rotation
and limited elastic deformation. Aggressive transformations that distort the
canal or foraminal morphology beyond plausible acquisition variation are
excluded.

Laterality-aware augmentation updates all side-specific labels. If a horizontal
flip changes anatomical left/right, left and right target labels and graph node
identity are swapped in the same transformation. Augmentation is disabled for
validation and test data except in a separately labelled test-time-augmentation
experiment.

\subsection{Model Selection}
\label{sec:method-model-selection}

Only validation data are used to choose architecture depth, routing
hyperparameters, loss weights, early stopping and calibration. The held-out
public test and fixed Rizgary test sets are not inspected for iterative model
selection. If many exploratory experiments are performed, only the pre-specified
final comparison is used for confirmatory statistical claims; exploratory
results are labelled as such.


% ===========================================================================
\section{Experimental Ladder and Ablation Programme}
\label{sec:method-ablation}

The main experiment is not one comparison between ``AMOG-Net'' and a baseline.
It is a ladder in which each additional assumption is isolated.

\begin{longtable}{p{1.1cm}p{5.0cm}p{7.2cm}}
\toprule
\textbf{ID} & \textbf{Configuration} & \textbf{Question isolated}\\
\midrule
\endfirsthead
\toprule
\textbf{ID} & \textbf{Configuration} & \textbf{Question isolated}\\
\midrule
\endhead
E0 & Independent ROI classifier & How well can anatomically localised targets be graded without relational or adaptive fusion?\\
E1 & + DICOM-aligned multi-sequence ROIs & What is gained by correct cross-plane correspondence under fixed fusion?\\
E2 & + disease-conditioned sequence routing & Does target-specific routing outperform fixed fusion?\\
E3 & + modality dropout / missing-sequence training & Does explicit training for absence improve robustness?\\
E4 & + anatomical cross-sequence SSL & Does anatomy-defined pretraining improve representation and label efficiency?\\
E5 & + homogeneous target graph & Does generic message passing help?\\
E6 & + typed heterogeneous disease--anatomy graph & Do clinically defined edge types matter beyond graph capacity?\\
E7 & + ordinal/cost-sensitive/calibrated heads & Can clinically consequential errors and probability reliability be improved?\\
E8 & Selected complete system + external transfer & Do the selected components survive institutional shift and adapt efficiently?\\
\bottomrule
\end{longtable}

The ordering may be adjusted for engineering practicality, but the comparisons
themselves remain. In particular, E6 must be compared with both E5 and a
shuffled-edge control. Otherwise any improvement can be attributed to additional
parameters rather than anatomical topology.

\subsection{Contribution-Specific Ablations}
\label{sec:method-contribution-ablation}

Three ablation families correspond directly to the doctoral contributions.

\textbf{Anatomical SSL ablation} compares no SSL, generic image SSL and
DICOM-aligned cross-sequence SSL under identical downstream training. It is
repeated at reduced label fractions.

\textbf{Routing ablation} compares fixed fusion, cross-attention, target-
conditioned routing and target-conditioned routing with modality dropout.
Performance is reported under complete and missing-sequence input.

\textbf{Graph ablation} compares independent heads, an ordered five-level
transformer, homogeneous graph, heterogeneous graph and random/shuffled edges.
Edge-family removal further identifies whether adjacency, bilateral symmetry or
same-level cross-condition interaction carries the signal.

\subsection{Negative Results}
\label{sec:method-negative-results}

A component that fails to improve the pre-specified outcome is not silently
removed from the research record. The thesis reports it, including confidence
intervals and target-specific effects. This is particularly important for
ordinal objectives and graph modelling, because both are theoretically
attractive and therefore vulnerable to confirmation bias.

A component enters the final integrated system only if it adds a reproducible
benefit, improves an outcome that the thesis has defined as clinically relevant,
or yields a meaningful robustness/calibration advantage despite similar
aggregate discrimination.

% ===========================================================================
\section{External Validation at Rizgary}
\label{sec:method-external}

\subsection{Definition of Zero-Shot Transfer}
\label{sec:method-zero-shot}

Zero-shot transfer is defined strictly. The selected model, including encoder,
router, graph, classifier and calibration strategy developed on public data, is
frozen before any Rizgary test label is used. The local images are passed
through the same preprocessing and localisation pipeline, with only
non-learned DICOM reconstruction and pre-specified per-volume normalisation
allowed.

No local fine-tuning, threshold optimisation or calibration is permitted in the
zero-shot result. If a local intensity harmonisation method is fitted on Rizgary
data, that constitutes a separate domain-adaptation experiment and is labelled
accordingly.

The primary zero-shot analysis is restricted to the central-canal targets for
which the local reports provide a defensible mapping. The principal comparison
is therefore not the full 25-target benchmark score versus a five-target local
score; it is public held-out versus local external performance on the
\emph{same central-canal task}.

\subsection{Quantifying Domain Shift}
\label{sec:method-domain-shift}

Performance degradation is reported as both absolute and relative change:

\begin{align}
\Delta M &= M_{\mathrm{Rizgary}}-M_{\mathrm{Public}},\\
\Delta M_{\%} &=
100\,
\frac{M_{\mathrm{Rizgary}}-M_{\mathrm{Public}}}
{M_{\mathrm{Public}}},
\end{align}

for each primary metric $M$. The thesis does not interpret a large negative
$\Delta M$ as failed research; the size and structure of the drop are among the
main outcomes of RQ4.

The shift is decomposed by:

\begin{itemize}
  \item lumbar level;
  \item severity class;
  \item available sequence configuration;
  \item scanner/vendor/field strength where the local cohort contains enough
        variation;
  \item slice thickness/in-plane resolution;
  \item localisation success;
  \item report/reference-standard confidence.
\end{itemize}

This analysis separates representational failure from data-quality failure
where possible. For example, if grading declines but localisation remains
stable, the pattern supports the detection-versus-grading asymmetry documented
in Chapter~2.

\subsection{Local Error Review}
\label{sec:method-local-error-review}

A blinded error audit is performed after the zero-shot metrics are locked.
Incorrect local predictions are grouped into:

\begin{itemize}
  \item wrong level/localisation;
  \item adjacent-grade disagreement;
  \item distant Severe$\leftrightarrow$Normal/Mild disagreement;
  \item apparent label/report ambiguity;
  \item missing or atypical sequence;
  \item image-quality/artifact issue;
  \item plausible model error despite adequate evidence.
\end{itemize}

Where reader time permits, a subset of model/report disagreements is re-read
without exposing the model prediction. The purpose is not to alter the test
labels opportunistically but to estimate how much apparent model error may
reflect reference-standard variability.

% ===========================================================================
\section{Few-Shot Domain Adaptation}
\label{sec:method-adaptation}

\subsection{Adaptation Sizes}
\label{sec:method-adaptation-sizes}

The adaptation pool is sampled at

\[
N\in\{10,25,50,100\}
\]

patient-level labelled cases, subject to the final number of eligible local
studies. Each $N$ is sampled multiple times, stratified where feasible by
central-canal severity. The same fixed local test set is used after every
adaptation.

The experiment therefore estimates a function

\begin{equation}
M(N)=
\text{test performance after adaptation with }N\text{ local cases},
\label{eq:adaptation-curve}
\end{equation}

rather than a single ``fine-tuned'' number.

\subsection{Adaptation Strategies}
\label{sec:method-adaptation-strategies}

Four strategies are compared:

\begin{enumerate}
  \item \textbf{No adaptation.} The zero-shot baseline.
  \item \textbf{Non-learned/simple harmonisation.} Intensity preprocessing is
        adjusted using only the adaptation pool where a defensible method exists.
  \item \textbf{Parameter-efficient adaptation.} Adapter layers or low-rank
        updates are trained while most of the public-data representation remains
        fixed.
  \item \textbf{Full fine-tuning.} All or most parameters are updated, where
        computationally feasible and where $N$ is large enough to avoid
        immediate overfitting.
\end{enumerate}

The exact PEFT implementation is chosen after a small validation comparison and
is not itself claimed as novel. Its role is to ask an operational question:
how many local labels are needed when a hospital cannot afford to retrain a full
model from scratch?

\subsection{Annotation-Efficiency Summary}
\label{sec:method-annotation-efficiency}

For each method, the recovery curve reports macro-F1, weighted $\kappa$, severe
recall and calibration as a function of $N$. A useful summary is the area under
the performance-versus-log-annotation curve, but raw points with confidence
intervals remain primary because the curve may be non-linear.

No arbitrary threshold such as ``95\% recovery'' defines success. If 100 local
cases recover little performance, that is a meaningful negative finding about
transportability. If 10 cases recover most of the loss, that is evidence for
annotation-efficient adaptation.

\subsection{Avoiding Adaptation Overfit}
\label{sec:method-adaptation-overfit}

Few-shot training is particularly vulnerable to overfitting and selection bias.
The study therefore:

\begin{itemize}
  \item repeats each $N$ across several sampled adaptation sets;
  \item uses regularisation and early stopping based only on an adaptation
        validation split or internal cross-validation within the adaptation
        pool;
  \item never selects the adaptation epoch from the fixed external test;
  \item reports the variance across sampled local subsets;
  \item compares PEFT with full fine-tuning rather than assuming fewer trainable
        parameters are always superior.
\end{itemize}

% ===========================================================================
\section{Evaluation Metrics}
\label{sec:method-metrics}

No single metric is treated as sufficient because the task is imbalanced,
ordinal and multi-target.

\subsection{Macro F1 and Balanced Accuracy}
\label{sec:method-macro}

For class $k$, precision and recall are

\begin{align}
P_k &= \frac{TP_k}{TP_k+FP_k},\\
R_k &= \frac{TP_k}{TP_k+FN_k},
\end{align}

and

\begin{equation}
F1_k=
2\frac{P_kR_k}{P_k+R_k}.
\end{equation}

Macro-F1 is the unweighted class mean,

\begin{equation}
F1_{\mathrm{macro}}
=
\frac{1}{K}\sum_{k=1}^{K}F1_k,
\label{eq:macro-f1}
\end{equation}

so the severe class cannot be hidden by the Normal/Mild majority.

Balanced accuracy is

\begin{equation}
BA=
\frac{1}{K}\sum_{k=1}^{K}R_k.
\label{eq:balanced-accuracy}
\end{equation}

These are primary discrimination metrics.

\subsection{Quadratic Weighted Kappa}
\label{sec:method-kappa}

Quadratic weighted $\kappa$ is reported because the labels are ordered and the
literature supplies human agreement values in $\kappa$. For observed confusion
matrix $O$ and expected matrix $E$, with weight

\[
w_{ij}=\frac{(i-j)^2}{(K-1)^2},
\]

the statistic is

\begin{equation}
\kappa_w
=
1-
\frac{
\sum_{ij}w_{ij}O_{ij}
}{
\sum_{ij}w_{ij}E_{ij}
}.
\label{eq:weighted-kappa}
\end{equation}

Exact agreement and within-one-grade agreement are also reported so that a high
$\kappa$ is not the only view of ordinal error.

\subsection{Severe-Class Outcomes}
\label{sec:method-severe}

The clinically decisive class is reported separately:

\begin{itemize}
  \item Severe recall/sensitivity;
  \item Severe precision;
  \item Severe-versus-rest AUROC and AUPRC where appropriate;
  \item Severe$\rightarrow$Normal/Mild error rate;
  \item Severe$\rightarrow$Moderate error rate.
\end{itemize}

The separation distinguishes distant under-grading from adjacent boundary error.

\subsection{AUROC and AUPRC}
\label{sec:method-auroc}

For three-class outcomes, AUROC is calculated one-versus-rest and macro-averaged.
AUPRC is reported alongside AUROC for rare Severe classes because it is more
sensitive to prevalence and false positives. Curves are constructed at patient-
target level, while confidence intervals respect patient clustering.

\subsection{Calibration}
\label{sec:method-calibration-metrics}

Calibration evaluation includes:

\begin{itemize}
  \item Brier score;
  \item Expected Calibration Error (ECE);
  \item classwise reliability diagrams;
  \item negative log likelihood where useful;
  \item risk--coverage curves under selective prediction.
\end{itemize}

Calibration is reported both internally and externally because a model can
retain discrimination while becoming overconfident after domain shift.

\subsection{Localisation Metrics}
\label{sec:method-localisation-metrics}

Localisation is evaluated separately using:

\begin{itemize}
  \item mean, median and 95th-percentile Euclidean error in millimetres;
  \item PCK at physically meaningful tolerances;
  \item Dice coefficient for available masks;
  \item surface-distance metrics where mask quality supports them;
  \item complete localisation-failure rate per study.
\end{itemize}

\subsection{Efficiency Metrics}
\label{sec:method-efficiency}

For each model configuration the thesis reports:

\begin{itemize}
  \item trainable and total parameter count;
  \item FLOPs/MACs where they can be estimated consistently;
  \item peak VRAM;
  \item training wall time and GPU hours;
  \item inference time per study and per target;
  \item storage size of weights.
\end{itemize}

Efficiency comparisons use the same hardware and batch/inference conditions.
Timing from different machines is not compared as if it were an architectural
property.

% ===========================================================================
\section{Statistical Analysis}
\label{sec:method-statistics}

\subsection{Patient-Level Bootstrap Confidence Intervals}
\label{sec:method-bootstrap}

Confidence intervals for primary metrics are obtained by bootstrap resampling
\emph{patients}, not individual targets. On each replicate, a patient is sampled
with replacement and all of that patient's target predictions enter together.
This preserves within-study dependence.

The default is a sufficiently large number of replicates to stabilise the
2.5th and 97.5th percentiles; the final number is
\torecord{bootstrap replicate count}. Ninety-five per cent percentile or
bias-corrected intervals are reported consistently.

\subsection{Paired Model Comparisons}
\label{sec:method-paired-tests}

Ablation models are evaluated on the same patients, so comparisons are paired.
For scalar patient-level losses or correctness indicators, a paired permutation
test or Wilcoxon signed-rank test is used where distributional assumptions are
not justified. AUCs are compared with an appropriate paired procedure such as
DeLong's test when its assumptions are satisfied. Bootstrap differences are
used for metrics without a simple analytic test.

The primary quantity is the paired difference

\[
\Delta M=M_A-M_B
\]

with a 95\% confidence interval. Point estimates alone are not used to claim
one component is superior.

\subsection{Multiple Comparisons}
\label{sec:method-multiple}

The ablation programme contains many target-level analyses. A limited set of
primary comparisons is therefore declared in advance:

\begin{enumerate}
  \item anatomical SSL versus generic/self-supervised baseline;
  \item routed+dropout fusion versus ordinary cross-attention;
  \item heterogeneous graph versus independent heads and homogeneous graph;
  \item zero-shot public versus local performance;
  \item PEFT versus no adaptation across annotation sizes.
\end{enumerate}

Target-by-target and subgroup analyses are secondary/exploratory. Where many
formal hypotheses are tested, false-discovery-rate control is applied. The
thesis distinguishes confirmatory comparisons from exploratory patterns instead
of presenting every $p<0.05$ as independent evidence.

\subsection{Repeated Training Seeds}
\label{sec:method-seeds}

Deep-learning optimisation is stochastic. Final ablation comparisons are
therefore repeated across multiple fixed seeds where computationally feasible.
The thesis reports mean and standard deviation/interval across seeds in addition
to the patient-level test uncertainty. A one-off training run is insufficient
to attribute a small improvement to a method.

\subsection{Subgroup Analysis}
\label{sec:method-subgroups}

Subgroup analysis is conducted only where sample size permits and includes:

\begin{itemize}
  \item lumbar level;
  \item severity;
  \item sex/age bands where metadata and governance allow;
  \item scanner/vendor/field strength;
  \item acquisition resolution or slice thickness;
  \item presence/absence of individual sequences.
\end{itemize}

Sparse cells are reported descriptively rather than subjected to unstable
significance testing. The subgroup analysis is intended to expose failure modes,
not to make causal claims about demographic biology.

% ===========================================================================
\section{Robustness Experiments}
\label{sec:method-robustness}

\subsection{Missing-Sequence Matrix}
\label{sec:method-missing-matrix}

The selected public-data models are evaluated under every supported sequence
configuration. For three sequences, the core matrix includes complete input,
three two-sequence combinations and three single-sequence configurations where
the required anatomical target remains observable.

Performance is reported per target family rather than only in aggregate. A
model can remain strong on central canal grading while collapsing on foraminal
narrowing when sagittal T1 is removed; an overall mean would obscure the
clinical meaning.

\subsection{Geometry Perturbation}
\label{sec:method-geometry-robust}

The reliance on DICOM coordinates creates a specific failure mode. Controlled
perturbations therefore test sensitivity to small localisation/geometry errors:
axial slice selection shifted by one or two slices, modest ROI-centre offsets,
and simulated spacing variation. The purpose is not to train on corrupted
geometry by default, but to measure how brittle the system is to the imperfect
localisation expected in practice.

\subsection{Image-Quality and Acquisition Perturbation}
\label{sec:method-acquisition-robust}

Where feasible, synthetic intensity perturbations approximating scanner shift
are applied to the public held-out set: contrast scaling, noise, bias field and
resolution degradation. These do not replace real external validation but help
identify which changes reproduce the pattern observed at Rizgary.

\subsection{Graph Perturbation}
\label{sec:method-graph-robust}

The graph is challenged by:

\begin{itemize}
  \item random edge deletion;
  \item shuffled edge types;
  \item random rewiring with matched edge count;
  \item removal of one relation family.
\end{itemize}

A robust anatomical graph should outperform these controls without depending on
one fragile edge. If random topology performs equivalently, the anatomical
interpretation is rejected.

% ===========================================================================
\section{Interpretability and Failure Analysis}
\label{sec:method-interpretability}

\subsection{Local Visual Evidence}
\label{sec:method-visual-evidence}

Grad-CAM or equivalent saliency may be used as a gross sanity check that the
encoder responds within the relevant ROI, but saliency is not presented as a
proof of reasoning. The stronger explanations in this thesis arise from
structural quantities that are part of the model itself:

\begin{itemize}
  \item which sequence received routing weight for a specific target;
  \item which graph relation contributed to target refinement;
  \item how uncertainty changed before and after relational reasoning;
  \item how prediction changed under controlled removal of a sequence or edge.
\end{itemize}

These are still model diagnostics, not causal explanations, but they are more
directly linked to the proposed mechanisms than a heatmap alone.

\subsection{Case-Level Audit Template}
\label{sec:method-case-audit}

A fixed template is used for qualitative audit of representative correct and
incorrect cases:

\begin{itemize}
  \item ground-truth target and grade;
  \item predicted probability distribution and calibrated grade;
  \item localisation overlay;
  \item available MRI sequences;
  \item routing weights;
  \item dominant incoming graph relations/attention where extractable;
  \item uncertainty;
  \item nearest error category;
  \item whether the discrepancy is within one grade.
\end{itemize}

Cases are sampled from predefined categories rather than cherry-picked from the
best-looking outputs.

% ===========================================================================
\section{Reproducibility and Software Engineering}
\label{sec:method-reproducibility}

\subsection{Environment}
\label{sec:method-environment}

The implementation will use a reproducible Python/PyTorch medical-imaging
environment with DICOM tooling such as pydicom, SimpleITK and/or MONAI.
The final thesis records exact versions of Python, PyTorch, CUDA, cuDNN, GPU
hardware and all key libraries. Environment specification is stored in a
lockfile or container recipe.

\subsection{Experiment Configuration}
\label{sec:method-config}

Each run is configured from a machine-readable file containing:

\begin{itemize}
  \item dataset manifest and split version;
  \item selected sequence IDs;
  \item preprocessing and resampling;
  \item ROI dimensions;
  \item encoder architecture;
  \item SSL/routing/graph settings;
  \item loss coefficients;
  \item optimiser and scheduler;
  \item seed;
  \item checkpoint and early-stopping rule.
\end{itemize}

The configuration is copied into the run directory with metrics and model
weights so that a table in the thesis can be traced to an executable experiment.

\subsection{Data Provenance}
\label{sec:method-provenance}

Raw data are never edited in place. A manifest records the path/identifier and
processing status of every study. Derived crops can be regenerated from the
manifest and DICOM source, and each derived sample retains its pseudonymous
patient, level, condition, side, sequence and source-series identifier.

The local reference matrix is versioned independently of model output. If a
label is corrected after clinical review, the change is documented and the
affected experiments are rerun rather than silently updating test truth.

\subsection{Code and Model Release}
\label{sec:method-release}

Code for the public-data methodology will be released where institutional and
dataset licences permit. Pretrained weights are released if the training
dataset terms allow redistribution. Local Rizgary images and report-derived
individual-level data are not included unless a separate hospital approval
authorises public release.

The repository will include:

\begin{itemize}
  \item preprocessing and geometry reconstruction;
  \item split-generation/verification code;
  \item localisation/ROI pipeline;
  \item SSL, routing and graph modules;
  \item training and evaluation scripts;
  \item configuration files for primary experiments;
  \item instructions for reproducing public benchmark results.
\end{itemize}

% ===========================================================================
\section{Methodological Safeguards and Threats to Validity}
\label{sec:method-threats}

\subsection{Reference-Standard Shift}
\label{sec:method-threat-label}

LumbarDISC uses structured expert grades; Rizgary uses routine narrative
reports. A performance drop can therefore reflect both imaging-domain shift and
reference-standard shift. The thesis mitigates this by restricting transfer to
defensible overlap, auditing report labels and, where possible, re-reading a
stratified subset. It does not claim that the two reference standards are
identical.

\subsection{Single-Centre Local External Validation}
\label{sec:method-threat-single-centre}

Rizgary is a genuinely independent institution but still only one local centre.
External performance therefore demonstrates transport to \emph{that} setting,
not universal Middle Eastern generalisation. The thesis states this boundary
explicitly. A future multi-hospital regional study is required before making a
population-wide deployment claim.

\subsection{Model Complexity}
\label{sec:method-threat-complexity}

The full architecture contains several interacting components. Complexity can
produce apparent gains merely through parameter count and tuning effort. The
ablation ladder, matched backbone, homogeneous graph and shuffled-edge controls
are specifically designed to separate mechanism from capacity.

\subsection{Localisation Error}
\label{sec:method-threat-localisation}

The graph and routing modules cannot recover evidence absent from a bad crop.
Localisation is therefore measured and failures are reported separately.
Downstream results may additionally be stratified by localisation error to
estimate how much grading depends on the enabling stage.

\subsection{Class Imbalance}
\label{sec:method-threat-imbalance}

Severe grades are uncommon, so accuracy can be misleading and few-shot local
subsets may contain few Severe examples. The primary metrics are macro-F1,
balanced accuracy, weighted $\kappa$ and per-class recall, with patient-level
confidence intervals. Adaptation subsets are stratified where feasible and
their severity counts are always reported.

\subsection{Hyperparameter Multiplicity}
\label{sec:method-threat-hyperparameters}

A sufficiently large search can overfit the validation set even without touching
test data. The study limits the number of tunable components in each stage,
uses a fixed primary backbone for the core ablation, records the search space
and freezes choices before external testing.

\subsection{Clinical Utility}
\label{sec:method-threat-clinical}

This methodology evaluates technical grading, calibration and portability. It
does not demonstrate that radiologists become faster, safer or more accurate
when using the model. Such a claim requires a prospective reader-assistance or
impact study with independent clinical approval. The thesis therefore uses the
language of ``decision-support potential'' rather than clinical benefit.

% ===========================================================================
\section{Mapping Research Questions to Evidence}
\label{sec:method-rq-mapping}

\begin{longtable}{p{1.0cm}p{4.8cm}p{4.4cm}p{3.8cm}}
\toprule
\textbf{RQ} & \textbf{Primary comparison} & \textbf{Primary outcomes} & \textbf{Failure interpretation}\\
\midrule
\endfirsthead
\toprule
\textbf{RQ} & \textbf{Primary comparison} & \textbf{Primary outcomes} & \textbf{Failure interpretation}\\
\midrule
\endhead
RQ1 & independent/level transformer/homogeneous graph vs typed heterogeneous graph & macro-F1, $\kappa$, severe recall, edge ablations & graph topology carries little useful signal or local evidence dominates\\
RQ2 & generic pretraining vs anatomical cross-sequence SSL at 10/25/50/100\% labels & macro-F1, $\kappa$, label efficiency, external drop & anatomy-defined positives do not improve transferable representation\\
RQ3 & fixed fusion/cross-attention vs routing $\pm$ modality dropout & missing-sequence degradation, target-specific performance, calibration & fixed fusion is sufficient or learned routing is unstable\\
RQ4 & public held-out vs fixed Rizgary test; no adaptation vs PEFT/full tuning & zero-shot $\Delta$, recovery curve, calibration shift & domain shift is larger or less annotation-efficient than expected\\
RQ5 & CE vs ordinal/cost-aware objectives; calibrated vs uncalibrated prediction & Severe$\rightarrow$Normal rate, grade distance, Brier/ECE, risk--coverage & ordinal/cost formulation offers no advantage over CE\\
\bottomrule
\end{longtable}

This table defines what evidence can answer each question. It also makes the
thesis robust to partial negative results: RQ1 can be answered if the graph does
not help, and RQ2 can be answered if generic pretraining is equal or superior.
The contribution is the controlled experiment and the resulting knowledge, not
a requirement that every proposed mechanism win.

% ===========================================================================
\section{Planned Reporting Standard}
\label{sec:method-reporting}

The final manuscript/thesis reporting will follow the current applicable
medical-AI and diagnostic/prediction reporting checklists at the time of
submission. At minimum the report will include:

\begin{itemize}
  \item complete cohort flow diagrams and inclusion/exclusion reasons;
  \item reference-standard source, annotator qualifications and agreement where
        available;
  \item exact patient-level split logic and IDs reproducible within the approved
        environment;
  \item all primary metrics with confidence intervals;
  \item per-class Severe performance and ordinal error distances;
  \item calibration and external-domain performance;
  \item ablations for every claimed novel component;
  \item repeated-seed variability;
  \item compute/hardware cost;
  \item explicit limitations and negative results;
  \item a statement of what was not tested clinically.
\end{itemize}

The reporting checklist itself will be versioned with the manuscript because
CLAIM/STARD/TRIPOD+AI guidance may be updated during the candidature.

% ===========================================================================
\section{Chapter Summary}
\label{sec:method-summary}

This chapter translates the research gap of Chapter~\ref{ch:litreview} into a
controlled experimental design. The study begins from a structured definition
of lumbar disease: twenty-five level--condition--laterality targets on the
public benchmark rather than a generic image-level class. DICOM patient-space
geometry is preserved so that sagittal and axial evidence can be aligned
physically; localisation and disease-specific 2.5D ROIs are treated as shared
enabling technology; and all data partitions are defined at patient level.

Three methodological contributions are then tested independently. Anatomically
aligned cross-sequence self-supervision asks whether the same spinal level
viewed under different MRI sequences provides a stronger representation-learning
signal than generic augmentation. Disease-conditioned routing asks whether one
model can allocate different sequence evidence to different targets while
remaining functional when modalities are absent. The heterogeneous
disease--anatomy graph asks whether explicit typed relations among adjacent
levels, co-occurring compartments and bilateral targets improve grading beyond
independent heads, an inter-level transformer or a homogeneous graph. Ordinal,
cost-sensitive and uncertainty-aware outputs remain supporting methods and are
retained only if their contribution is measurable.

The translational design is deliberately separated from model development.
Rizgary is not mixed into public training. A fixed local test set establishes
zero-shot domain shift, while a separate adaptation pool supplies controlled
few-shot subsets. The primary local comparison is restricted to central-canal
stenosis, where the benchmark and narrative local reports can support a
defensible common target; local herniation morphology remains a separate task
rather than being forced into the RSNA schema.

Finally, the methodology defines success as an answer to the research
questions, not as attainment of a predetermined accuracy. A graph that adds
nothing, an ordinal loss that fails to beat cross-entropy, or a large
cross-institutional performance drop are all informative if demonstrated under
the same controlled protocol. This principle is what turns the architecture
from an engineering assembly into a doctoral experiment: every claimed
mechanism is falsifiable, every comparison is tied to a research question, and
the external cohort is used to test portability rather than to decorate an
internal result.

\printbibliography

\end{document}
 from the main thesis.
% ===========================================================================

\documentclass[12pt,a4paper]{report}

\usepackage{lmodern}
\usepackage[T1]{fontenc}
\usepackage[utf8]{inputenc}
\usepackage[english]{babel}
\usepackage{microtype}
\usepackage[margin=2.5cm]{geometry}
\usepackage{setspace}
\onehalfspacing

\usepackage{amsmath,amssymb}
\usepackage{booktabs}
\usepackage{longtable}
\usepackage{array}
\usepackage{graphicx}
\usepackage[hidelinks]{hyperref}

\usepackage[
  backend=biber,
  style=numeric-comp,
  sorting=none,
  maxbibnames=6,
  minbibnames=3,
  giveninits=true,
  doi=true,
  url=false,
  isbn=false
]{biblatex}

% Same bibliography used by Chapter 2.
\addbibresource{../lumbar_spine_mri_ai_literature_inventory.bib}

\AtEveryBibitem{%
  \clearfield{note}%
  \clearfield{comment}%
  \clearfield{annotation}%
  \clearfield{abstract}%
  \clearfield{file}%
}

\newcommand{\toconfirm}[1]{\textbf{[TO CONFIRM: #1]}}
\newcommand{\torecord}[1]{\textbf{[TO RECORD AFTER IMPLEMENTATION: #1]}}

\begin{document}

\chapter{Methodology}
\label{ch:methodology}

% ===========================================================================
\section{Introduction and Methodological Position}
\label{sec:method-intro}

Chapter~\ref{ch:litreview} established that automated lumbar MRI grading has
moved beyond the stage at which novelty can be claimed from replacing one image
backbone with another. Anatomical localisation, multi-sequence regions of
interest, inter-level context and ordinal grading have all been demonstrated in
recent work. The methodological question in this thesis is therefore not whether
a sufficiently large network can classify lumbar MRI. It is whether the
\emph{structure of the prediction problem} can be represented more faithfully:
which disease is being graded, at which level and side, which MRI sequence
contains the relevant evidence, how neighbouring targets are related, and how
the resulting model behaves when moved to a hospital whose scanners, protocols,
population and reporting conventions were not represented during development.

The study is designed as a staged, controlled experimental programme around
three methodological contributions and one translational validation programme.
The first contribution is an anatomically aligned cross-sequence
self-supervised representation in which positive relationships are defined by
DICOM patient-space correspondence rather than by generic image similarity.
The second is a disease-conditioned sequence router that learns how much each
available MRI sequence contributes to each anatomical target and is trained
explicitly to remain operational when one or more sequences are absent. The
third is a heterogeneous disease--anatomy graph whose nodes correspond to
level--condition--laterality targets and whose typed edges encode anatomical
adjacency, within-level relationships and bilateral symmetry. The translational
programme then separates development from deployment: the model is trained on a
public multinational benchmark, frozen, evaluated zero-shot on an independent
Rizgary Teaching Hospital cohort, and only afterwards adapted with controlled
amounts of local labelled data.

Several components are deliberately treated as \emph{enabling technology}
rather than as doctoral contributions. Anatomical localisation, segmentation,
2.5D region construction, conventional convolutional or transformer encoders,
class-imbalance handling, ordinal losses and calibration are all necessary to
make the comparison scientifically fair, but none is advertised as new by
itself. This distinction follows directly from the novelty boundary established
in Chapter~\ref{ch:litreview}: the thesis tests a new formulation of the
structured task rather than presenting a collection of familiar modules as if
their combination alone constituted originality.

The chapter is written as a protocol-grade methodology. Wherever an
implementation choice has not yet been fixed by evidence, it is specified as a
controlled comparison rather than invented retrospectively. Wherever an
institutional fact remains unresolved---for example the final number of
Rizgary studies surviving data reconciliation or the formal ethics reference---
it is marked explicitly. This is intentional. A doctoral methodology should
describe what is known, what is pre-specified and what will be decided by a
transparent rule; it should not convert planning assumptions into apparently
observed facts.

\subsection{Research Questions}
\label{sec:method-rqs}

The five research questions derived in Chapter~\ref{ch:litreview} are
operationalised as follows.

\begin{description}
  \item[RQ1 --- Relational disease modelling.] Does representing lumbar disease
  targets as a typed heterogeneous graph improve grading over independent target
  heads, an ordered five-level transformer and a homogeneous graph, and are any
  gains concentrated where single-target evidence is weakest?

  \item[RQ2 --- Anatomically aligned self-supervision.] Does cross-sequence
  pretraining based on DICOM-defined anatomical correspondence improve grading,
  label efficiency and transfer compared with ImageNet initialisation, generic
  medical pretraining and ordinary augmentation-based contrastive learning?

  \item[RQ3 --- Disease-conditioned modality routing.] Can one model learn
  target- and patient-dependent sequence weights while remaining robust to
  arbitrary missing-sequence configurations, without requiring a separately
  trained network for every combination of sagittal T1, sagittal T2/STIR and
  axial T2?

  \item[RQ4 --- Cross-institutional transfer.] What is the zero-shot degradation
  when a benchmark-trained model is applied to a previously unseen Middle
  Eastern clinical cohort, and how does performance recover as the amount of
  local annotation available for adaptation increases?

  \item[RQ5 --- Clinically asymmetric error and uncertainty.] Do ordinal and
  cost-sensitive objectives reduce clinically consequential distant errors,
  particularly Severe$\rightarrow$Normal/Mild, and can calibrated uncertainty
  support selective prediction without overstating clinical safety?
\end{description}

The order is not arbitrary. RQ2 and RQ3 determine how evidence is represented;
RQ1 determines how target predictions communicate; RQ5 determines how the
ordered output and its uncertainty are represented; and RQ4 asks whether those
choices survive the transition from a benchmark to real local clinical data.

\subsection{Design Principles}
\label{sec:method-principles}

Six principles govern all experiments.

\textbf{First, anatomy precedes grading.} Every classifier and ablation receives
the same anatomically localised input so that a comparison of representations is
not confounded by one model being given a better crop than another. This follows
the consistent evidence that localisation is a precondition of reliable
per-level grading \cite{lu2018deepspine,pan2021automatically,batra2025mscan}.

\textbf{Second, DICOM geometry is preserved.} The primary scientific
representation is reconstructed from DICOM patient-space coordinates. PNG
conversion may be used for visual inspection but is not the basis of
cross-sequence alignment, because it discards the spatial metadata required to
determine which axial slices correspond to a sagittal disc level.

\textbf{Third, patient identity defines independence.} No images, slices,
levels, sides or sequences from one patient may cross between training,
validation and test partitions. The unit of statistical independence for
uncertainty intervals and resampling is the patient, not the image or the disc
level.

\textbf{Fourth, the external cohort is not a tuning set.} The primary model is
developed without Rizgary images. Zero-shot performance is measured before any
local tuning. Few-shot adaptation then uses a separately defined local
adaptation pool while evaluation remains on a fixed local test set.

\textbf{Fifth, every claimed contribution requires an ablation.} A complete
model score does not establish that anatomical self-supervision, routing or
graph topology contributed anything. Each component is added and removed under
a controlled experiment, and a negative result is reported as such.

\textbf{Sixth, clinical interpretation is bounded by the reference standard.}
The model is trained to reproduce radiological grades whose inter-reader
reliability is imperfect. Performance is therefore interpreted against human
agreement and not as evidence that the network has discovered a more objective
biological truth.

% ===========================================================================
\section{Overall Research Design}
\label{sec:method-design}

\subsection{Study Type}
\label{sec:method-study-type}

The project is a retrospective, multi-dataset medical-imaging methods study with
three linked stages: public benchmark development, controlled methodological
ablation, and independent cross-institutional transfer. It does not prospectively
alter patient care and does not use model output to make treatment decisions.
The Rizgary component is an external retrospective evaluation of imaging and
radiology reports already acquired in routine clinical care.

The core development dataset is the RSNA Lumbar Degenerative Imaging Spine
Classification benchmark (LumbarDISC), which contains multi-sequence MRI and
expert per-level severity labels \cite{rsna2024rsna,richards2026the}. SPIDER is
used as an anatomical localisation/segmentation resource where its licence and
task definition permit \cite{vandergraaf2024lumbar}. The local Rizgary cohort is
reserved for external transfer and domain-adaptation analysis rather than being
mixed into the benchmark training data.

This separation is central to the inference the study is intended to support. A
random split of one pooled dataset can estimate internal generalisation under
that dataset's acquisition distribution. It cannot answer whether a model
trained on international benchmark data transports to a different hospital.
The latter requires the hospital to remain unseen during development.

\subsection{Unit of Analysis}
\label{sec:method-unit}

The primary unit of prediction on LumbarDISC is a
\emph{level--condition--laterality target}. At five disc levels
(L1--L2, L2--L3, L3--L4, L4--L5 and L5--S1), the benchmark defines five targets:
central canal stenosis, left and right neural foraminal narrowing, and left and
right subarticular stenosis. Full benchmark scope therefore contains
$5\times5=25$ targets per examination. Each target receives the ordered label
Normal/Mild, Moderate or Severe \cite{richards2026the}.

The unit of \emph{statistical independence}, however, remains the patient.
Twenty-five outputs from the same examination are correlated measurements, not
twenty-five independent observations. This distinction affects data splitting,
confidence intervals, significance tests and error analysis throughout the
chapter.

For the local Rizgary cohort, the target space is narrower. The reports contain
disc morphology and central-canal descriptions but do not reproduce the entire
RSNA schema reliably. The primary cross-institutional comparison is therefore
restricted to the five central-canal targets, L1--L2 through L5--S1, for which a
defensible overlap can be created. Bulge, protrusion and extrusion are analysed
as a separate local morphology task rather than being relabelled as if they were
RSNA stenosis grades. Foraminal and subarticular transfer are excluded unless
the local reference standard is expanded by new expert grading.

\subsection{Experimental Staging}
\label{sec:method-staging}

The methodology is intentionally cumulative:

\begin{enumerate}
  \item establish reproducible DICOM reconstruction, patient-level splits,
  localisation and an independent ROI baseline;
  \item establish anatomically correct cross-sequence correspondence;
  \item test disease-conditioned sequence routing and missing-modality training;
  \item test anatomically aligned cross-sequence self-supervision;
  \item test homogeneous and heterogeneous graph reasoning against independent
        target prediction;
  \item test ordinal/cost-sensitive output heads and probability calibration;
  \item freeze the selected public-data model and perform zero-shot external
        evaluation at Rizgary;
  \item adapt with prespecified numbers of local cases and estimate the
        annotation--performance recovery curve.
\end{enumerate}

The staging protects the thesis from a common failure mode of complex medical AI
systems: if the final model underperforms, a monolithic experiment cannot reveal
whether the failure arose from localisation, fusion, representation, graph
reasoning, loss design or domain shift. Here each stage creates a measurable
intermediate object and a direct baseline.

\subsection{Pre-Specified Hypotheses}
\label{sec:method-hypotheses}

The principal hypotheses are directional but not expressed as arbitrary target
accuracies.

\begin{description}
  \item[H1] Anatomical cross-sequence self-supervision will improve
  label-efficient grading and cross-institutional robustness relative to generic
  pretraining, particularly when only a fraction of benchmark labels is used.

  \item[H2] Disease-conditioned routing with modality dropout will degrade less
  under missing-sequence experiments than fixed concatenation or ordinary
  cross-attention trained only on complete studies.

  \item[H3] A typed heterogeneous graph will outperform independent target heads
  and a homogeneous graph if the anatomical edge types carry useful relational
  information; performance under a shuffled-edge control should deteriorate if
  the topology itself matters.

  \item[H4] Zero-shot performance on Rizgary will be lower than internal
  benchmark performance because of scanner, protocol, population and reference-
  standard shift; the magnitude and pattern of that drop are treated as findings
  rather than failures of the project.

  \item[H5] Local adaptation will improve performance as annotation increases,
  but no predetermined percentage of ``recovery'' is required. The shape of the
  recovery curve is the result.

  \item[H6] Cost-sensitive/ordinal objectives may reduce distant
  Severe$\rightarrow$Normal/Mild errors even if they do not improve aggregate
  macro-F1; plain cross-entropy remains the mandatory baseline because previous
  lumbar work has shown that ordinal objectives do not automatically help
  \cite{niemeyer2021a}.
\end{description}

% ===========================================================================
\section{Data Sources}
\label{sec:method-data}

\subsection{Public Development Cohort: LumbarDISC}
\label{sec:method-lumbardisc}

The primary development data are drawn from the publicly released RSNA 2024
Lumbar Spine Degenerative Classification dataset. The published LumbarDISC
resource comprises 2,697 patients and 8,593 MRI series collected across eight
institutions in six countries \cite{richards2026the}. The locally held labelled
competition training subset contains approximately 1,975 studies and is used for
core development. The distinction between the published cohort size and the
available labelled subset is retained explicitly throughout the thesis rather
than treating the two numbers as interchangeable.

Eligible benchmark examinations contain sagittal T2-like imaging, sagittal T1
and axial T2, with expert severity grading at the five lumbar levels
\cite{rsna2024rsna,richards2026the}. Coordinate/localiser annotations are also
available for relevant disease locations. These properties make the dataset
appropriate for all three methodological contributions: it is multi-sequence,
level-resolved, multi-target, geographically heterogeneous and sufficiently
large to separate representation learning from local external validation.

The benchmark serves three roles:

\begin{enumerate}
  \item supervised development of per-target grading;
  \item construction of anatomically corresponding cross-sequence training pairs;
  \item internal/site-aware robustness analysis before the model is exposed to
        Rizgary.
\end{enumerate}

All dataset versions, downloaded files, checksums where feasible and the final
patient split lists will be archived in the reproducibility record. The exact
number of studies remaining after local quality-control exclusions will be
\torecord{final LumbarDISC development, validation and held-out counts}.

\subsection{Public Anatomical Resource: SPIDER}
\label{sec:method-spider}

SPIDER provides manually delineated lumbar anatomy on sagittal MRI from multiple
hospitals, including vertebrae, intervertebral discs and spinal canal
\cite{vandergraaf2024lumbar}. It is used only for anatomical localisation,
segmentation benchmarking or pretraining where permitted by its licence. It is
not used to manufacture disease labels that are absent from the dataset.

The practical purpose is to avoid making localisation the doctoral novelty.
Rather than spend the main experimental budget designing another segmentation
network, the thesis will use or reproduce a strong established anatomical model
and quantify its localisation error. If a SPIDER-pretrained model is used to
initialise the localiser, the external disease-grading experiments remain
separate because the disease targets originate in LumbarDISC.

\subsection{Independent Local Cohort: Rizgary Teaching Hospital}
\label{sec:method-rizgary}

The local cohort consists of lumbar MRI examinations and corresponding English
radiology reports supplied by Rizgary Teaching Hospital. Current project
inventory identifies 299 narrative reports and 294 candidate cases with matching
multi-sequence imaging, while a wider raw DICOM audit contains additional
studies. These counts are not treated as final until a one-to-one cohort
reconciliation is complete.

The raw imaging inspected during protocol development confirms the sequence
families required by the thesis---sagittal T1, sagittal T2 and axial T2---and
preserves standard DICOM spatial metadata. A sample examination was acquired on
a Siemens Avanto 1.5-T system; scanner/vendor distribution across the full local
cohort will be reported after complete metadata extraction rather than inferred
from one case.

The local cohort is valuable precisely because it is not curated to reproduce
the public benchmark. Its scanner, acquisition protocol, prevalence, reporting
language and label completeness differ from LumbarDISC. Those differences are
the domain shift the transfer experiment is intended to measure, not nuisances
to be eliminated by silently harmonising labels or discarding difficult cases.

\subsection{Local Radiology Reports and Historical Spreadsheet}
\label{sec:method-local-reports}

The narrative radiology reports are the primary textual source for the local
reference standard. The existing spreadsheet is treated as a secondary
historical transcription aid rather than as unquestioned ground truth. It
contains patient-level variables and comma-separated level references, but it
does not constitute a complete level-resolved RSNA-style matrix and has known
transcription inconsistencies. The methodology therefore reconstructs the local
reference matrix from source reports and uses the spreadsheet only for
cross-checking.

The canonical local matrix will preserve, at minimum:

\begin{itemize}
  \item pseudonymous case identifier;
  \item lumbar level;
  \item central-canal stenosis severity where explicitly stated;
  \item disc bulge, protrusion and extrusion as separate morphology fields;
  \item foraminal findings and side only where explicitly stated;
  \item nerve-root pressure/impingement where reported;
  \item an explicit missing/not-reported state distinct from a negative finding.
\end{itemize}

This last distinction is essential. Absence of a phrase in a narrative report is
not automatically equivalent to ``normal''. A field may be unreported because
the radiologist considered it clinically irrelevant, because reporting practice
varied, or because the report was not structured around that compartment.
``Not reported'' is therefore encoded separately and excluded from target-
specific external evaluation unless a clinician has explicitly adjudicated the
absence as normal.

\subsection{Cohort Reconciliation}
\label{sec:method-reconciliation}

Before any local model evaluation, a case-flow table will reconcile the raw
archive, reports and analysable imaging. For each candidate case the audit will
record:

\begin{itemize}
  \item whether a report exists;
  \item whether sagittal T1 exists and is usable;
  \item whether sagittal T2/STIR exists and is usable;
  \item whether axial T2 exists and is usable;
  \item whether DICOM geometry is internally consistent;
  \item whether a central-canal reference label can be derived;
  \item whether laterality-dependent labels are available;
  \item whether the case is duplicated or represents repeat imaging;
  \item the reason for any exclusion.
\end{itemize}

The final thesis will report the flow from raw studies to the fixed external
test cohort, rather than presenting the current 299/294/raw-study counts without
explanation. Repeat examinations from one patient, if identified, remain in the
same partition to preserve patient independence.

% ===========================================================================
\section{Ethics, Governance and De-Identification}
\label{sec:method-ethics}

The Rizgary component uses retrospective clinical data. No MRI is acquired for
the purposes of this study, no patient undergoes an additional intervention, and
model output is not returned to the clinical record. Nevertheless, the data are
sensitive medical information and are subject to institutional governance.

The final thesis will state the formal approval mechanism and reference:
\toconfirm{hospital research committee/IRB approval number, date and responsible
institution}. It will also state the legal/ethical basis for retrospective use:
\toconfirm{consent waiver, broad consent or other applicable basis}. These items
must be completed before the local study is described as ethically approved.

The raw local DICOM archive contains direct identifiers and therefore cannot be
given to students or uploaded to external compute in its original state.
De-identification is performed on a controlled copy before analysis. At minimum,
the process removes direct patient identifiers, replaces case identifiers with
study pseudonyms, strips or audits private tags, and handles dates according to
the approved policy while preserving only the temporal information required for
scientific analysis. The de-identification procedure must also check for
burned-in identifiers where applicable. The re-identification key, if one must
exist for hospital governance, is stored separately from the research
repository and is not available to the model-training pipeline.

External/cloud processing of local images is prohibited unless explicitly
allowed by hospital policy:
\toconfirm{whether external cloud/GPU processing is permitted and under what
conditions}. Public benchmark data are handled according to their published
licence and competition terms.

Only aggregate results are reported. Illustrative local images may be included
in the thesis only after de-identification review and only if institutional
permission permits publication. The local cohort is not assumed to be publicly
shareable; any future data release would require separate approval and a more
stringent disclosure review.


% ===========================================================================
\section{Reference Standard and Target Encoding}
\label{sec:method-reference}

\subsection{Benchmark Target Representation}
\label{sec:method-benchmark-targets}

For LumbarDISC, the target set is represented explicitly rather than collapsed
into a generic three-class image label. Let the five lumbar levels be
\[
\mathcal{L}=\{L1\!-\!L2,L2\!-\!L3,L3\!-\!L4,L4\!-\!L5,L5\!-\!S1\},
\]
and let the five condition/laterality targets be
\[
\mathcal{C}=\{
\mathrm{Canal},
\mathrm{LF},
\mathrm{RF},
\mathrm{LSA},
\mathrm{RSA}
\},
\]
where LF/RF denote left/right neural foraminal narrowing and LSA/RSA
left/right subarticular stenosis. One examination therefore has the target set
\[
\mathcal{T}=\mathcal{L}\times\mathcal{C}, \qquad |\mathcal{T}|=25.
\]

For target $t\in\mathcal{T}$ the observed grade is
\[
y_t\in\{0,1,2\},
\]
corresponding to Normal/Mild, Moderate and Severe. Encoding the targets in this
form is not merely convenient notation. The pair $(l,c)$ is retained throughout
data loading, feature extraction, graph construction and evaluation, which
prevents a severity label from losing the anatomical meaning that generated it.

Where the benchmark provides an absent or unusable annotation, a target mask
$m_t\in\{0,1\}$ is carried alongside the grade. Masked targets do not contribute
to the supervised loss. Missing labels are never imputed from neighbouring
levels for the main experiments, because doing so would introduce exactly the
relational assumption the thesis is intended to test.

\subsection{Local Reference Matrix}
\label{sec:method-local-reference}

The Rizgary reference standard is constructed at the same anatomical level but
not forced into the same disease schema. A canonical row represents one
case--level pair. The minimum schema is:

\begin{center}
\begin{tabular}{p{2.4cm}p{2.0cm}p{8.1cm}}
\toprule
\textbf{Field} & \textbf{Type} & \textbf{Interpretation}\\
\midrule
Case ID & categorical & pseudonymous patient/study identifier\\
Level & categorical & L1--L2 through L5--S1\\
Canal stenosis & ordinal/missing & explicit report-derived severity, if stated\\
Bulge & binary/missing & circumferential/generalised bulge stated at level\\
Protrusion & binary/missing & protrusion stated at level\\
Extrusion & binary/missing & extrusion stated at level\\
Foraminal finding & categorical/missing & narrowing/pressure effect if stated\\
Laterality & left/right/bilateral/NR & side only where stated\\
Nerve-root effect & categorical/missing & pressure, displacement or compression\\
Other morphology & text/coded & migration, dehydration, height loss and related findings\\
\bottomrule
\end{tabular}
\end{center}

The local matrix keeps \emph{not reported} (NR) separate from \emph{explicitly
absent}. This rule is especially important for foraminal and subarticular
findings, because local reporting completeness differs from the structured RSNA
annotation scheme. A target is eligible for external evaluation only when the
reference standard provides a defensible class label.

\subsection{Report Extraction and Human Verification}
\label{sec:method-report-verification}

Automated report extraction, if available from the parallel clinical-NLP work,
may accelerate construction of the matrix but does not define the gold
standard. The final labels used for Selar's external evaluation are manually
verified against the original report text.

A two-stage verification procedure is preferred:

\begin{enumerate}
  \item \textbf{Structured extraction audit.} Every non-normal/non-empty finding
  is traced back to the exact report phrase, level and laterality. Multi-level
  expressions such as ``L2--3 and L4--5, L5--S1 shows circumferential disc
  bulge'' are expanded into separate level records without changing their
  meaning.

  \item \textbf{Clinical reliability subset.} A pre-specified stratified subset
  of local studies is independently re-read by two qualified readers using an
  explicit grading sheet. Disagreements are adjudicated by a more senior reader
  where available. Inter-reader agreement is reported with weighted
  $\kappa$ for ordinal grades and ordinary agreement/appropriate $\kappa$ for
  binary morphology. The size of the re-read subset will be fixed before model
  results are examined, based on reader availability and the precision required
  for the agreement estimate.
\end{enumerate}

If a second-reader study cannot be completed, the report-derived matrix remains
usable as a real-world reference standard, but the absence of independent
regrading is stated explicitly as a limitation and no claim is made that local
labels reproduce the RSNA grading protocol exactly.

\subsection{Schema Alignment Between LumbarDISC and Rizgary}
\label{sec:method-schema-alignment}

Cross-institutional evaluation requires equality of \emph{task}, not merely
similar words. The mapping is therefore conservative.

\begin{center}
\begin{tabular}{p{4.1cm}p{4.0cm}p{5.0cm}}
\toprule
\textbf{LumbarDISC target} & \textbf{Rizgary evidence} & \textbf{Transfer status}\\
\midrule
Central canal stenosis & level-resolved mild/moderate/severe statements where present
& Primary external target\\
Left/right foraminal narrowing & pressure/narrowing often incompletely lateralised
& Exclude unless laterality and grade are adequate or freshly re-read\\
Left/right subarticular stenosis & not routinely represented in current reports
& Exclude unless new expert annotation is created\\
Disc bulge/protrusion/extrusion & directly described locally but not an RSNA stenosis target
& Separate local morphology task\\
\bottomrule
\end{tabular}
\end{center}

No local disc-herniation morphology is converted into a stenosis grade by
heuristic rule. Likewise, ``pressure on exit foramina'' is not automatically
treated as a left/right foraminal severity label unless the report or expert
review supplies the information required by that target. This prevents
apparent external validation from being created by label re-interpretation.

% ===========================================================================
\section{Data Partitioning and Leakage Control}
\label{sec:method-splits}

\subsection{Patient-Level Independence}
\label{sec:method-patient-split}

All splits are made at patient/study level before slice extraction or data
augmentation. For a patient $p$, every sagittal series, every axial series,
every disc level, every side and every augmented crop belongs to one and only
one partition. A programmatic assertion is run before training to confirm that
the intersection of patient identifiers across partitions is empty.

The split record is version-controlled as a list of pseudonymous IDs. Data
loaders consume those fixed lists rather than performing a new random split each
time the training script is executed.

\subsection{Public Development, Validation and Held-Out Evaluation}
\label{sec:method-public-split}

Where the available public dataset exposes institution identifiers, a
site-aware evaluation is preferred because it tests a more meaningful shift
than a random split. The exact public partitioning scheme will be chosen
according to the metadata available in the held dataset and will be frozen
before model comparison. At minimum there will be:

\begin{itemize}
  \item a development/training partition;
  \item a validation partition for hyperparameter choice, early stopping and
        calibration;
  \item a held-out public evaluation partition not used for model selection.
\end{itemize}

If site identifiers permit leave-one-site-out or institution-held-out analysis,
that analysis is reported in addition to the ordinary benchmark split. It is not
a replacement for Rizgary, because the local hospital introduces a different
reporting/reference-standard domain as well as a different imaging domain.

\subsection{Fixed Local Test Set and Adaptation Pool}
\label{sec:method-local-split}

The local transfer study is protected against adaptation leakage by separating
Rizgary into two parts \emph{before} the public model is evaluated:

\begin{enumerate}
  \item a \textbf{fixed local test set}, which is never used for adaptation,
        calibration or threshold selection; and
  \item a \textbf{local adaptation pool}, from which the $N=10,25,50,100$
        few-shot subsets are sampled.
\end{enumerate}

The exact number allocated to each part will be determined after the audited
class distribution is known. The split will be stratified by central-canal
severity and lumbar level where feasible without producing clinically
implausible or extremely small strata. The fixed local test set is used for the
zero-shot baseline and for every subsequent adapted model, making the recovery
curves directly comparable.

For each adaptation size $N$, several independent stratified subsets are sampled
from the adaptation pool. Results are averaged across these draws with variance
reported, because one set of ten cases can be unrepresentative by chance. The
test set itself never changes across adaptation sizes.

\subsection{Prevention of Information Leakage During Self-Supervision}
\label{sec:method-ssl-leakage}

Self-supervised learning does not remove the need for split discipline. A model
pretrained on the held-out test patients, even without labels, has already seen
their anatomy and acquisition characteristics. Therefore anatomical
self-supervision uses only the development partition during internal
experiments. The held-out public test set and fixed Rizgary test set remain
unseen in both supervised and self-supervised training.

The same restriction applies to intensity normalisation parameters and learned
harmonisation models: any transformation whose parameters are estimated from
data is fitted on the training/development partition and then applied unchanged
to validation/test data.

% ===========================================================================
\section{DICOM Reconstruction and Quality Control}
\label{sec:method-dicom}

\subsection{Series Identification}
\label{sec:method-series-identification}

Each study is first organised by DICOM \texttt{SeriesInstanceUID}. Candidate
series are classified into sagittal T1, sagittal T2/STIR and axial T2 using a
combination of DICOM metadata and image-orientation checks. Series descriptions
are informative but not trusted alone because naming conventions differ across
institutions. The classifier therefore considers fields such as
\texttt{SeriesDescription}, \texttt{ProtocolName}, sequence parameters where
available, and the orientation derived from \texttt{ImageOrientationPatient}.

Ambiguous or duplicated candidate series are flagged for manual review rather
than selected by filename order. The selected series identifiers are stored in
a manifest so that every experiment uses the same source series.

\subsection{Patient-Space Geometry}
\label{sec:method-geometry}

For each DICOM slice, the image plane is reconstructed from
\texttt{ImagePositionPatient}, \texttt{ImageOrientationPatient} and
\texttt{PixelSpacing}. Let $\mathbf{s}$ be the patient-space location of the
upper-left pixel, $\mathbf{r}$ and $\mathbf{c}$ the row and column direction
cosines, and $\Delta_r,\Delta_c$ the corresponding pixel spacings. The physical
coordinate of pixel $(i,j)$ is

\begin{equation}
\mathbf{x}(i,j)=
\mathbf{s}
+i\,\Delta_r\,\mathbf{r}
+j\,\Delta_c\,\mathbf{c}.
\label{eq:dicom-patient-coordinate}
\end{equation}

The slice normal is

\begin{equation}
\mathbf{n}=\mathbf{r}\times\mathbf{c}.
\label{eq:slice-normal}
\end{equation}

Slices are ordered by projection onto $\mathbf{n}$ rather than by filename.
Spacing between consecutive slice origins is inspected for discontinuities.
This procedure preserves patient coordinates across sagittal and axial
acquisitions, allowing a level localised in one plane to be projected into the
other.

For a sagittal target point $\mathbf{q}$ and candidate axial slice $k$ with
origin $\mathbf{s}_k$ and normal $\mathbf{n}_k$, the perpendicular distance to
that slice plane is

\begin{equation}
d_k =
\left|
\mathbf{n}_k^\top(\mathbf{q}-\mathbf{s}_k)
\right|.
\label{eq:plane-distance}
\end{equation}

The nearest anatomically corresponding axial slices are the slices with the
smallest $d_k$, subject to coverage and orientation checks. The method is an
implementation of physical correspondence, not a learned registration shortcut,
and follows the geometric principle used by prior lumbar pipelines
\cite{pan2021automatically,batra2025mscan}.

\subsection{Orientation Standardisation}
\label{sec:method-orientation}

Volumes are reoriented to a common anatomical convention for tensor processing
while retaining the transform to original patient coordinates. Left/right must
never be inferred after arbitrary array transposition; all laterality-dependent
targets are checked against the DICOM orientation transform. This is particularly
important for foraminal and subarticular labels, where a left/right inversion
would create a scientifically silent but catastrophic error.

\subsection{Voxel Spacing and Resampling}
\label{sec:method-resampling}

The acquisition spacing varies by institution, scanner and plane. Inputs are
therefore resampled only after anatomical coordinates have been reconstructed.
The target spacing is chosen from the training-set distribution so that it
preserves structures relevant to grading without creating unnecessary memory
cost. Sagittal and axial streams may use different in-plane resolutions because
their native geometry and diagnostic roles differ.

Interpolation uses a continuous method (e.g.\ linear or higher-order interpolation)
for MR intensities and nearest-neighbour interpolation for discrete masks or
labels. Original spacing is retained in metadata for any measurement reported
in millimetres. The exact target spacing will be \torecord{final sagittal and
axial resampling spacings selected from the development cohort}.

\subsection{Intensity Processing}
\label{sec:method-intensity}

MRI intensities are not quantitatively standardised across scanners. The
preprocessing therefore removes gross scale differences without imposing a
global histogram learned from test data. A default pipeline is:

\begin{enumerate}
  \item suppress non-finite values and obvious padding/background values;
  \item define a foreground mask from the body/spine region;
  \item clip extreme foreground intensities using training-derived robust
        percentiles;
  \item perform per-volume or per-ROI robust z-score normalisation;
  \item optionally compare with bias-field correction or histogram
        standardisation as a robustness ablation.
\end{enumerate}

Any learned harmonisation is fitted only on the development data. The external
Rizgary zero-shot result is reported both with the baseline preprocessing and,
if a separate harmonisation experiment is conducted, with that method clearly
labelled as an intervention rather than folded into the zero-shot definition.

\subsection{Quality-Control Flags}
\label{sec:method-qc}

Automated QC records the following per series:

\begin{itemize}
  \item slice count and ordering consistency;
  \item missing or duplicated instance positions;
  \item orientation consistency within series;
  \item in-plane spacing and slice spacing;
  \item field of view and lumbar coverage;
  \item intensity dynamic range;
  \item severe motion or truncation where detectable;
  \item presence of all sequences required by the particular experiment.
\end{itemize}

Cases failing geometry checks are not silently forced into the pipeline. They
are reviewed, excluded with a recorded reason, or assigned to a robustness
analysis if the failure itself represents a realistic missing-data condition.
A final inclusion/exclusion table will report every such decision.

% ===========================================================================
\section{Anatomical Localisation and ROI Construction}
\label{sec:method-localisation}

\subsection{Role of Localisation}
\label{sec:method-localisation-role}

Localisation exists to make the subsequent comparison fair. The thesis does not
claim that identifying L1--L2 through L5--S1 is itself novel; Chapter~2 showed
that this problem is mature and that strong systems already attain highly
reliable level assignment \cite{pan2021automatically}. The requirement here is
that the localiser be accurate enough that representation and graph experiments
are not dominated by wrong crops.

An established model is therefore preferred. Candidate implementations include
a SPIDER-pretrained nnU-Net-style segmentation network
\cite{isensee2021nnunet,vandergraaf2024lumbar}, a heatmap-regression detector,
or another state-of-the-art anatomical parser available when implementation
begins. The choice is made on validation localisation error and robustness, not
on whether it is architecturally fashionable.

\subsection{Outputs}
\label{sec:method-localisation-output}

The localiser produces, for each study:

\begin{itemize}
  \item centres of L1--L2 through L5--S1 in patient coordinates;
  \item vertebral/disc/canal masks or landmarks where the selected model provides
        them;
  \item a confidence or quality flag for each detected level;
  \item the transform between voxel coordinates and DICOM patient space.
\end{itemize}

A failed or low-confidence localisation is retained as a distinct failure mode.
It is not corrected manually for the primary automated pipeline without being
reported, because doing so would inflate downstream grading performance by
removing a real deployment error.

\subsection{Localisation Evaluation}
\label{sec:method-localisation-eval}

Where point annotations are available, centre localisation error is measured in
millimetres:

\begin{equation}
e_l = \|\hat{\mathbf{q}}_l-\mathbf{q}_l\|_2.
\label{eq:localisation-error}
\end{equation}

Mean/median error and percentile distributions are reported by level.
Percentage-of-correct-keypoint (PCK) or an equivalent tolerance-based measure is
also reported at pre-specified physical thresholds. Where full masks are
available, Dice and surface-distance metrics are reported. These results are
separated from disease grading, consistent with the literature showing that
detection/localisation and fine-grained grading generalise differently.

\subsection{Disease-Specific 2.5D Regions of Interest}
\label{sec:method-roi}

One rectangular crop is not assumed to be equally appropriate for every target.
The crop definition is conditioned on anatomical compartment.

For central canal stenosis, the model receives sagittal T2/STIR context around
the disc level together with corresponding axial T2 slices. For neural
foraminal narrowing, the sagittal T1 stream receives greater parasagittal
coverage and the ROI is side-aware, with sagittal T2 retained as complementary
context. For subarticular stenosis, axial T2 supplies the dominant local
evidence with sagittal context. These choices follow the clinical sequence--
pathology correspondence established in Chapter~2 rather than being learned
from arbitrary crop templates.

Each 2.5D stream contains a small ordered stack around the anatomically closest
slice,

\[
\mathcal{I}_l =
\{I_{k-r},\ldots,I_k,\ldots,I_{k+r}\},
\]

where $k$ is the level-centred slice and $r$ is a small radius determined on the
development set. A five-slice stack ($r=2$) is the initial reference
configuration because it captures through-plane context at modest memory cost;
the final radius is reported after sensitivity testing rather than treated as
a novel contribution.

Crops are defined in physical dimensions where possible, then resampled to the
network input resolution. This avoids a fixed 100-pixel crop representing
different anatomical widths on scanners with different pixel spacing.

\subsection{ROI Quality Control}
\label{sec:method-roi-qc}

A validation subset is visually inspected using overlays of the detected level,
crop boundary and corresponding axial slices. The inspection records:

\begin{itemize}
  \item correct lumbar level;
  \item inclusion of the relevant canal/foramen/recess;
  \item correct left/right orientation;
  \item adequate axial coverage;
  \item crop truncation;
  \item correspondence failure caused by unusual axial angulation.
\end{itemize}

The purpose is not to manually curate the full dataset but to detect systematic
geometry errors before they propagate into model training.

% ===========================================================================
\section{Baseline Representation and Sequence Encoders}
\label{sec:method-baseline}

\subsection{Independent ROI Classifier}
\label{sec:method-e0}

The mandatory baseline (E0) is an independent ROI classifier. Each target is
graded from its anatomically localised input without communication with other
targets. This represents the prevailing formulation criticised in
Chapter~\ref{ch:litreview} and provides the reference against which relational
reasoning is measured.

The baseline uses a well-characterised image encoder such as ResNet
\cite{he2016deep}, EfficientNet \cite{tan2019efficientnet}, Swin
\cite{liu2021swin}, or another contemporary backbone selected before final
experiments. The purpose is not to prove that one family is universally
superior. A primary backbone is selected for the ablation ladder so that
representation, routing and graph changes are measured under one stable
feature extractor. A second backbone may be used in a model-agnosticity check
for the strongest contribution.

\subsection{Sequence-Specific Feature Extraction}
\label{sec:method-sequence-encoders}

Each available MRI sequence is processed by a sequence-specific encoder
$E_m(\cdot)$,

\begin{equation}
\mathbf{f}_{p,l,m}=E_m(\mathcal{I}_{p,l,m}),
\label{eq:sequence-encoder}
\end{equation}

for patient $p$, level $l$ and modality/sequence $m$. The encoders may share
architecture while keeping separate normalisation statistics and weights, or
may partially share weights if an explicit ablation shows that doing so is
beneficial. Separate encoders are the default because T1, sagittal T2/STIR and
axial T2 have different contrast and orientation statistics.

To keep comparisons fair, the dimensionality of the output representation is
held constant across fusion methods. Fixed concatenation, simple averaging and
ordinary cross-attention all consume the same per-sequence features.

\subsection{Controlled Backbone Selection}
\label{sec:method-backbone-control}

Backbone choice is selected on the validation partition before the central
novelty experiments and then held fixed. The candidate comparison will use
identical crops, augmentations, optimiser budget and patient split. If one
backbone is selected because it gives the most stable validation performance
rather than the absolute highest point estimate, that decision is recorded.

This design prevents the thesis from becoming a repeated architecture search.
The main questions concern anatomical self-supervision, adaptive modality
allocation and target-level relational reasoning. A backbone is a controlled
feature extractor in those experiments.


% ===========================================================================
\section{Core Contribution I: Anatomically Aligned Cross-Sequence Self-Supervision}
\label{sec:method-ssl}

\subsection{Rationale}
\label{sec:method-ssl-rationale}

Conventional self-supervised vision learning defines related views by applying
augmentations to one image. That is useful but does not exploit a distinctive
property of multi-sequence MRI: sagittal T1, sagittal T2/STIR and axial T2 can
look very different while representing the same underlying patient and spinal
level. DICOM geometry supplies the correspondence directly.

The proposed objective therefore treats anatomical identity, rather than visual
similarity, as the primary source of positive pairs. For patient $p$, level
$l$, and two sequences $m_a$ and $m_b$, embeddings
$\mathbf{z}_{p,l,m_a}$ and $\mathbf{z}_{p,l,m_b}$ are encouraged to encode a
shared anatomical representation even though acquisition plane and contrast may
differ.

This idea is evaluated as a representation-learning hypothesis, not assumed to
be beneficial. The literature already shows that generic and longitudinal
self-supervision can reduce label dependence
\cite{jamaludin2017selfsupervised,chen2020simple}; the experiment here asks
whether explicitly anatomical cross-sequence correspondence is a stronger
signal for lumbar grading.

\subsection{Projection Head}
\label{sec:method-ssl-projection}

Each sequence feature is mapped through a projection head $P_m$:

\begin{equation}
\mathbf{z}_{p,l,m}=
\frac{
P_m(\mathbf{f}_{p,l,m})
}{
\|P_m(\mathbf{f}_{p,l,m})\|_2
}.
\label{eq:ssl-projection}
\end{equation}

The projection space is used for the self-supervised objective; the encoder
features before projection are retained for downstream grading. This prevents
the representation used for contrastive geometry from constraining the final
classifier unnecessarily.

\subsection{Positive-Pair Construction}
\label{sec:method-positive-pairs}

The strongest positive relation is:

\[
(p,l,m_a)\sim(p,l,m_b),
\qquad m_a\neq m_b,
\]

meaning the same patient and same anatomical level in two different sequences.
Correspondence is established through the DICOM reconstruction and level
localiser, not through nearest image appearance.

An optional softer relation may be tested for adjacent levels from the same
patient, because neighbouring levels share biomechanical context. It is not
included in the primary objective by default. Treating L3--L4 and L4--L5 as
equivalent positives would risk erasing genuine level-specific pathology, so
any adjacent-level term is isolated as a separate ablation.

Different patients at the same level are not defined as positives in the main
objective, because common anatomical level does not imply common disease state.
They may be used in a secondary supervised-contrastive variant once labels are
available, but this is conceptually distinct from the proposed anatomical
self-supervision.

\subsection{Contrastive Objective}
\label{sec:method-ssl-objective}

For anchor embedding $\mathbf{z}_i$ and its cross-sequence positive
$\mathbf{z}_{i^+}$, the reference InfoNCE-style loss is

\begin{equation}
\mathcal{L}_{\mathrm{ACSSL}}(i)
=
-\log
\frac{
\exp(\mathrm{sim}(\mathbf{z}_i,\mathbf{z}_{i^+})/\tau)
}{
\sum_{j\in\mathcal{B}\setminus\{i\}}
\exp(\mathrm{sim}(\mathbf{z}_i,\mathbf{z}_j)/\tau)
},
\label{eq:acssl}
\end{equation}

where $\mathrm{sim}$ is cosine similarity, $\tau$ is a temperature parameter
and $\mathcal{B}$ is the contrastive batch. When more than one anatomically
valid positive exists, the numerator is extended to the set of positives rather
than choosing one arbitrarily.

False negatives are a concern because two different patients may have very
similar normal anatomy. The implementation will therefore compare ordinary
InfoNCE with a variant that does not force all non-positive samples to be hard
negatives, such as a multi-positive or non-contrastive consistency objective.
The novelty claim is tied to the \emph{anatomical pairing}, not to a particular
brand of contrastive loss.

\subsection{Label-Efficiency Experiment}
\label{sec:method-label-efficiency}

The primary test of RQ2 uses fixed fractions of the labelled development set:
10\%, 25\%, 50\% and 100\%. For each fraction:

\begin{enumerate}
  \item the encoder is initialised from the candidate pretraining strategy;
  \item the same supervised grading head is fitted;
  \item patient splits and target distribution are held fixed;
  \item several random labelled subsets are evaluated where computationally
        feasible;
  \item macro-F1, weighted $\kappa$, severe-class recall and calibration are
        compared.
\end{enumerate}

The competing initialisations are:

\begin{itemize}
  \item training from random initialisation;
  \item ImageNet initialisation;
  \item generic medical/volumetric pretraining where a suitable public model is
        available;
  \item ordinary augmentation-based self-supervision;
  \item the proposed anatomically aligned cross-sequence self-supervision.
\end{itemize}

A successful result is not defined as beating every baseline at every label
fraction. The useful scientific question is where anatomical pairing helps,
where it does not, and whether any benefit is largest under label scarcity or
domain shift.

\subsection{Transfer Test of the Representation}
\label{sec:method-ssl-transfer}

The public-data model used for zero-shot Rizgary evaluation is trained once with
and once without anatomical self-supervision while the downstream architecture
is held constant. If the anatomically pretrained model loses less performance
externally despite comparable internal performance, that would support the
claim that the representation is less tied to dataset-specific appearance. If
both models degrade equally, the self-supervised contribution is limited to
label efficiency rather than domain robustness.

% ===========================================================================
\section{Core Contribution II: Disease-Conditioned Adaptive Sequence Routing}
\label{sec:method-routing}

\subsection{Problem Formulation}
\label{sec:method-routing-problem}

A conventional multi-sequence model applies the same fusion rule to every
target. That assumption is clinically implausible because the sequences are not
equally informative for each anatomical compartment. Sagittal T1 is particularly
useful for foraminal fat obliteration, sagittal T2/STIR for disc/canal context,
and axial T2 for canal and lateral recess morphology. The router turns that
known heterogeneity into a learnable, testable mechanism.

For patient $p$, target $t=(l,c)$ and available modalities
$\mathcal{M}_p$, let $\mathbf{f}_{p,l,m}$ be the feature from modality $m$.
The fused target feature is

\begin{equation}
\mathbf{u}_{p,t}
=
\sum_{m\in\mathcal{M}_p}
g_{p,t,m}\,
\mathbf{f}_{p,l,m},
\label{eq:routed-feature}
\end{equation}

where the non-negative routing weights satisfy

\begin{equation}
\sum_{m\in\mathcal{M}_p} g_{p,t,m}=1.
\label{eq:routing-simplex}
\end{equation}

The gate is conditioned on target identity and patient features:

\begin{equation}
g_{p,t,m}
=
\mathrm{softmax}_{m\in\mathcal{M}_p}
\left(
G[
\mathbf{f}_{p,l,m},
\mathbf{e}_{c},
\mathbf{e}_{l},
\mathbf{a}_{p}
]
\right),
\label{eq:routing-gate}
\end{equation}

where $\mathbf{e}_{c}$ and $\mathbf{e}_{l}$ are learnable condition and level
embeddings and $\mathbf{a}_{p}$ denotes optional acquisition/context metadata
that is available without leaking the target label. The initial implementation
uses a lightweight MLP or target-conditioned attention gate. A more complex
mixture-of-experts router is considered only if the simpler mechanism is
insufficient.

\subsection{Explicit Availability Mask}
\label{sec:method-routing-mask}

Missing sequences are represented by an availability mask
$a_{p,m}\in\{0,1\}$. Logits for unavailable modalities are masked before the
softmax:

\begin{equation}
g_{p,t,m}=0 \quad \mathrm{if}\quad a_{p,m}=0.
\label{eq:routing-missing-mask}
\end{equation}

The remaining weights are renormalised over the modalities that actually
exist. A missing sequence is therefore not represented by an all-zero image
whose meaning the network must infer. Availability is explicit in the model.

\subsection{Modality Dropout}
\label{sec:method-modality-dropout}

During training, complete studies are stochastically converted into incomplete
ones by removing one or more sequence streams. The dropout distribution includes:

\begin{itemize}
  \item all three sequences;
  \item sagittal T1 + sagittal T2;
  \item sagittal T2 + axial T2;
  \item sagittal T1 + axial T2 where clinically meaningful;
  \item each single-sequence configuration as an edge case.
\end{itemize}

The sampling probabilities are tuned so that the model sees enough complete
studies to learn complementary fusion while also receiving substantial exposure
to incomplete inputs. The exact probabilities are \torecord{final modality-
dropout schedule selected on validation data}.

No external test case is artificially repaired by copying another sequence into
a missing slot. The model's purpose is to remain functional under genuine
absence, not to conceal it.

\subsection{Routing Consistency}
\label{sec:method-routing-consistency}

An optional consistency term encourages the prediction from an incomplete
modality set to remain close to the complete-study prediction when the removed
modality is non-essential:

\begin{equation}
\mathcal{L}_{\mathrm{miss}}
=
D\!\left(
p_{\theta}(y\mid x_{\mathrm{full}}),
p_{\theta}(y\mid x_{\mathrm{subset}})
\right),
\label{eq:missing-consistency}
\end{equation}

where $D$ is a divergence such as KL divergence or Jensen--Shannon divergence.
This term is not applied blindly: for targets known to depend strongly on a
specific sequence, forcing invariance to its removal could suppress legitimate
diagnostic information. The term is therefore ablated by target and compared
with modality-dropout training alone.

\subsection{Routing Baselines}
\label{sec:method-routing-baselines}

RQ3 is evaluated against four progressively stronger controls:

\begin{enumerate}
  \item fixed feature concatenation;
  \item simple equal-weight averaging;
  \item ordinary cross-attention over available sequence features;
  \item disease-conditioned routing without modality dropout;
  \item disease-conditioned routing with modality dropout.
\end{enumerate}

For each method, encoder capacity, ROI inputs and training budget are matched as
closely as possible. Missing-sequence robustness is evaluated on the same
patients under the same artificial absence pattern.

\subsection{Interpreting Routing Weights}
\label{sec:method-routing-interpretation}

The gate weights are model allocations, not causal estimates of sequence
importance. A high axial T2 weight does not prove that axial T2 caused a correct
decision. Clinical plausibility is assessed in two complementary ways.

First, aggregate routing weights are summarised by target. The expected pattern
is that foraminal targets allocate more weight to sagittal T1 and subarticular/
canal targets more to axial T2, but this expectation is not built into the gate
as a hard prior.

Second, controlled input ablation tests whether removing the sequence with high
routing weight causes a greater performance loss than removing a low-weight
sequence. Agreement between learned allocation and intervention strengthens the
interpretation; disagreement is reported rather than rationalised post hoc.

% ===========================================================================
\section{Core Contribution III: Heterogeneous Disease--Anatomy Graph}
\label{sec:method-graph}

\subsection{Graph Construction}
\label{sec:method-graph-construction}

The proposed graph represents \emph{prediction targets}, not pixels and not a
3D disc surface. At full LumbarDISC scope, the node set is

\begin{equation}
V=\{v_{l,c}:l\in\mathcal{L},\,c\in\mathcal{C}\},
\qquad |V|=25.
\label{eq:graph-nodes}
\end{equation}

Each node begins with the routed visual representation for that target together
with embeddings that identify level and condition:

\begin{equation}
\mathbf{h}^{(0)}_{l,c}
=
[
\mathbf{u}_{l,c}
\Vert
\mathbf{e}_{l}
\Vert
\mathbf{e}_{c}
].
\label{eq:node-initial}
\end{equation}

The graph topology is defined anatomically before training. It is not learned
from correlations in the test data.

\subsection{Typed Edge Families}
\label{sec:method-edge-types}

Three primary edge types are used.

\textbf{Adjacent-level edges} connect the same target type across neighbouring
lumbar levels, for example

\[
v_{L3-L4,\mathrm{Canal}}
\leftrightarrow
v_{L4-L5,\mathrm{Canal}}.
\]

These encode the ordered chain of lumbar motion segments.

\textbf{Same-level cross-condition edges} connect anatomical compartments that
coexist at one level, for example central canal to left/right foraminal and
subarticular targets. The edge expresses that the targets share local
morphology, not that one grade deterministically implies another.

\textbf{Bilateral edges} connect left and right counterparts of the same
condition, for example

\[
v_{L4-L5,\mathrm{LF}}
\leftrightarrow
v_{L4-L5,\mathrm{RF}}.
\]

This represents bilateral anatomical symmetry without forcing equal grades.

Additional cross-type edges are added only if clinically justified in the
pre-specified topology. A denser graph is not automatically a better graph; the
study tests whether explicit relation types add value over homogeneous message
passing.

\subsection{Relation-Aware Message Passing}
\label{sec:method-message-passing}

For layer $q$, node $i$ receives messages from neighbours $j\in\mathcal{N}(i)$.
A relation-aware attention formulation is

\begin{align}
\mathbf{q}_i &= W_Q\mathbf{h}^{(q)}_i,\\
\mathbf{k}_{j,r} &= W^r_K\mathbf{h}^{(q)}_j,\\
\mathbf{v}_{j,r} &= W^r_V\mathbf{h}^{(q)}_j,
\end{align}

with attention

\begin{equation}
\alpha_{ijr}
=
\mathrm{softmax}_{j,r}
\left(
\frac{
\mathbf{q}_i^\top\mathbf{k}_{j,r}
}{
\sqrt{d}
}
+b_r
\right),
\label{eq:relational-attention}
\end{equation}

and update

\begin{equation}
\tilde{\mathbf{h}}^{(q+1)}_i
=
\sum_{(j,r)\in\mathcal{N}(i)}
\alpha_{ijr}\mathbf{v}_{j,r}.
\label{eq:graph-message}
\end{equation}

Residual connection, normalisation and feed-forward transformation produce the
final layer output. Relation-specific transformations $W^r_K,W^r_V$ allow an
adjacent-level neighbour to influence a node differently from its contralateral
counterpart.

The exact graph transformer implementation may be replaced by a computationally
simpler relation-aware GAT if validation shows no benefit from the more complex
version. The contribution is the explicit typed target graph and its controlled
evaluation, not a requirement to use the largest graph architecture.

\subsection{Graph Baselines}
\label{sec:method-graph-baselines}

The graph contribution is tested against:

\begin{enumerate}
  \item \textbf{Independent heads:} no message passing between targets.
  \item \textbf{Ordered level transformer:} features communicate across the five
        lumbar levels but do not form 25 typed condition/laterality nodes.
  \item \textbf{Homogeneous graph:} the same 25 nodes and adjacency pattern but
        one undifferentiated edge type.
  \item \textbf{Heterogeneous graph:} typed adjacent-level, cross-condition and
        bilateral relations.
  \item \textbf{Random/shuffled-edge control:} graph capacity retained while
        anatomical edges are permuted.
\end{enumerate}

The shuffled-edge control is essential. If a graph with arbitrary topology
performs as well as the anatomical graph, the study has shown that extra
message-passing capacity helps, not that the encoded anatomy matters.

\subsection{Edge-Family Ablation}
\label{sec:method-edge-ablation}

The full graph is further decomposed:

\begin{itemize}
  \item remove adjacent-level edges;
  \item remove same-level cross-condition edges;
  \item remove bilateral edges;
  \item retain only one family at a time.
\end{itemize}

Results are reported by target family. It is plausible, for example, that
bilateral edges help lateral-compartment grading but add little to central-canal
classification. A single average graph improvement would conceal this structure.

\subsection{Graph Consistency and Uncertainty}
\label{sec:method-graph-uncertainty}

Graph information is not used to hard-code clinically deterministic rules.
Severe stenosis at L4--L5 does not require moderate stenosis at L3--L4, and the
model is allowed to preserve isolated pathology. Relational message passing is
therefore soft.

For uncertainty analysis, node disagreement can be measured between an
independent-head prediction and the graph-refined prediction, or between
neighbour-conditioned and local evidence. This disagreement is explored as an
auxiliary signal for selective prediction, but it is not equated with calibrated
clinical uncertainty unless validated against error.

% ===========================================================================
\section{Supporting Method: Ordinal and Cost-Sensitive Grading}
\label{sec:method-ordinal}

\subsection{Categorical Baseline}
\label{sec:method-ce}

Standard three-class cross-entropy is retained as the mandatory reference:

\begin{equation}
\mathcal{L}_{\mathrm{CE}}
=
-\sum_{t}m_t
\sum_{k=0}^{2}
\mathbb{I}(y_t=k)\log p_{t,k}.
\label{eq:cross-entropy}
\end{equation}

Class weighting or balanced sampling is applied in separate controlled variants.
This baseline is necessary because previous lumbar work found that explicitly
ordinal objectives did not consistently outperform cross-entropy
\cite{niemeyer2021a}.

\subsection{Ordinal Threshold Formulation}
\label{sec:method-ordinal-threshold}

An ordinal head represents a three-grade target as two ordered threshold
questions:

\[
P(y>0\mid\mathbf{h}), \qquad
P(y>1\mid\mathbf{h}).
\]

A CORAL/CORN-style implementation may be used provided the ordering constraint
is respected. The exact ordinal formulation is selected before the final
ablation and compared directly with categorical cross-entropy under identical
features.

Performance is evaluated not only by exact class agreement but also by error
distance:

\begin{equation}
d_t=|\hat{y}_t-y_t|.
\label{eq:grade-distance}
\end{equation}

The proportions with $d_t=0$, $d_t=1$ and $d_t=2$ are reported separately.

\subsection{Clinically Asymmetric Cost}
\label{sec:method-cost}

A cost matrix $C$ allows distant under-grading to receive a larger training
penalty:

\begin{equation}
C=
\begin{bmatrix}
0 & c_{01} & c_{02}\\
c_{10} & 0 & c_{12}\\
c_{20} & c_{21} & 0
\end{bmatrix},
\label{eq:cost-matrix}
\end{equation}

with the clinically motivated condition

\[
c_{20}>c_{21},
\]

so that Severe$\rightarrow$Normal/Mild is more costly than
Severe$\rightarrow$Moderate. The exact numerical values are not invented by the
model developer. They are either aligned with a defensible benchmark/clinical
weighting or selected through a sensitivity analysis across a small
pre-specified family of matrices.

One implementation uses the expected cost under the predicted distribution:

\begin{equation}
\mathcal{L}_{\mathrm{cost}}
=
\sum_t m_t
\sum_{k=0}^{2}
C_{y_t,k}\,p_{t,k}.
\label{eq:expected-cost}
\end{equation}

The cost term is evaluated alongside, not instead of, standard calibration and
discrimination metrics. A model is not considered better merely because it
optimises the chosen matrix.

\subsection{Class Imbalance}
\label{sec:method-imbalance}

Because severe grades are uncommon in LumbarDISC, the study compares a small
set of established strategies rather than stacking all of them simultaneously:

\begin{itemize}
  \item unweighted cross-entropy;
  \item inverse-frequency or effective-number class weighting;
  \item weighted sampling;
  \item focal loss \cite{lin2017focal};
  \item the cost-sensitive/ordinal variant above.
\end{itemize}

MixUp/CutMix are not the default for ordinal grading because mixed labels may
have no clinical interpretation. Geometric augmentation is constrained so that
left/right labels are updated consistently; horizontal flipping is prohibited
for laterality-dependent targets unless labels and orientation are transformed
together.

% ===========================================================================
\section{Supporting Method: Calibration, Uncertainty and Selective Prediction}
\label{sec:method-calibration}

\subsection{Probability Calibration}
\label{sec:method-temperature}

Raw neural-network probabilities are not treated as calibrated confidence.
Temperature scaling is fitted on the validation set:

\begin{equation}
p_k^{(T)}
=
\frac{
\exp(z_k/T)
}{
\sum_j \exp(z_j/T)
},
\label{eq:temperature}
\end{equation}

where $T>0$ is selected without access to the test labels. Calibration is then
evaluated with Expected Calibration Error, Brier score and reliability diagrams.

\subsection{Predictive Uncertainty}
\label{sec:method-uncertainty}

The primary feasible uncertainty methods are MC dropout and/or a small deep
ensemble. If $S$ stochastic predictions are available, the predictive mean is

\begin{equation}
\bar{\mathbf{p}}
=
\frac{1}{S}\sum_{s=1}^{S}\mathbf{p}^{(s)},
\label{eq:predictive-mean}
\end{equation}

and uncertainty can be summarised by predictive entropy,

\begin{equation}
H(\bar{\mathbf{p}})
=
-\sum_k \bar{p}_k\log\bar{p}_k,
\label{eq:entropy}
\end{equation}

together with between-member variance. The study reports whether uncertainty is
higher on incorrect predictions, distant grade errors, low-agreement targets
and external-domain cases.

\subsection{Selective Prediction}
\label{sec:method-selective}

A selective system abstains when uncertainty exceeds threshold $\delta$:

\[
U(x)>\delta \Rightarrow \mathrm{refer/abstain}.
\]

Rather than choosing one arbitrary threshold, the complete risk--coverage curve
is reported: as increasingly uncertain cases are withheld, what fraction of the
cohort remains covered and what error remains among the retained predictions?
The area under, or summary points from, this curve can compare methods.

The terminology ``refer'' describes model behaviour only. It does not claim a
validated clinical workflow; a prospective reader study would be required to
show that abstention improves patient care.

% ===========================================================================
\section{Composite Training Objective}
\label{sec:method-total-loss}

The complete model may use the objective

\begin{equation}
\mathcal{L}_{\mathrm{total}}
=
\mathcal{L}_{\mathrm{grade}}
+\lambda_{\mathrm{ssl}}\mathcal{L}_{\mathrm{ACSSL}}
+\lambda_{\mathrm{miss}}\mathcal{L}_{\mathrm{miss}}
+\lambda_{\mathrm{cost}}\mathcal{L}_{\mathrm{cost}}
+\lambda_{\mathrm{reg}}\mathcal{L}_{\mathrm{reg}},
\label{eq:total-loss}
\end{equation}

where $\mathcal{L}_{\mathrm{grade}}$ is categorical or ordinal supervised loss,
$\mathcal{L}_{\mathrm{ACSSL}}$ is the anatomical cross-sequence objective,
$\mathcal{L}_{\mathrm{miss}}$ is optional missing-modality consistency,
$\mathcal{L}_{\mathrm{cost}}$ is the optional clinical cost term and
$\mathcal{L}_{\mathrm{reg}}$ collects standard regularisation.

Graph message passing is normally part of the forward model rather than a
separate loss. If an explicit graph-consistency regulariser is introduced, it
receives its own coefficient and ablation.

Every non-zero $\lambda$ is treated as a hyperparameter whose contribution must
be demonstrated. A loss term is removed if it adds no reproducible value. The
final thesis reports the selected values and the search range:
\torecord{$\lambda$ values, temperature, optimiser and regularisation settings}.

% ===========================================================================
\section{Optimisation and Training Protocol}
\label{sec:method-training}

\subsection{Training Phases}
\label{sec:method-training-phases}

Training is separated into conceptually distinct phases:

\begin{enumerate}
  \item \textbf{Anatomical/localisation training or loading.}
        The localiser is trained/reproduced and then fixed for the grading
        experiments unless a joint-training ablation is explicitly conducted.

  \item \textbf{Self-supervised pretraining.}
        Sequence encoders are trained with anatomical cross-sequence pairs using
        development patients only.

  \item \textbf{Supervised grading.}
        Encoders are fine-tuned with the routing and graph modules on
        LumbarDISC labels.

  \item \textbf{Calibration.}
        Temperature/uncertainty parameters are fitted on validation data after
        model selection.

  \item \textbf{External evaluation.}
        The selected model is frozen and evaluated on the fixed Rizgary test
        set.

  \item \textbf{Few-shot adaptation.}
        New model copies are adapted using only the local adaptation pool and
        re-evaluated on the same fixed local test set.
\end{enumerate}

\subsection{Optimiser and Schedules}
\label{sec:method-optimiser}

AdamW is the default optimiser for transformer/graph components; SGD or AdamW
may be compared for CNN baselines if required. Learning rate is selected on the
validation set within a pre-specified range, with warm-up and cosine decay or a
plateau scheduler used consistently across the main ablation. Early stopping is
based on a prevalence-robust validation metric such as macro-F1 or weighted
$\kappa$, not raw accuracy.

The final report records:

\begin{itemize}
  \item optimiser and version;
  \item initial learning rate and scheduler;
  \item batch size / effective batch size;
  \item number of epochs and early-stopping patience;
  \item weight decay and dropout;
  \item input resolution and slice-stack radius;
  \item precision mode (FP32/BF16/FP16);
  \item random seeds;
  \item hardware and software versions;
  \item wall-clock and GPU time.
\end{itemize}

No final hyperparameter value is fabricated in this protocol chapter; each is
\torecord{executed training configuration and model-selection criterion}.

\subsection{Augmentation}
\label{sec:method-augmentation}

Training augmentation is designed to simulate plausible MRI variability without
altering disease anatomy. Candidate transforms include modest intensity scaling,
gamma/contrast change, bias-field simulation, noise, small translation/rotation
and limited elastic deformation. Aggressive transformations that distort the
canal or foraminal morphology beyond plausible acquisition variation are
excluded.

Laterality-aware augmentation updates all side-specific labels. If a horizontal
flip changes anatomical left/right, left and right target labels and graph node
identity are swapped in the same transformation. Augmentation is disabled for
validation and test data except in a separately labelled test-time-augmentation
experiment.

\subsection{Model Selection}
\label{sec:method-model-selection}

Only validation data are used to choose architecture depth, routing
hyperparameters, loss weights, early stopping and calibration. The held-out
public test and fixed Rizgary test sets are not inspected for iterative model
selection. If many exploratory experiments are performed, only the pre-specified
final comparison is used for confirmatory statistical claims; exploratory
results are labelled as such.


% ===========================================================================
\section{Experimental Ladder and Ablation Programme}
\label{sec:method-ablation}

The main experiment is not one comparison between ``AMOG-Net'' and a baseline.
It is a ladder in which each additional assumption is isolated.

\begin{longtable}{p{1.1cm}p{5.0cm}p{7.2cm}}
\toprule
\textbf{ID} & \textbf{Configuration} & \textbf{Question isolated}\\
\midrule
\endfirsthead
\toprule
\textbf{ID} & \textbf{Configuration} & \textbf{Question isolated}\\
\midrule
\endhead
E0 & Independent ROI classifier & How well can anatomically localised targets be graded without relational or adaptive fusion?\\
E1 & + DICOM-aligned multi-sequence ROIs & What is gained by correct cross-plane correspondence under fixed fusion?\\
E2 & + disease-conditioned sequence routing & Does target-specific routing outperform fixed fusion?\\
E3 & + modality dropout / missing-sequence training & Does explicit training for absence improve robustness?\\
E4 & + anatomical cross-sequence SSL & Does anatomy-defined pretraining improve representation and label efficiency?\\
E5 & + homogeneous target graph & Does generic message passing help?\\
E6 & + typed heterogeneous disease--anatomy graph & Do clinically defined edge types matter beyond graph capacity?\\
E7 & + ordinal/cost-sensitive/calibrated heads & Can clinically consequential errors and probability reliability be improved?\\
E8 & Selected complete system + external transfer & Do the selected components survive institutional shift and adapt efficiently?\\
\bottomrule
\end{longtable}

The ordering may be adjusted for engineering practicality, but the comparisons
themselves remain. In particular, E6 must be compared with both E5 and a
shuffled-edge control. Otherwise any improvement can be attributed to additional
parameters rather than anatomical topology.

\subsection{Contribution-Specific Ablations}
\label{sec:method-contribution-ablation}

Three ablation families correspond directly to the doctoral contributions.

\textbf{Anatomical SSL ablation} compares no SSL, generic image SSL and
DICOM-aligned cross-sequence SSL under identical downstream training. It is
repeated at reduced label fractions.

\textbf{Routing ablation} compares fixed fusion, cross-attention, target-
conditioned routing and target-conditioned routing with modality dropout.
Performance is reported under complete and missing-sequence input.

\textbf{Graph ablation} compares independent heads, an ordered five-level
transformer, homogeneous graph, heterogeneous graph and random/shuffled edges.
Edge-family removal further identifies whether adjacency, bilateral symmetry or
same-level cross-condition interaction carries the signal.

\subsection{Negative Results}
\label{sec:method-negative-results}

A component that fails to improve the pre-specified outcome is not silently
removed from the research record. The thesis reports it, including confidence
intervals and target-specific effects. This is particularly important for
ordinal objectives and graph modelling, because both are theoretically
attractive and therefore vulnerable to confirmation bias.

A component enters the final integrated system only if it adds a reproducible
benefit, improves an outcome that the thesis has defined as clinically relevant,
or yields a meaningful robustness/calibration advantage despite similar
aggregate discrimination.

% ===========================================================================
\section{External Validation at Rizgary}
\label{sec:method-external}

\subsection{Definition of Zero-Shot Transfer}
\label{sec:method-zero-shot}

Zero-shot transfer is defined strictly. The selected model, including encoder,
router, graph, classifier and calibration strategy developed on public data, is
frozen before any Rizgary test label is used. The local images are passed
through the same preprocessing and localisation pipeline, with only
non-learned DICOM reconstruction and pre-specified per-volume normalisation
allowed.

No local fine-tuning, threshold optimisation or calibration is permitted in the
zero-shot result. If a local intensity harmonisation method is fitted on Rizgary
data, that constitutes a separate domain-adaptation experiment and is labelled
accordingly.

The primary zero-shot analysis is restricted to the central-canal targets for
which the local reports provide a defensible mapping. The principal comparison
is therefore not the full 25-target benchmark score versus a five-target local
score; it is public held-out versus local external performance on the
\emph{same central-canal task}.

\subsection{Quantifying Domain Shift}
\label{sec:method-domain-shift}

Performance degradation is reported as both absolute and relative change:

\begin{align}
\Delta M &= M_{\mathrm{Rizgary}}-M_{\mathrm{Public}},\\
\Delta M_{\%} &=
100\,
\frac{M_{\mathrm{Rizgary}}-M_{\mathrm{Public}}}
{M_{\mathrm{Public}}},
\end{align}

for each primary metric $M$. The thesis does not interpret a large negative
$\Delta M$ as failed research; the size and structure of the drop are among the
main outcomes of RQ4.

The shift is decomposed by:

\begin{itemize}
  \item lumbar level;
  \item severity class;
  \item available sequence configuration;
  \item scanner/vendor/field strength where the local cohort contains enough
        variation;
  \item slice thickness/in-plane resolution;
  \item localisation success;
  \item report/reference-standard confidence.
\end{itemize}

This analysis separates representational failure from data-quality failure
where possible. For example, if grading declines but localisation remains
stable, the pattern supports the detection-versus-grading asymmetry documented
in Chapter~2.

\subsection{Local Error Review}
\label{sec:method-local-error-review}

A blinded error audit is performed after the zero-shot metrics are locked.
Incorrect local predictions are grouped into:

\begin{itemize}
  \item wrong level/localisation;
  \item adjacent-grade disagreement;
  \item distant Severe$\leftrightarrow$Normal/Mild disagreement;
  \item apparent label/report ambiguity;
  \item missing or atypical sequence;
  \item image-quality/artifact issue;
  \item plausible model error despite adequate evidence.
\end{itemize}

Where reader time permits, a subset of model/report disagreements is re-read
without exposing the model prediction. The purpose is not to alter the test
labels opportunistically but to estimate how much apparent model error may
reflect reference-standard variability.

% ===========================================================================
\section{Few-Shot Domain Adaptation}
\label{sec:method-adaptation}

\subsection{Adaptation Sizes}
\label{sec:method-adaptation-sizes}

The adaptation pool is sampled at

\[
N\in\{10,25,50,100\}
\]

patient-level labelled cases, subject to the final number of eligible local
studies. Each $N$ is sampled multiple times, stratified where feasible by
central-canal severity. The same fixed local test set is used after every
adaptation.

The experiment therefore estimates a function

\begin{equation}
M(N)=
\text{test performance after adaptation with }N\text{ local cases},
\label{eq:adaptation-curve}
\end{equation}

rather than a single ``fine-tuned'' number.

\subsection{Adaptation Strategies}
\label{sec:method-adaptation-strategies}

Four strategies are compared:

\begin{enumerate}
  \item \textbf{No adaptation.} The zero-shot baseline.
  \item \textbf{Non-learned/simple harmonisation.} Intensity preprocessing is
        adjusted using only the adaptation pool where a defensible method exists.
  \item \textbf{Parameter-efficient adaptation.} Adapter layers or low-rank
        updates are trained while most of the public-data representation remains
        fixed.
  \item \textbf{Full fine-tuning.} All or most parameters are updated, where
        computationally feasible and where $N$ is large enough to avoid
        immediate overfitting.
\end{enumerate}

The exact PEFT implementation is chosen after a small validation comparison and
is not itself claimed as novel. Its role is to ask an operational question:
how many local labels are needed when a hospital cannot afford to retrain a full
model from scratch?

\subsection{Annotation-Efficiency Summary}
\label{sec:method-annotation-efficiency}

For each method, the recovery curve reports macro-F1, weighted $\kappa$, severe
recall and calibration as a function of $N$. A useful summary is the area under
the performance-versus-log-annotation curve, but raw points with confidence
intervals remain primary because the curve may be non-linear.

No arbitrary threshold such as ``95\% recovery'' defines success. If 100 local
cases recover little performance, that is a meaningful negative finding about
transportability. If 10 cases recover most of the loss, that is evidence for
annotation-efficient adaptation.

\subsection{Avoiding Adaptation Overfit}
\label{sec:method-adaptation-overfit}

Few-shot training is particularly vulnerable to overfitting and selection bias.
The study therefore:

\begin{itemize}
  \item repeats each $N$ across several sampled adaptation sets;
  \item uses regularisation and early stopping based only on an adaptation
        validation split or internal cross-validation within the adaptation
        pool;
  \item never selects the adaptation epoch from the fixed external test;
  \item reports the variance across sampled local subsets;
  \item compares PEFT with full fine-tuning rather than assuming fewer trainable
        parameters are always superior.
\end{itemize}

% ===========================================================================
\section{Evaluation Metrics}
\label{sec:method-metrics}

No single metric is treated as sufficient because the task is imbalanced,
ordinal and multi-target.

\subsection{Macro F1 and Balanced Accuracy}
\label{sec:method-macro}

For class $k$, precision and recall are

\begin{align}
P_k &= \frac{TP_k}{TP_k+FP_k},\\
R_k &= \frac{TP_k}{TP_k+FN_k},
\end{align}

and

\begin{equation}
F1_k=
2\frac{P_kR_k}{P_k+R_k}.
\end{equation}

Macro-F1 is the unweighted class mean,

\begin{equation}
F1_{\mathrm{macro}}
=
\frac{1}{K}\sum_{k=1}^{K}F1_k,
\label{eq:macro-f1}
\end{equation}

so the severe class cannot be hidden by the Normal/Mild majority.

Balanced accuracy is

\begin{equation}
BA=
\frac{1}{K}\sum_{k=1}^{K}R_k.
\label{eq:balanced-accuracy}
\end{equation}

These are primary discrimination metrics.

\subsection{Quadratic Weighted Kappa}
\label{sec:method-kappa}

Quadratic weighted $\kappa$ is reported because the labels are ordered and the
literature supplies human agreement values in $\kappa$. For observed confusion
matrix $O$ and expected matrix $E$, with weight

\[
w_{ij}=\frac{(i-j)^2}{(K-1)^2},
\]

the statistic is

\begin{equation}
\kappa_w
=
1-
\frac{
\sum_{ij}w_{ij}O_{ij}
}{
\sum_{ij}w_{ij}E_{ij}
}.
\label{eq:weighted-kappa}
\end{equation}

Exact agreement and within-one-grade agreement are also reported so that a high
$\kappa$ is not the only view of ordinal error.

\subsection{Severe-Class Outcomes}
\label{sec:method-severe}

The clinically decisive class is reported separately:

\begin{itemize}
  \item Severe recall/sensitivity;
  \item Severe precision;
  \item Severe-versus-rest AUROC and AUPRC where appropriate;
  \item Severe$\rightarrow$Normal/Mild error rate;
  \item Severe$\rightarrow$Moderate error rate.
\end{itemize}

The separation distinguishes distant under-grading from adjacent boundary error.

\subsection{AUROC and AUPRC}
\label{sec:method-auroc}

For three-class outcomes, AUROC is calculated one-versus-rest and macro-averaged.
AUPRC is reported alongside AUROC for rare Severe classes because it is more
sensitive to prevalence and false positives. Curves are constructed at patient-
target level, while confidence intervals respect patient clustering.

\subsection{Calibration}
\label{sec:method-calibration-metrics}

Calibration evaluation includes:

\begin{itemize}
  \item Brier score;
  \item Expected Calibration Error (ECE);
  \item classwise reliability diagrams;
  \item negative log likelihood where useful;
  \item risk--coverage curves under selective prediction.
\end{itemize}

Calibration is reported both internally and externally because a model can
retain discrimination while becoming overconfident after domain shift.

\subsection{Localisation Metrics}
\label{sec:method-localisation-metrics}

Localisation is evaluated separately using:

\begin{itemize}
  \item mean, median and 95th-percentile Euclidean error in millimetres;
  \item PCK at physically meaningful tolerances;
  \item Dice coefficient for available masks;
  \item surface-distance metrics where mask quality supports them;
  \item complete localisation-failure rate per study.
\end{itemize}

\subsection{Efficiency Metrics}
\label{sec:method-efficiency}

For each model configuration the thesis reports:

\begin{itemize}
  \item trainable and total parameter count;
  \item FLOPs/MACs where they can be estimated consistently;
  \item peak VRAM;
  \item training wall time and GPU hours;
  \item inference time per study and per target;
  \item storage size of weights.
\end{itemize}

Efficiency comparisons use the same hardware and batch/inference conditions.
Timing from different machines is not compared as if it were an architectural
property.

% ===========================================================================
\section{Statistical Analysis}
\label{sec:method-statistics}

\subsection{Patient-Level Bootstrap Confidence Intervals}
\label{sec:method-bootstrap}

Confidence intervals for primary metrics are obtained by bootstrap resampling
\emph{patients}, not individual targets. On each replicate, a patient is sampled
with replacement and all of that patient's target predictions enter together.
This preserves within-study dependence.

The default is a sufficiently large number of replicates to stabilise the
2.5th and 97.5th percentiles; the final number is
\torecord{bootstrap replicate count}. Ninety-five per cent percentile or
bias-corrected intervals are reported consistently.

\subsection{Paired Model Comparisons}
\label{sec:method-paired-tests}

Ablation models are evaluated on the same patients, so comparisons are paired.
For scalar patient-level losses or correctness indicators, a paired permutation
test or Wilcoxon signed-rank test is used where distributional assumptions are
not justified. AUCs are compared with an appropriate paired procedure such as
DeLong's test when its assumptions are satisfied. Bootstrap differences are
used for metrics without a simple analytic test.

The primary quantity is the paired difference

\[
\Delta M=M_A-M_B
\]

with a 95\% confidence interval. Point estimates alone are not used to claim
one component is superior.

\subsection{Multiple Comparisons}
\label{sec:method-multiple}

The ablation programme contains many target-level analyses. A limited set of
primary comparisons is therefore declared in advance:

\begin{enumerate}
  \item anatomical SSL versus generic/self-supervised baseline;
  \item routed+dropout fusion versus ordinary cross-attention;
  \item heterogeneous graph versus independent heads and homogeneous graph;
  \item zero-shot public versus local performance;
  \item PEFT versus no adaptation across annotation sizes.
\end{enumerate}

Target-by-target and subgroup analyses are secondary/exploratory. Where many
formal hypotheses are tested, false-discovery-rate control is applied. The
thesis distinguishes confirmatory comparisons from exploratory patterns instead
of presenting every $p<0.05$ as independent evidence.

\subsection{Repeated Training Seeds}
\label{sec:method-seeds}

Deep-learning optimisation is stochastic. Final ablation comparisons are
therefore repeated across multiple fixed seeds where computationally feasible.
The thesis reports mean and standard deviation/interval across seeds in addition
to the patient-level test uncertainty. A one-off training run is insufficient
to attribute a small improvement to a method.

\subsection{Subgroup Analysis}
\label{sec:method-subgroups}

Subgroup analysis is conducted only where sample size permits and includes:

\begin{itemize}
  \item lumbar level;
  \item severity;
  \item sex/age bands where metadata and governance allow;
  \item scanner/vendor/field strength;
  \item acquisition resolution or slice thickness;
  \item presence/absence of individual sequences.
\end{itemize}

Sparse cells are reported descriptively rather than subjected to unstable
significance testing. The subgroup analysis is intended to expose failure modes,
not to make causal claims about demographic biology.

% ===========================================================================
\section{Robustness Experiments}
\label{sec:method-robustness}

\subsection{Missing-Sequence Matrix}
\label{sec:method-missing-matrix}

The selected public-data models are evaluated under every supported sequence
configuration. For three sequences, the core matrix includes complete input,
three two-sequence combinations and three single-sequence configurations where
the required anatomical target remains observable.

Performance is reported per target family rather than only in aggregate. A
model can remain strong on central canal grading while collapsing on foraminal
narrowing when sagittal T1 is removed; an overall mean would obscure the
clinical meaning.

\subsection{Geometry Perturbation}
\label{sec:method-geometry-robust}

The reliance on DICOM coordinates creates a specific failure mode. Controlled
perturbations therefore test sensitivity to small localisation/geometry errors:
axial slice selection shifted by one or two slices, modest ROI-centre offsets,
and simulated spacing variation. The purpose is not to train on corrupted
geometry by default, but to measure how brittle the system is to the imperfect
localisation expected in practice.

\subsection{Image-Quality and Acquisition Perturbation}
\label{sec:method-acquisition-robust}

Where feasible, synthetic intensity perturbations approximating scanner shift
are applied to the public held-out set: contrast scaling, noise, bias field and
resolution degradation. These do not replace real external validation but help
identify which changes reproduce the pattern observed at Rizgary.

\subsection{Graph Perturbation}
\label{sec:method-graph-robust}

The graph is challenged by:

\begin{itemize}
  \item random edge deletion;
  \item shuffled edge types;
  \item random rewiring with matched edge count;
  \item removal of one relation family.
\end{itemize}

A robust anatomical graph should outperform these controls without depending on
one fragile edge. If random topology performs equivalently, the anatomical
interpretation is rejected.

% ===========================================================================
\section{Interpretability and Failure Analysis}
\label{sec:method-interpretability}

\subsection{Local Visual Evidence}
\label{sec:method-visual-evidence}

Grad-CAM or equivalent saliency may be used as a gross sanity check that the
encoder responds within the relevant ROI, but saliency is not presented as a
proof of reasoning. The stronger explanations in this thesis arise from
structural quantities that are part of the model itself:

\begin{itemize}
  \item which sequence received routing weight for a specific target;
  \item which graph relation contributed to target refinement;
  \item how uncertainty changed before and after relational reasoning;
  \item how prediction changed under controlled removal of a sequence or edge.
\end{itemize}

These are still model diagnostics, not causal explanations, but they are more
directly linked to the proposed mechanisms than a heatmap alone.

\subsection{Case-Level Audit Template}
\label{sec:method-case-audit}

A fixed template is used for qualitative audit of representative correct and
incorrect cases:

\begin{itemize}
  \item ground-truth target and grade;
  \item predicted probability distribution and calibrated grade;
  \item localisation overlay;
  \item available MRI sequences;
  \item routing weights;
  \item dominant incoming graph relations/attention where extractable;
  \item uncertainty;
  \item nearest error category;
  \item whether the discrepancy is within one grade.
\end{itemize}

Cases are sampled from predefined categories rather than cherry-picked from the
best-looking outputs.

% ===========================================================================
\section{Reproducibility and Software Engineering}
\label{sec:method-reproducibility}

\subsection{Environment}
\label{sec:method-environment}

The implementation will use a reproducible Python/PyTorch medical-imaging
environment with DICOM tooling such as pydicom, SimpleITK and/or MONAI.
The final thesis records exact versions of Python, PyTorch, CUDA, cuDNN, GPU
hardware and all key libraries. Environment specification is stored in a
lockfile or container recipe.

\subsection{Experiment Configuration}
\label{sec:method-config}

Each run is configured from a machine-readable file containing:

\begin{itemize}
  \item dataset manifest and split version;
  \item selected sequence IDs;
  \item preprocessing and resampling;
  \item ROI dimensions;
  \item encoder architecture;
  \item SSL/routing/graph settings;
  \item loss coefficients;
  \item optimiser and scheduler;
  \item seed;
  \item checkpoint and early-stopping rule.
\end{itemize}

The configuration is copied into the run directory with metrics and model
weights so that a table in the thesis can be traced to an executable experiment.

\subsection{Data Provenance}
\label{sec:method-provenance}

Raw data are never edited in place. A manifest records the path/identifier and
processing status of every study. Derived crops can be regenerated from the
manifest and DICOM source, and each derived sample retains its pseudonymous
patient, level, condition, side, sequence and source-series identifier.

The local reference matrix is versioned independently of model output. If a
label is corrected after clinical review, the change is documented and the
affected experiments are rerun rather than silently updating test truth.

\subsection{Code and Model Release}
\label{sec:method-release}

Code for the public-data methodology will be released where institutional and
dataset licences permit. Pretrained weights are released if the training
dataset terms allow redistribution. Local Rizgary images and report-derived
individual-level data are not included unless a separate hospital approval
authorises public release.

The repository will include:

\begin{itemize}
  \item preprocessing and geometry reconstruction;
  \item split-generation/verification code;
  \item localisation/ROI pipeline;
  \item SSL, routing and graph modules;
  \item training and evaluation scripts;
  \item configuration files for primary experiments;
  \item instructions for reproducing public benchmark results.
\end{itemize}

% ===========================================================================
\section{Methodological Safeguards and Threats to Validity}
\label{sec:method-threats}

\subsection{Reference-Standard Shift}
\label{sec:method-threat-label}

LumbarDISC uses structured expert grades; Rizgary uses routine narrative
reports. A performance drop can therefore reflect both imaging-domain shift and
reference-standard shift. The thesis mitigates this by restricting transfer to
defensible overlap, auditing report labels and, where possible, re-reading a
stratified subset. It does not claim that the two reference standards are
identical.

\subsection{Single-Centre Local External Validation}
\label{sec:method-threat-single-centre}

Rizgary is a genuinely independent institution but still only one local centre.
External performance therefore demonstrates transport to \emph{that} setting,
not universal Middle Eastern generalisation. The thesis states this boundary
explicitly. A future multi-hospital regional study is required before making a
population-wide deployment claim.

\subsection{Model Complexity}
\label{sec:method-threat-complexity}

The full architecture contains several interacting components. Complexity can
produce apparent gains merely through parameter count and tuning effort. The
ablation ladder, matched backbone, homogeneous graph and shuffled-edge controls
are specifically designed to separate mechanism from capacity.

\subsection{Localisation Error}
\label{sec:method-threat-localisation}

The graph and routing modules cannot recover evidence absent from a bad crop.
Localisation is therefore measured and failures are reported separately.
Downstream results may additionally be stratified by localisation error to
estimate how much grading depends on the enabling stage.

\subsection{Class Imbalance}
\label{sec:method-threat-imbalance}

Severe grades are uncommon, so accuracy can be misleading and few-shot local
subsets may contain few Severe examples. The primary metrics are macro-F1,
balanced accuracy, weighted $\kappa$ and per-class recall, with patient-level
confidence intervals. Adaptation subsets are stratified where feasible and
their severity counts are always reported.

\subsection{Hyperparameter Multiplicity}
\label{sec:method-threat-hyperparameters}

A sufficiently large search can overfit the validation set even without touching
test data. The study limits the number of tunable components in each stage,
uses a fixed primary backbone for the core ablation, records the search space
and freezes choices before external testing.

\subsection{Clinical Utility}
\label{sec:method-threat-clinical}

This methodology evaluates technical grading, calibration and portability. It
does not demonstrate that radiologists become faster, safer or more accurate
when using the model. Such a claim requires a prospective reader-assistance or
impact study with independent clinical approval. The thesis therefore uses the
language of ``decision-support potential'' rather than clinical benefit.

% ===========================================================================
\section{Mapping Research Questions to Evidence}
\label{sec:method-rq-mapping}

\begin{longtable}{p{1.0cm}p{4.8cm}p{4.4cm}p{3.8cm}}
\toprule
\textbf{RQ} & \textbf{Primary comparison} & \textbf{Primary outcomes} & \textbf{Failure interpretation}\\
\midrule
\endfirsthead
\toprule
\textbf{RQ} & \textbf{Primary comparison} & \textbf{Primary outcomes} & \textbf{Failure interpretation}\\
\midrule
\endhead
RQ1 & independent/level transformer/homogeneous graph vs typed heterogeneous graph & macro-F1, $\kappa$, severe recall, edge ablations & graph topology carries little useful signal or local evidence dominates\\
RQ2 & generic pretraining vs anatomical cross-sequence SSL at 10/25/50/100\% labels & macro-F1, $\kappa$, label efficiency, external drop & anatomy-defined positives do not improve transferable representation\\
RQ3 & fixed fusion/cross-attention vs routing $\pm$ modality dropout & missing-sequence degradation, target-specific performance, calibration & fixed fusion is sufficient or learned routing is unstable\\
RQ4 & public held-out vs fixed Rizgary test; no adaptation vs PEFT/full tuning & zero-shot $\Delta$, recovery curve, calibration shift & domain shift is larger or less annotation-efficient than expected\\
RQ5 & CE vs ordinal/cost-aware objectives; calibrated vs uncalibrated prediction & Severe$\rightarrow$Normal rate, grade distance, Brier/ECE, risk--coverage & ordinal/cost formulation offers no advantage over CE\\
\bottomrule
\end{longtable}

This table defines what evidence can answer each question. It also makes the
thesis robust to partial negative results: RQ1 can be answered if the graph does
not help, and RQ2 can be answered if generic pretraining is equal or superior.
The contribution is the controlled experiment and the resulting knowledge, not
a requirement that every proposed mechanism win.

% ===========================================================================
\section{Planned Reporting Standard}
\label{sec:method-reporting}

The final manuscript/thesis reporting will follow the current applicable
medical-AI and diagnostic/prediction reporting checklists at the time of
submission. At minimum the report will include:

\begin{itemize}
  \item complete cohort flow diagrams and inclusion/exclusion reasons;
  \item reference-standard source, annotator qualifications and agreement where
        available;
  \item exact patient-level split logic and IDs reproducible within the approved
        environment;
  \item all primary metrics with confidence intervals;
  \item per-class Severe performance and ordinal error distances;
  \item calibration and external-domain performance;
  \item ablations for every claimed novel component;
  \item repeated-seed variability;
  \item compute/hardware cost;
  \item explicit limitations and negative results;
  \item a statement of what was not tested clinically.
\end{itemize}

The reporting checklist itself will be versioned with the manuscript because
CLAIM/STARD/TRIPOD+AI guidance may be updated during the candidature.

% ===========================================================================
\section{Chapter Summary}
\label{sec:method-summary}

This chapter translates the research gap of Chapter~\ref{ch:litreview} into a
controlled experimental design. The study begins from a structured definition
of lumbar disease: twenty-five level--condition--laterality targets on the
public benchmark rather than a generic image-level class. DICOM patient-space
geometry is preserved so that sagittal and axial evidence can be aligned
physically; localisation and disease-specific 2.5D ROIs are treated as shared
enabling technology; and all data partitions are defined at patient level.

Three methodological contributions are then tested independently. Anatomically
aligned cross-sequence self-supervision asks whether the same spinal level
viewed under different MRI sequences provides a stronger representation-learning
signal than generic augmentation. Disease-conditioned routing asks whether one
model can allocate different sequence evidence to different targets while
remaining functional when modalities are absent. The heterogeneous
disease--anatomy graph asks whether explicit typed relations among adjacent
levels, co-occurring compartments and bilateral targets improve grading beyond
independent heads, an inter-level transformer or a homogeneous graph. Ordinal,
cost-sensitive and uncertainty-aware outputs remain supporting methods and are
retained only if their contribution is measurable.

The translational design is deliberately separated from model development.
Rizgary is not mixed into public training. A fixed local test set establishes
zero-shot domain shift, while a separate adaptation pool supplies controlled
few-shot subsets. The primary local comparison is restricted to central-canal
stenosis, where the benchmark and narrative local reports can support a
defensible common target; local herniation morphology remains a separate task
rather than being forced into the RSNA schema.

Finally, the methodology defines success as an answer to the research
questions, not as attainment of a predetermined accuracy. A graph that adds
nothing, an ordinal loss that fails to beat cross-entropy, or a large
cross-institutional performance drop are all informative if demonstrated under
the same controlled protocol. This principle is what turns the architecture
from an engineering assembly into a doctoral experiment: every claimed
mechanism is falsifiable, every comparison is tied to a research question, and
the external cohort is used to test portability rather than to decorate an
internal result.

\printbibliography

\end{document}
